

\medterm{B And T Cell Screen} A blood test that measures B cells and T cells, two major groups of immune cells. It may help evaluate suspected immune deficiency, unusual blood-cell counts, or some blood cancers, but results must be interpreted with other tests.

\medterm{B And T Cells} B cells make antibodies, while T cells coordinate immune responses or destroy infected and abnormal cells. Both are types of lymphocytes, a group of white blood cells that protect against infection and disease.
\medterm{B And T Lymphocytes} Two major types of immune cells. B lymphocytes can make antibodies and immune memory, while T lymphocytes help control immune responses or destroy infected and abnormal cells.

\medterm{B Cell} A type of white blood cell that develops in the bone marrow. After activation, it can become an antibody-producing cell or a memory cell that helps the immune system respond faster to a later infection.

\synonyms
B Lymphocyte; Bursa-Equivalent Lymphocyte

\medterm{B Hemophilia} Alternate name; see \textit{Hemophilia B}.

\medterm{B Lymphocyte} Alternate name; see \textit{B Cell}.

\medterm{B Virus} A herpes virus carried mainly by macaques. It can cause a rare but severe infection of the brain and spinal cord in people after a bite, scratch, or contact with infected animal tissue.

\medterm{B-Cell Activation} The process in which a B cell receives signals after recognizing a foreign substance, then multiplies and develops into antibody-producing cells or long-lasting memory cells.

\medterm{B-Cell Leukemia/Lymphoma Panel} A laboratory test that examines cells from blood, bone marrow, or tissue for abnormal groups of B cells. It uses cell features and marker proteins to help identify or classify some blood cancers.

\medterm{B-Cell Lymphoma} A group of cancers that begin in B lymphocytes. Different types grow at different speeds and may affect lymph nodes or organs; diagnosis and treatment depend on the exact type, extent, and person's health.

\medterm{B-Cell Receptor} A protein on the surface of a B cell that binds a particular foreign substance. This binding sends signals that can make the cell multiply, mature, and produce antibodies.

\medterm{B-Lymphocyte} Alternate spelling; see \textit{B Cell}.

\medterm{B-Lynch Suture} A surgical stitch placed around the uterus to control severe bleeding after childbirth when the uterus remains relaxed and medicines have not stopped the bleeding. It compresses the uterus while usually preserving the possibility of future pregnancy.

\medterm{B-Mode Ultrasound} The usual two-dimensional ultrasound image created from echoes of sound waves sent into the body. It shows the size, shape, and structure of organs and soft tissues without using radiation.

\medterm{B-Type Natriuretic Peptide} A hormone released mainly by the heart when its chambers are stretched or under pressure. A blood test for BNP or the related NT-proBNP can support assessment of heart failure, but age, kidney disease, rhythm, body size, and other illnesses affect the result.

\synonyms
BNP; Brain Natriuretic Peptide; NT-proBNP

\medterm{Ba} The chemical symbol for barium, a metallic element used in compounds such as barium sulfate. Barium sulfate can be swallowed or given as an enema to outline parts of the digestive tract on X-rays.

\medterm{Babcock Clamp} A surgical instrument with smooth, rounded jaws used to hold delicate tubular tissue gently. It may be used on the intestine, fallopian tubes, or other tissues that could be injured by a crushing clamp.

\medterm{Babesia} A group of parasites transmitted mainly by certain ticks. They infect red blood cells and can cause fever, tiredness, anaemia, low platelets, jaundice, or dark urine; severe illness is more likely without a spleen or with weak immunity.

\medterm{Babesia Divergens} A species of parasite that infects red blood cells and can cause babesiosis, especially in Europe. It may cause severe destruction of red cells, particularly in people without a spleen or with a weakened immune system.

\medterm{Babesia Microti} A parasite and the most common cause of babesiosis in North America. It is spread by black-legged ticks, infects red blood cells, and may cause no symptoms or severe anaemia, jaundice, and organ failure.

\medterm{Babesiosis} An infection caused by Babesia parasites, usually spread by ticks. The parasites enter red blood cells and may cause fever, tiredness, anaemia, low platelets, jaundice, or dark urine; severe disease can affect several organs.

\medterm{Babinski Reflex} A reflex tested by stroking the sole of the foot. In babies, the big toe normally moves upward and the other toes spread; in older children and adults, the toes usually curl downward. An upward response later in life may indicate a brain or spinal-cord problem.

\medterm{Babinski Response} Alternate name; see \textit{Babinski Reflex}.

\medterm{Babinski Sign} Alternate name; see \textit{Babinski Reflex}.

\medterm{Baby Acne} Small red or white bumps, usually on an infant's face, caused by temporary effects of hormones. It commonly improves without treatment; gentle washing is sufficient, and adult acne products are not used on infant skin unless a clinician advises.

\medterm{Baby Blues} A common short period of tearfulness, irritability, anxiety, or mood changes beginning in the first days after childbirth. It usually improves within two weeks; persistent, severe, or dangerous thoughts may indicate postnatal depression or an emergency.

\medterm{Baby Health} An infant's physical growth, development, feeding, sleep, immunizations, screening, and protection from illness or injury. Regular care also monitors hearing, vision, movement, communication, and social development.

\medterm{Baby Holding Techniques} Safe ways to support an infant's head, neck, back, and body while carrying, feeding, or moving the baby. Newborns need full head and neck support; never shake an infant or leave one unattended on a raised surface.

\medterm{Babyhood} The period of infancy, from birth through the first year of life. During this time, babies undergo rapid growth and development of movement, communication, feeding skills, immunity, and social responses.

\medterm{Bacillaceae} A family of rod-shaped bacteria that includes species able to form resistant spores. Some members are harmless or useful, while others can cause infections such as anthrax or food poisoning.

\medterm{Bacillary Angiomatosis} A bacterial infection, usually caused by Bartonella, that produces red or purple skin bumps in people with weakened immunity. It can also affect the liver, spleen, bones, or lymph nodes and may resemble other skin tumours.

\medterm{Bacillary Dysentery} An intestinal infection, commonly caused by Shigella bacteria, that produces fever, abdominal cramps, and frequent diarrhoea that may contain blood or mucus. Severe dehydration or complications require medical care.

\medterm{Bacille Calmette-Guérin} A weakened form of a mycobacterium used as a vaccine against tuberculosis in many countries. It is also placed into the bladder as an immune treatment for some early bladder cancers.

\medterm{Bacilluria} The presence of rod-shaped bacteria in urine, detected by microscopy or culture. It may represent contamination, harmless colonization, or a urinary infection; symptoms, the amount of bacteria, and the organism determine its significance.

\medterm{Bacillus} A rod-shaped bacterium. The word can describe bacterial shape or members of the Bacillus genus, many of which form resistant spores; some are harmless, while others cause disease.

\medterm{Bacillus Anthracis} A spore-forming bacterium that causes anthrax in people and animals. Infection may affect the skin, digestive tract, or lungs and can rapidly become life-threatening, requiring urgent specialist treatment and public-health action.

\medterm{Bacillus Cereus} A spore-forming bacterium that can cause food poisoning after contaminated food is eaten. It may cause sudden vomiting or watery diarrhoea and, rarely, serious infections in people with weakened immunity.

\medterm{Bacillus Cereus Food Poisoning} A foodborne illness caused by toxins produced by Bacillus cereus, manifesting in two distinct clinical forms: an emetic syndrome characterized by acute nausea and vomiting developing 1 to 5 hours after ingestion (classically linked to fried or reheated rice), and a diarrheal syndrome with abdominal cramps and watery diarrhea developing 8 to 16 hours after ingestion (associated with contaminated meats, vegetables, or sauces).

\medterm{Bacitracin} An antibiotic applied to the skin that prevents susceptible bacteria from building their protective cell wall. It is used for selected superficial infections, but it can cause skin irritation or allergic reactions.

\medterm{Bacitracin Overdose} Harm caused by too much exposure to bacitracin. Skin use usually causes local irritation, while swallowing a large amount or receiving excessive injected medicine can cause vomiting, allergic reactions, or kidney injury.

\medterm{Bacitracin Zinc Overdose} Harm caused by excessive exposure to bacitracin zinc. The effects depend on the amount and route; large exposures may cause vomiting, allergic reactions, or kidney injury and require prompt medical treatment.

\medterm{Back Injury} Damage to the muscles, ligaments, joints, discs, bones, or nerves of the back after trauma, lifting, or repeated strain. Pain and stiffness are common; weakness, numbness, loss of bladder or bowel control, or major-trauma pain requires urgent assessment.

\medterm{Back Pain} Pain or discomfort in any part of the back, from the neck to the lower back. It is often caused by muscles, joints, or discs, but may also arise from nerves or internal organs.

\synonyms
Dorsalgia; Backache; Spinal Pain

\medterm{Back Pressure} Pressure that builds up behind an obstruction or against the normal direction of flow. It may occur in an airway, blood vessel, duct, or hollow organ and can impair movement of air, blood, or fluid.

\medterm{Back-Blows} Firm blows delivered between the shoulder blades as first aid for severe choking in a responsive person. They are used with age-appropriate choking procedures and do not replace emergency help.

\medterm{Backache} Pain or discomfort in the back, from the neck to the buttocks. It is a symptom rather than a disease and may result from muscles, joints, discs, nerves, or pain referred from an internal organ.

\medterm{Backwash Ileitis} Mild, superficial inflammation of the terminal ileum that occurs in patients with extensive ulcerative colitis (pancolitis). It results from an open, incompetent ileocecal valve that allows inflamed colonic contents and bacteria to reflux backward into the small intestine. Unlike Crohn's disease, the inflammation is continuous with the cecum, is confined to the mucosa without strictures or granulomas, and resolves after surgical removal of the colon. Its presence is associated with severe colitis and an increased long-term risk of colorectal dysplasia.

\synonyms
Reflux Ileitis; Ulcerative Colitis Backwash Ileitis

\medterm{Baclofen} A medicine that reduces muscle spasm by lowering nerve signals in the spinal cord. It can cause sleepiness, weakness, dizziness, or confusion; stopping it suddenly after regular use can cause dangerous withdrawal symptoms.

\medterm{Bacteraemia} Alternate spelling; see \textit{Bacteremia}.

\medterm{Bacteremia} The presence of living bacteria in the bloodstream. It may be temporary, intermittent, or continuous. Some cases cause no illness, while others produce fever and can progress to sepsis or infection of distant organs.

\medterm{Bacteria} Microscopic single-celled organisms. Many live harmlessly or help the body, while others cause infections; bacteria can live alone or in communities and may develop resistance to antibiotics.

\medterm{Bacterial Adherence} The process by which bacteria attach to body cells or surfaces, often the first step in infection. Surface proteins and tiny hair-like structures help bacteria stay in place and resist removal.

\medterm{Bacterial Biofilm} A community of bacteria attached to a surface and enclosed in a protective material they produce. Biofilms can resist immune defenses and antibiotics and cause persistent infections on wounds, catheters, implants, or teeth.

\medterm{Bacterial Capsule} A protective outer coating around some bacteria, often made of sugars. It can help bacteria avoid being swallowed by immune cells, survive in the body, and cause disease.

\medterm{Bacterial Cell Wall} A strong outer layer that gives bacteria their shape and prevents them from bursting. Many bacterial cell walls contain peptidoglycan, a material made from linked sugars and short protein chains.

\medterm{Bacterial Classification} The grouping of bacteria according to their genetic features, appearance, chemistry, and evolutionary relationships. Classification helps identify organisms, understand relatedness, and guide diagnosis, infection control, and treatment.

\medterm{Bacterial Conjugation} The transfer of DNA from one bacterium to another through direct contact. The transferred DNA may include genes that help bacteria resist antibiotics or cause disease.

\medterm{Bacterial Conjunctivitis} A contagious bacterial infection of the thin surface covering the eye and inner eyelids. It commonly causes redness, irritation, gritty discomfort, and thick discharge that sticks the eyelids together. Eye pain, light sensitivity, reduced vision, or contact-lens use needs prompt assessment.

\medterm{Bacterial Dermatitis} Skin inflammation caused by bacterial infection. The skin may become red, warm, swollen, painful, or covered with pus or crusts. It can start in healthy skin or develop when bacteria infect an existing rash or wound.

\medterm{Bacterial Dormancy} A low-activity state that allows bacteria to survive conditions in which they cannot grow. Some form highly resistant spores; others remain as ordinary cells that temporarily tolerate stress and may become active again when conditions improve.

\medterm{Bacterial Endocarditis} A serious infection of the heart’s inner lining or valves. It may cause fever, chills, tiredness, a new or changing murmur, or weight loss. Bacteria can damage valves or break off and block blood vessels; diagnosis usually requires blood cultures and heart imaging.

\medterm{Bacterial Endocarditis Vegetations} Clumps of bacteria, clotting material, and platelets that form on a heart valve or the heart's inner lining during bacterial endocarditis. They can develop when bacteria in the blood attach to a damaged or abnormal heart surface.

\medterm{Bacterial Endospores} Highly resistant resting structures made by some bacteria, including Bacillus and Clostridium, when nutrients are scarce. They can survive drying, heat, chemicals, and long periods of time, then grow into active bacteria when conditions improve.

\medterm{Bacterial Flagella} Thread-like structures that help some bacteria move through liquid or across surfaces. A rotating base works like a motor and turns the attached filament; flagella can also influence how bacteria attach and spread.

\medterm{Bacterial Gastroenteritis} Acute inflammation of the stomach and intestines caused by pathogenic bacteria such as
Salmonella, Campylobacter, Shigella, or pathogenic Escherichia coli. It commonly causes
diarrhea, abdominal cramps, fever, and sometimes bloody stools or vomiting.

\medterm{Bacterial Genetics} The study of bacterial genes, DNA, inheritance, and genetic change. Bacteria usually have one main chromosome and may carry extra DNA rings called plasmids that can include antibiotic-resistance genes.

\medterm{Bacterial Growth Phases} The stages bacteria pass through when growing in a culture. They first adapt, then multiply rapidly, later reach a balance as nutrients fall and waste builds up, and finally decline as conditions become unsuitable.

\medterm{Bacterial Infection} An illness caused by single-celled microscopic organisms called bacteria entering, multiplying, and spreading in body tissues. While millions of harmless or beneficial bacteria normally live on the skin and in the gut (the microbiome), infections happen when disease-causing bacteria invade the body or when normal bacteria get into areas that are usually sterile, such as the lungs, kidneys, bladder, or bloodstream. Bacteria damage tissues directly, produce poisons called toxins, or trigger intense immune responses that cause swelling, redness, fever, and pus. Common examples include strep throat, urinary tract infections, bacterial pneumonia, skin boils, and food poisoning; severe infections can spread to the blood and cause life-threatening sepsis. Unlike viral infections, bacterial infections can be treated and cured with antibiotics, but taking antibiotics correctly is essential to prevent bacteria from becoming resistant to medicines.

\synonyms
Bacterial Disease; Bacterial Illness; Bacteriosis

\medterm{Bacterial Meningitis} A serious infection of the membranes around the brain and spinal cord. It can cause fever, severe headache, stiff neck, vomiting, confusion, seizures, or a rash and requires immediate hospital treatment.

\medterm{Bacterial Metabolism} The chemical reactions bacteria use to obtain energy, build cell material, grow, and survive. Different bacteria use oxygen, fermentation, light, or other substances, and some of these processes are targets for antibiotics.

\medterm{Bacterial Morphology} The size, shape, and arrangement of bacteria seen under a microscope. Common forms are spheres, rods, and spirals; bacteria may occur singly, in pairs, chains, clusters, or other patterns.

\medterm{Bacterial Overgrowth} Small-intestinal bacterial overgrowth is an excessive number of bacteria in the small intestine, where relatively few normally live. It can cause bloating, abdominal pain, diarrhoea, weight loss, or vitamin deficiencies and may result from poor intestinal movement or structural changes.

\synonyms
Small Intestinal Bacterial Overgrowth; SIBO

\medterm{Bacterial Pathogenesis} The ways bacteria enter the body, attach to cells, invade tissues, avoid immune defenses, and cause illness. Damage may result from direct invasion, toxins, inflammation, or disruption of normal body functions.

\medterm{Bacterial Peritonitis} A bacterial infection of the lining of the abdominal cavity. It may follow bowel perforation, surgery, another abdominal infection, or infection of abdominal fluid and can cause severe pain, fever, vomiting, and sepsis.

\medterm{Bacterial Pharyngitis} Bacterial infection and inflammation of the throat, often caused by group A Streptococcus. It may cause sudden sore throat, fever, painful swallowing, swollen neck glands, and sometimes white patches on the tonsils.

\medterm{Bacterial Pili} Tiny hair-like structures on the surface of some bacteria. They help bacteria attach to body cells or surfaces; some also help them move or take up DNA from their surroundings.

\medterm{Bacterial Pneumonia} A bacterial infection of the lungs that causes inflammation and may fill air spaces with fluid or pus. It can cause fever, cough, phlegm, chest pain, shortness of breath, and low oxygen levels.

\medterm{Bacterial Prostatitis} Acute or chronic bacterial infection of the prostate gland, typically caused by uropathogenic Gram-negative bacilli (such as \textit{Escherichia coli}), presenting with pelvic pain, dysuria, urinary frequency, fever, and tender prostate.

\synonyms
Acute Bacterial Prostatitis; Chronic Bacterial Prostatitis; Infectious Prostatitis

\medterm{Bacterial Protein} A protein produced by a bacterium that may form part of its structure, support its metabolism, help it attach to host cells, or contribute to disease.

\medterm{Bacterial Spore} A dormant, highly resistant form made by some bacteria when conditions are harsh. Spores can survive drying, heat, and many disinfectants and may later grow into active bacteria when conditions improve.

\medterm{Bacterial Toxemia} The presence or effects of harmful bacterial substances in the bloodstream. It can occur during severe infection and may cause widespread inflammation, low blood pressure, shock, or organ injury.

\medterm{Bacterial Toxin} A harmful substance made or released by bacteria. It may kill or injure cells, interfere with body functions, or trigger severe inflammation; toxins may be secreted or released when bacteria break apart.

\medterm{Bacterial Toxins} Harmful substances produced or released by bacteria. Some are secreted proteins that target particular cells, while others are parts of certain bacterial outer membranes that can trigger widespread inflammation when released.

\medterm{Bacterial Tracheitis} A bacterial infection and swelling of the windpipe. It can cause fever, a harsh cough, noisy or difficult breathing, and rapid narrowing of the airway, making urgent treatment necessary.

\medterm{Bacterial Transduction} The transfer of DNA from one bacterium to another by a virus that infects bacteria. The transferred DNA may give the receiving bacterium new traits, including toxin production or antibiotic resistance.

\medterm{Bacterial Transformation} The uptake of loose DNA from the environment by a bacterium able to absorb it. The new DNA may change the bacterium's features, including how it causes disease or responds to antibiotics.

\medterm{Bacterial Vaginosis} A common change in the vaginal bacteria in which protective lactobacilli decrease and other bacteria increase. It may cause a thin, grey or white discharge and a fish-like smell, but many people have no symptoms.

\medterm{Bacterial Virulence Factors} Features that help bacteria cause disease, such as attaching to cells, entering tissues, avoiding immune defenses, producing toxins, or obtaining nutrients. Examples include surface proteins, protective coatings, enzymes, and toxins.

\medterm{Bactericide} A substance that kills bacteria rather than only stopping them from multiplying. Whether it works depends on the organism, medicine, dose, infection site, and the person's immune defenses.

\medterm{Bacteriocin} A protein-like substance made by some bacteria that inhibits or kills other susceptible bacteria nearby. It gives the producing bacterium a competitive advantage, but its medical importance depends on the organism and setting.

\medterm{Bacteriology} The study of bacteria, including their structure, genes, growth, spread, effects on health, and role in disease. It includes laboratory methods used to identify bacteria and test antibiotic sensitivity.

\medterm{Bacteriophage} A virus that infects bacteria and uses bacterial cells to make more copies of itself. Some phages destroy the bacteria they infect and are being studied for selected antibiotic-resistant infections.

\medterm{Bacteriophage Lambda} A virus that infects the bacterium Escherichia coli. It may multiply and burst the bacterial cell or insert its genetic material into the cell’s DNA and remain inactive until later.

\medterm{Bacteriophage T4} A virus that infects Escherichia coli, makes many copies inside the bacterium, and usually breaks the cell open to release them.

\medterm{Bacteriophages} Viruses that infect bacteria and use them to make more viruses. They may destroy the bacterial cell or insert their genetic material into it and remain inactive for a time.

\medterm{Bacteriorhodopsin} A light-sensitive protein in some microorganisms. It uses light energy to move charged particles across the cell membrane, helping the organism produce energy.

\medterm{Bacteriostasis} The stopping or slowing of bacterial growth without necessarily killing the bacteria. The immune system can then remove the remaining organisms; this differs from an agent that directly kills bacteria.

\medterm{Bacteriostatic} Describes a substance that stops or slows bacteria from multiplying without directly killing them. The body's immune system can then help clear the remaining bacteria.

\medterm{Bacterium} A single bacterial cell. It usually lacks a nucleus, has genetic material in the cell interior, and may have an outer cell wall; some bacteria are helpful, while others cause infection.

\medterm{Bacteriuria} The presence of bacteria in urine detected by testing. It may occur with urinary symptoms as part of an infection or without symptoms; the result must be interpreted with the sample quality, amount of bacteria, pregnancy, and other medical factors.

\medterm{Bacteroidaceae} A family of bacteria that commonly live in the intestine and grow without oxygen. They are usually harmless in the gut but can cause abscesses or other infections if they enter wounds or body areas that should be free of bacteria.

\medterm{Bacteroides} A group of bacteria that normally live in the intestine and grow without oxygen. They are usually harmless there but can cause abscesses or bloodstream infection if bowel contents enter wounds or other sterile areas.

\medterm{Bacteroides Caccae} A bacterium that normally lives in the intestine without oxygen. It is usually harmless there but may be found in mixed abdominal or bloodstream infections when intestinal bacteria escape into other tissues.

\medterm{Bacteroides Dorei} A bacterium that normally lives in the human intestine and grows without oxygen. It is generally considered a harmless member of the gut microbiota.

\medterm{Bacteroides Fragilis} A bacterium that normally lives in the intestine without oxygen. If bowel contents leak into the abdomen, it is an important cause of abscesses, abdominal infection, bloodstream infection, and sepsis.

\medterm{Bacteroides Gingivalis} An older name for Porphyromonas gingivalis, a bacterium that grows without oxygen in dental plaque. It is associated with inflammation and destruction of the tissues supporting the teeth.

\medterm{Bacteroides Ovatus} An anaerobic Gram-negative bacterium that is part of the intestinal microbiota and can contribute to opportunistic infection if it enters damaged or normally sterile tissue.

\medterm{Bacteroides Thetaiotaomicron} An anaerobic Gram-negative bacterium that is an important member of the human gut microbiota and helps break down complex dietary carbohydrates; it can cause opportunistic infection outside the intestine.

\medterm{Bacteroides Ureolyticus} A bacterial name historically applied to an anaerobic, urease-producing organism now generally classified as Campylobacter ureolyticus, which can be associated with wound or bloodstream infection.

\medterm{Bacteroides Vulgatus} An anaerobic Gram-negative bacterium commonly found in the human intestine that can contribute to infection when the intestinal barrier is disrupted.

\medterm{Bad Breath} An unpleasant odor from the mouth, medically called halitosis. It most often results from bacteria, food debris, gum disease, tooth decay, dry mouth, or tonsil problems; less commonly, it may be related to nose, throat, digestive, or other medical conditions.

\synonyms
Halitosis; Oral Malodor; Fetor Oris

\medterm{BAER} Alternate name; see \textit{Brainstem Auditory Evoked Response}.

\medterm{Bag-Valve Mask} A hand-operated resuscitation device, also called a manual resuscitator, consisting of a self-inflating bag, one-way valves, a face mask or airway connector, and often an oxygen reservoir. It delivers positive-pressure ventilation and may be used with a face mask, supraglottic airway, or endotracheal tube.

\synonyms
BVM
Manual Resuscitator
Bag-Valve-Mask Device

\medterm{Bag-Valve-Mask Ventilation} A fundamental emergency airway and resuscitation procedure that delivers manual positive-pressure breathing support to a patient who has stopped breathing (apnea) or is breathing inadequately. A rescuer seals a cushioned anatomical mask tightly over the patient's nose and mouth using the specialized ``E-C clamp'' hand technique while maintaining an open airway, and rhythmically squeezes a self-inflating bag connected to high-flow oxygen. Each gentle squeeze forces oxygenated air past the vocal cords into the lungs, causing the chest to visibly rise. Bag-valve-mask ventilation is the essential first-line rescue bridge in emergency life support and CPR, sustaining vital oxygenation and carbon dioxide elimination until spontaneous breathing returns or an advanced artificial airway (such as an endotracheal tube) is secured.

\synonyms
BVM Ventilation; Bag-Mask Ventilation; Ambu Bag Ventilation; Manual Resuscitator Ventilation

\medterm{Bagassosis} An allergic lung inflammation caused by breathing dust from mouldy bagasse, the fibrous waste left after sugarcane juice is extracted. It may cause cough, fever, breathlessness, and fatigue after workplace exposure.

\medterm{Bags Under Eyes} Mild puffiness or loose skin beneath the eyes. It commonly develops with aging, inherited facial features, lack of sleep, allergies, or fluid retention and is usually harmless.

\synonyms
Puffy Eyes

\medterm{Bainbridge Reflex} An increase in heart rate that can occur when extra blood stretches the upper chambers of the heart. It helps the heart move increased venous blood forward.

\medterm{Bak Protein} A protein on the outer membrane of mitochondria, the cell’s energy-producing structures, that can help start controlled cell death. When activated, it makes openings that release cell-death signals; abnormal control may contribute to cancer.

\medterm{Baker Cyst} A fluid-filled swelling behind the knee caused by fluid from the knee joint collecting in a small sac. It often occurs with arthritis, a meniscus injury, or another condition causing knee inflammation.

\synonyms
Popliteal Cyst; Baker's Cyst

\medterm{Baker's Cyst} A fluid-filled lump behind the knee, usually caused by excess fluid from an irritated knee joint. It may cause tightness or pain and can occasionally rupture, producing calf swelling that can resemble a blood clot.

\medterm{Baker's Cyst} Alternate spelling; see \textit{Baker Cyst}.

\medterm{Baking Powder Overdose} Illness after swallowing a large amount of baking powder. Depending on its ingredients and amount, it may cause vomiting, abdominal swelling, excess sodium, disturbed body salts, or an abnormal blood-alkali level.

\medterm{Baking Soda Overdose} Harm caused by swallowing too much sodium bicarbonate. It can cause vomiting, abdominal swelling, excess sodium, disturbed body salts, abnormal blood acidity, seizures, or heart complications and requires prompt medical treatment.

\medterm{Balamuthia} A genus of free-living amoebas found in soil and water. Some species can rarely infect people, usually causing slowly progressive inflammation of the brain and sometimes skin disease.

\medterm{Balamuthia Mandrillaris} A free-living amoeba found in soil and water that can rarely infect humans. It may enter through broken skin or the nose and cause severe, gradually worsening infection of the brain or skin.

\medterm{Balamuthia Mandrillaris Infection} A rare, usually fatal infection caused by the free-living amoeba Balamuthia mandrillaris. It may
begin as a skin lesion and progress to granulomatous amoebic encephalitis with headache, confusion, seizures, or
focal neurologic deficits. Diagnosis is difficult and requires specialized testing; urgent infectious-disease and
neurologic evaluation is essential.

\medterm{Balance} The ability to keep the body steady while sitting, standing, or moving. It depends on information from the inner ear, eyes, muscles, joints, and brain; problems with any of these systems can increase the risk of falls.

\medterm{Balance Assessment} An evaluation of steadiness during sitting, standing, walking, and changing position. It may include movement tests and review of vision, medicines, strength, sensation, and dizziness to identify fall risk and guide support.

\medterm{Balance Disorders} Problems with steadiness or coordinated movement that may cause dizziness, unsteadiness, or falls. Causes include inner-ear disease, vision problems, brain or nerve disorders, muscle weakness, medicines, and blood-pressure changes.

\medterm{Balance Disorders And Vertigo} Conditions that cause unsteadiness, dizziness, or a false feeling of movement. Causes may involve the inner ear, brain, nerves, vision, muscles, medicines, or blood pressure and require assessment when severe or new.

\medterm{Balance Impairment} Reduced ability to remain steady while sitting, standing, or walking. It may cause swaying, stumbling, or falls and can result from inner-ear disease, vision problems, weakness, nerve disease, medicines, or brain disorders.

\medterm{Balance Problems} Difficulty maintaining body stability while standing or moving, often caused by disorders of the inner ear, nervous system, vision, or muscles.

\medterm{Balance Training} Exercises that improve steadiness and reduce falls by practising safe posture, weight shifting, walking, and responses to movement. A clinician or therapist can adapt exercises to a person's strength, dizziness, and medical conditions.

\medterm{Balanced Study} A clinical study design in which participant characteristics are evenly distributed between comparison groups at baseline, either by design or by randomization, to reduce confounding and allow fair comparison of interventions.

\medterm{Balanitis} Inflammation of the glans penis, often involving the prepuce (balanoposthitis). Causes
include poor hygiene, irritants, infection, and dermatologic conditions such as lichen
planus, psoriasis, and lichen sclerosus. It may cause erythema, swelling, discharge,
pruritus, and discomfort; treatment depends on the cause.

\medterm{Balanitis Xerotica Obliterans} A long-term inflammatory skin disorder affecting the foreskin and head of the penis. It may cause pale tight skin, itching, pain, painful erections, difficulty passing urine, and narrowing of the opening; it needs medical treatment and follow-up.

\synonyms
BZO; Penile Lichen Sclerosus

\medterm{Balanitis, Non-Infectious / Infectious} Alternate index label; see \textit{Balanitis}.

\medterm{Balanoposthitis} Inflammation of both the glans penis and the prepuce. It may cause pain, erythema,
swelling, discharge, or difficulty retracting the foreskin and can result from
infection, irritation, or inflammatory skin disease.

\medterm{Balantidiasis} Balantidiasis is intestinal infection caused by the protozoan Balantidium coli, producing diarrhea, abdominal pain, or colitis, especially after fecal exposure.

\medterm{Balantidium} A genus of large, single-celled parasites that can infect the large intestine. Infection may cause diarrhoea, abdominal pain, or inflammation of the colon, although many infected people have no symptoms.

\medterm{Balantidium Coli} A single-celled parasite that can infect the large intestine, usually after contact with food or water contaminated by animal or human stool. It may cause diarrhoea, abdominal pain, or colitis.

\medterm{Balbiani Body} A temporary cluster of cell parts and genetic messages near the nucleus of a developing egg cell. It helps organize materials that the early embryo needs after fertilization.

\medterm{Baldness} Commonly used term for substantial hair loss, especially androgenetic alopecia, also called pattern
hair loss. In this condition, genetically susceptible hair follicles gradually become smaller under the influence of
androgens.

\synonyms
Pattern Hair Loss

\medterm{Balint's Syndrome} A rare disorder caused by damage to both back portions of the brain. It can cause difficulty seeing more than one object at a time, reaching accurately for objects while looking at them, and directing visual attention.

\medterm{Ballance Sign} A physical examination finding characterized by fixed dullness to percussion in the left flank and shifting dullness in the right flank. It indicates a ruptured spleen with a localized left subcapsular or perisplenic hematoma and free blood in the peritoneal cavity.

\medterm{Ballard Score} A newborn examination used to estimate how many weeks of pregnancy had passed at birth. It assesses physical features and muscle tone or movement; the estimate is useful when pregnancy dates or early ultrasound information are uncertain.

\medterm{Ballism} A movement disorder causing sudden, involuntary, large, flinging movements of an arm or leg. It usually results from damage or dysfunction in deep movement-control areas of the brain.

\medterm{Balloon Angioplasty} A procedure in which a small balloon on a catheter is inflated inside a narrowed artery to improve blood flow. A stent may be placed to help keep the artery open.

\medterm{Balloon Aortic Valvuloplasty} A catheter procedure in which a balloon is inflated across a narrowed aortic valve to improve blood flow from the heart. The benefit may be temporary, and the valve can narrow again; some people later need valve replacement.

\medterm{Balloon Brachytherapy} A form of internal radiation treatment in which a balloon catheter is placed in a body space, often the cavity left after breast-cancer surgery, and radioactive material is temporarily put inside it to treat nearby tissue.

\medterm{Balloon Catheter} A thin flexible tube with an inflatable tip. It can widen a narrowed blood vessel or body passage, open a stent, stop bleeding by applying pressure, or help perform another procedure.

\medterm{Balloon Catheter, Vascular} An inflatable balloon mounted at the distal tip of a flexible vascular catheter used to dilate stenotic or occluded blood vessels during angioplasty, deliver and expand endovascular stents, or temporarily occlude vascular flow.

\synonyms
Angioplasty Balloon; Vascular Balloon Dilator; Balloon Catheter

\medterm{Balloon Dilation} A procedure that widens a narrowed blood vessel or body passage by inflating a small balloon at the narrowed area. The result and risks depend on the site and underlying cause.

\synonyms
Balloon Dilatation; Balloon Angioplasty When Performed in a Blood Vessel

\medterm{Balloon Occlusion} Temporary blockage of a blood vessel or other passage using an inflatable balloon on a catheter. It may control bleeding, test blood flow, or support another procedure; possible complications include reduced blood supply, clotting, and vessel injury.

\medterm{Balloon Tamponade} Temporary control of bleeding or blockage by inflating a balloon inside or against a body passage or blood vessel. It applies pressure to the bleeding area but can injure tissue and requires close monitoring.

\medterm{Balneology} The study of mineral waters, bathing, and spa treatments and their possible health effects. These treatments may ease some symptoms but do not replace proven care, and heat or water exposure can be unsafe for some people.

\medterm{Balneotherapy} Treatment using mineral or thermal water, baths, muds, or related spa methods. It may relieve symptoms in selected conditions, but evidence varies and heat, infection, falls, or heart strain can make it unsafe for some people.

\medterm{Balsam} A thick, resin-containing substance from certain plants, used in some fragrances, skin products, and traditional remedies. It can irritate the skin or cause allergic contact dermatitis in sensitive people.

\medterm{Band Cell} An immature neutrophil, a white blood cell that helps fight infection, whose nucleus has a curved band shape. Increased numbers may occur when the bone marrow responds to infection, inflammation, or severe stress.

\medterm{Band Keratopathy} A horizontal strip of calcium deposits in the clear front surface of the eye. It may result from long-term inflammation or high calcium levels and can cause irritation, blurred vision, or reduced sight.

\medterm{Bandage} Material placed over or around a body part to protect a wound, absorb fluid, apply pressure, or support an injured area. It may be made of gauze, adhesive material, elastic fabric, or a tubular sleeve.

\medterm{BANF} Alternate name; see \textit{Neurofibromatosis Type 2}.

\medterm{Bankart Lesion} A tear in the front-lower rim of the shoulder socket, usually after the upper arm has dislocated. It can make the shoulder unstable and increase the chance of repeated dislocation.

\medterm{Bankart Repair} An open or arthroscopic surgical procedure performed to reattach and reconstruct a torn anterior-inferior glenoid labrum and stretched joint capsule (Bankart lesion) to the glenoid rim to restore joint stability following recurrent anterior shoulder dislocations.

\medterm{Bannayan-Riley-Ruvalcaba Syndrome} A rare inherited condition that can cause many noncancerous growths, a large head size, intestinal polyps, extra freckles around the penis, and increased risk of certain cancers. It results from changes in the PTEN gene.

\medterm{Banti Syndrome} An older term for an enlarged spleen and low blood-cell counts caused by high pressure in the veins that carry blood through the liver. Modern evaluation looks for the liver, vein, or blood disorder responsible.

\synonyms
Banti Disease; Banti's Syndrome

\medterm{Bar Graph} A chart that represents values with rectangular bars whose lengths or heights correspond to the quantities being compared.

\medterm{Barakat Syndrome} An inherited disorder that can cause low parathyroid hormone, hearing loss, and kidney abnormalities. These problems may occur together because of genetic changes affecting the development of several organs.

\medterm{Barber Itch} An inflammatory infection of hair follicles in the beard area, often caused by bacteria or dermatophyte fungi and associated with itchy bumps or pustules.

\synonyms
Barber's Itch

\medterm{Barber Say Syndrome} A rare congenital disorder caused by changes in the TWIST2 gene, characterized by excessive facial hair growth (hypertrichosis), loose, atrophic skin, a wide mouth, distinctive facial features, and eyelid abnormalities such as ectropion.

\medterm{Barbiturate} A group of medicines that slow brain activity and can cause relaxation, sleep, anaesthesia, or seizure control. They can suppress breathing, cause dependence, and be fatal in overdose, so use requires careful medical supervision.

\medterm{Barbiturate Intoxication And Overdose} Poisoning caused by too much of a barbiturate medicine. It can cause extreme sleepiness, confusion, slow or stopped breathing, low blood pressure, coma, or death and requires emergency treatment.

\medterm{Barbiturates} Medicines that reduce activity in the brain and spinal cord. They may be used for anaesthesia or selected seizure disorders, but their risk of breathing suppression, dependence, and overdose means safer alternatives are often preferred.

\medterm{Bardet-Biedl Syndrome} A rare inherited disorder that may cause retinal degeneration, obesity, extra fingers or toes, kidney disease, and developmental or reproductive problems.

\synonyms
BBS

\medterm{Bare-Metal Stent} An expandable metal mesh scaffold without an antiproliferative drug coating, deployed via endovascular catheterization to maintain vascular patency in coronary or peripheral arteries following balloon angioplasty.

\synonyms
BMS; Uncoated Vascular Stent; Bare Stent

\medterm{Bariatric Embolization} A specialized catheter procedure that reduces blood flow to selected stomach arteries to alter hunger-related signaling and support weight loss. It remains limited or investigational in many settings and may cause bleeding, stomach injury, or other vascular complications.

\medterm{Bariatric Surgery} Operations that reduce stomach capacity or change digestion to support substantial weight loss and improve obesity-related disease. They require long-term changes in eating, supplements, follow-up, and monitoring for complications.

\medterm{Bariatrics} The medical field focused on preventing and treating obesity and related conditions. Care may include nutrition, physical activity, psychological support, medicines, and weight-loss surgery.

\medterm{Barium} A metallic element used medically mainly as barium sulfate, a substance that outlines the digestive tract during certain X-ray examinations. Soluble barium compounds are poisonous and are different from barium sulfate.

\medterm{Barium Enema} An X-ray examination of the large intestine after barium is put into the rectum. Sometimes air is added to show the bowel lining more clearly. It may be used to examine blockage, diverticula, or changes after surgery; other tests are now preferred for routine colon-cancer screening.

\medterm{Barium Meal} An X-ray examination in which swallowed barium coats the swallowing tube, stomach, and first part of the small intestine. It can show narrowing, ulcers, abnormal movement, or other structural changes.

\medterm{Barium Study} An X-ray examination that uses barium to outline parts of the digestive tract. Examples include a barium swallow, barium enema, and upper-GI series. It can show the shape and movement of the digestive tract.

\medterm{Barium Sulfate} A liquid or paste swallowed or placed in the rectum to outline the digestive tract during X-ray examinations. It is not absorbed normally, but breathing it into the lungs or leakage through a bowel hole can cause serious complications.

\medterm{Barium Swallow} An X-ray examination of the swallowing tube after a person drinks barium. It can show narrowing, abnormal shape, or problems with movement. A barium follow-through continues the examination into the small intestine.

\medterm{Barium-133} A radioactive form of barium used mainly to check and calibrate radiation detectors and imaging equipment. It is not ordinarily used as a diagnostic scan or treatment for patients.

\medterm{Barlow Maneuver} A clinical physical examination test used to detect developmental dysplasia of the hip in infants. The examiner adducts the flexed hip and applies gentle posterior pressure along the shaft of the femur; a palpable clunk or sensation of the femoral head slipping out of the acetabulum indicates hip instability.

\synonyms
Barlow Test; Barlow Dislocation Test

\medterm{Barlow Sign} A check for hip problems in newborns where the doctor gently presses the baby's bent thigh backward; feeling the hip slip out of place shows the joint is unstable or loose.

\medterm{Barlow Test} A physical examination maneuver for detecting hip instability in infants. With the hip flexed, gentle pressure is applied to see whether the femoral head can be displaced out of the socket, suggesting developmental dysplasia of the hip.

\medterm{Barlow's Disease} Barlow's disease is an older name for severe vitamin-C deficiency, or scurvy, causing bleeding gums, bruising, poor wound healing, and bone or joint symptoms.

\medterm{Barognosis} The ability to recognize and compare the weight of objects by holding them. Losing this ability despite normal basic touch sensation may indicate a problem in brain areas that interpret sensory information.

\medterm{Baroreceptor} A pressure sensor in the walls of major arteries, especially in the neck and chest. It detects blood-pressure changes and signals the brain to adjust heart rate and blood-vessel narrowing.

\medterm{Baroreceptor Reflex} A rapid automatic response that helps keep blood pressure stable. Pressure sensors in major arteries signal the brain to change heart rate and blood-vessel tightness when blood pressure rises or falls.

\medterm{Baroreceptors} Stretch-sensitive nerve endings in major arteries that detect changes in blood pressure. They send signals to the brain, which adjusts heart rate and blood-vessel narrowing to help maintain stable pressure.

\medterm{Baroreflex} A rapid reflex that helps regulate blood pressure by adjusting heart rate and blood-vessel tone in response to changes in arterial pressure.

\medterm{Barotitis Media} Alternate name; see \textit{Airplane Ear}.

\medterm{Barotrauma} Injury caused by a sudden or excessive difference in pressure, affecting the lungs, ears, sinuses, or other air-filled spaces. Lung barotrauma can cause a collapsed lung or air escaping into tissues.

\medterm{Barr Body} The tightly packed, inactive copy of an X chromosome seen inside the nucleus of many cells that contain more than one X chromosome. It helps balance gene activity between cells with different numbers of X chromosomes.

\medterm{Barraquer-Simons Disease} A rare acquired disorder in which fat gradually disappears from the face, arms, and upper body while remaining or increasing around the hips and legs. It may be associated with immune disease and kidney complications.

\medterm{Barrel Chest} A rounded, enlarged appearance of the chest caused by increased lung volume or changes in the chest wall.
It may occur with chronic obstructive pulmonary disease, severe asthma, or other long-standing
lung conditions.

\medterm{Barrett Esophagus} A change in the lower lining of the swallowing tube caused by long-term acid reflux. The changed lining can increase the risk of precancerous changes and cancer.

\medterm{Barrett Esophagus Surveillance} Planned repeat examinations of the lower swallowing tube, usually with an endoscope and tissue samples, to look for precancerous cell changes or early cancer. Timing depends on whether abnormal cell changes are present and their severity.

\medterm{Barrett's Esophagus} Alternate spelling; see \textit{Barrett Esophagus}.
Columnar Epithelium Lined Lower Esophagus

\medterm{Barrett's Esophagus} Alternate spelling; see \textit{Barrett Esophagus}.

\medterm{Barrett's Oesophagus} Alternate spelling; see \textit{Barrett Esophagus}.

\medterm{Barrier Creams} Skin products that form a protective layer against moisture, friction, irritants, and body fluids. They can help prevent or soothe irritation, but infected, broken, or worsening skin may need medical assessment.

\medterm{Barrier Nursing} Infection control and nursing protocols designed to prevent the transmission of pathogenic microorganisms between patients, healthcare personnel, and visitors through strict source isolation, rigorous hand hygiene, and personal protective equipment.

\synonyms
Isolation Nursing; Barrier Precautions; Transmission-Based Precautions

\medterm{Bart's Hemoglobinopathy} A severe inherited blood disorder caused when all four genes needed to make alpha-globin are missing or inactive. The fetus cannot make enough normal haemoglobin, causing severe anaemia, swelling, and often death before or soon after birth.

\medterm{Barth Syndrome} A rare inherited disorder of mitochondrial energy production caused by changes in the TAZ gene. It may cause heart-muscle disease, low muscle tone, weakness, poor growth, low neutrophil counts, and recurrent infections.

\medterm{Barthel Index} A scoring tool that measures how much help a person needs with daily activities such as eating, bathing, dressing, using the toilet, moving between surfaces, walking, and climbing stairs. It is used to track function and rehabilitation progress.

\medterm{Bartholin Abscess} A painful collection of pus caused by infection of a Bartholin gland near the vaginal opening. It can cause severe swelling and pain when sitting or walking and usually needs drainage; antibiotics may be needed in selected cases.

\synonyms
Infected Bartholin Cyst; Bartholin Gland Abscess

\medterm{Bartholin Cyst Or Abscess} A blocked gland near the vaginal opening can cause a painless fluid-filled cyst; infection can produce a painful collection of pus. A new lump in someone aged 40 or older needs specialist assessment.

\medterm{Bartholin Gland} One of two small glands beside the vaginal opening that release lubricating fluid. A blocked duct can cause a cyst, and infection can cause a painful abscess with swelling and difficulty sitting or walking.

\medterm{Bartholin Gland Cyst} A mucus-filled retention cyst resulting from obstruction of the main duct of a greater vestibular (Bartholin) gland located in the posterolateral labium majus. While typically painless when small, secondary infection leads to acute Bartholin abscess formation requiring drainage or marsupialization.

\synonyms
Bartholin Cyst; Bartholin Duct Cyst

\medterm{Bartholin Gland Hyperplasia} Enlargement caused by an increase in the cells of a Bartholin gland. It may appear as a lump near the vaginal opening and can resemble a cyst or tumour, so a persistent or unusual mass needs assessment.

\medterm{Bartholin Gland Mucinous Cyst} A mucus-filled lump near the vaginal opening caused by blockage of a Bartholin gland duct. It is often painless but can become infected and develop into an abscess.

\medterm{Bartholin's Cyst} Alternate spelling; see \textit{Bartholin Cyst}.

\synonyms
Bartholin Gland Cyst; Bartholin Cyst

\medterm{Bartholin's Glands} Alternate spelling; see \textit{Bartholin Gland}.

\medterm{Bartholinitis} Inflammation or infection of a Bartholin gland near the vaginal opening. It may cause pain, redness, swelling, or an abscess and can make sitting or walking uncomfortable.

\medterm{Bartonella} A group of bacteria usually spread by fleas, lice, ticks, or contact with animals. Different species can cause cat-scratch disease, trench fever, skin lesions, heart-valve infection, or bloodstream infection.

\medterm{Bartonella Elizabethae} A species of Bartonella bacteria that has rarely been linked to bloodstream infection and inflammation of the heart valves, especially in people with underlying health problems.

\medterm{Bartonella Henselae} The bacterium that most commonly causes cat-scratch disease. It is usually spread by a kitten scratch or bite and may cause a small skin bump, swollen nearby lymph nodes, fever, or, rarely, disease of other organs.

\medterm{Bartonella Infections} Infections caused by Bartonella bacteria. They may produce swollen lymph nodes, fever, skin lesions, anaemia, heart-valve infection, or inflammation of organs; symptoms and severity depend on the species and immune status.

\medterm{Bartonella Quintana} A Bartonella species spread mainly by body lice. It causes trench fever, with recurring fever, headache, bone pain, and weakness, and can sometimes cause bloodstream or heart-valve infection.

\medterm{Bartonellaceae} A family of bacteria that includes Bartonella species. Members can cause cat-scratch disease, trench fever, skin and blood-vessel lesions, bloodstream infection, or infection of the heart valves.

\medterm{Bartonellosis} Infection caused by a Bartonella bacterium. Depending on the species, it may cause swollen lymph nodes, recurring fever, skin lesions, anaemia, heart-valve infection, or disease affecting several organs.

\medterm{Bartter Disease} Alternate name; see \textit{Bartter Syndrome}.

\medterm{Bartter Syndrome} An inherited kidney disorder causing excessive loss of salt and potassium in urine. It may lead to dehydration, low blood pressure, muscle cramps or weakness, frequent urination, excessive thirst, and poor growth in children.

\medterm{Basal} Relating to the base of a structure or to a person's resting level of activity, such as resting metabolism or temperature measured after sleep.

\medterm{Basal Body} A small structure at the base of a cilium or flagellum, such as a cell hair or sperm tail. It anchors the structure to the cell and helps organize the parts that produce movement.

\medterm{Basal Cell} A cell in the deepest layer of an epithelial tissue, where it can divide and produce new cells. In the epidermis, basal cells help renew the skin and contain keratinocyte stem and progenitor cells.

\medterm{Basal Cell Carcinoma} The most common malignant skin tumor, arising from basal keratinocyte-lineage cells and
associated with cumulative ultraviolet exposure and Hedgehog-pathway abnormalities. It
commonly presents as a pearly papule or nodule with telangiectasia, ulceration, or a
rolled border.

\medterm{Basal Cell Skin Cancer} Alternate name; see \textit{Basal Cell Carcinoma}.

\medterm{Basal Ganglia} Deep brain structures that help control movement, habits, learning, and emotions. Problems in these structures can cause tremor, stiffness, slowed movement, or unwanted movements.

\medterm{Basal Ganglia Dysfunction} Abnormal function of deep brain structures that help control movement and learned habits. It may cause tremor, stiffness, slowed movement, twisting movements, or other involuntary movements, as in Parkinson disease or Huntington disease.

\medterm{Basal Joint Arthritis} Alternate name; see \textit{Thumb Arthritis}.

\medterm{Basal Lamina} A thin layer beneath many lining cells that anchors them to supporting tissue. It also helps control movement of substances between the lining and the tissue below.

\medterm{Basal Metabolic Rate} The amount of energy the body uses at complete rest to maintain essential functions such as breathing, circulation, temperature, and cell activity. It varies with body size, age, muscle mass, hormones, illness, and other factors.

\medterm{Basal Nuclei} Groups of nerve cells deep within the brain that help start, stop, and control movement. They also contribute to learning, habits, motivation, and other thinking processes.

\medterm{Basal Tear} A tear that is continuously present on the eye's surface. It keeps the cornea moist, smooth, and protected and helps wash away dust and microorganisms.

\medterm{Basal Temperature} Body temperature measured immediately after waking and before eating, exercise, or other activity. A small sustained rise often occurs after ovulation and can help track menstrual cycles.

\medterm{Base} A substance that can accept hydrogen ions or reduce acidity. In water, a base usually has a pH above 7; strong bases can burn tissue, while weaker bases are important in body fluids and chemical reactions.

\medterm{Base Excision Repair} A way cells repair small areas of DNA damage. Enzymes remove the faulty DNA building block, fill the gap with the correct one, and seal the DNA strand.

\medterm{Base Pair} Two matching building blocks joined inside DNA or RNA. In DNA, adenine pairs with thymine and cytosine pairs with guanine; in RNA, adenine usually pairs with uracil.

\medterm{Base Pairing} The matching of complementary DNA or RNA building blocks. This matching helps copy genetic information accurately and allows RNA molecules to fold into useful shapes or attach to other molecules.

\medterm{Base Sequence} The exact order of DNA or RNA building blocks in a molecule. It carries genetic information and can be examined to identify inherited changes, mutations, or infectious organisms.

\medterm{Basedow's Disease} Older name; see \textit{Graves disease}.

\medterm{Baseline} A starting measurement or description of a person's health before treatment, illness, or observation. Later findings are compared with it to identify improvement, worsening, or change.

\medterm{Baseline EKG} An electrocardiogram recorded before treatment, a procedure, or a change in health. It provides a starting record for comparison with later heart tracings.

\synonyms
Baseline ECG; Reference Electrocardiogram

\medterm{Basement Membrane} A thin supporting layer beneath many cells that line organs and body surfaces. It anchors the cells, helps control exchange with underlying tissue, and can act as a barrier to tumour spread.

\medterm{Basic First Aid} Immediate help given for illness or injury until professional care arrives. It may include making the area safe, calling emergency services, controlling bleeding, helping a choking person, starting resuscitation, and monitoring breathing.

\medterm{Basic Life Support} Emergency care for a person who is unresponsive and not breathing normally. It includes calling for help, chest compressions, rescue breaths when appropriate, and use of an automated defibrillator if advised.

\medterm{Basic Metabolic Panel} A routine blood test measuring eight vital substances: blood sugar (glucose), calcium, four key electrolytes (sodium, potassium, chloride, and bicarbonate), and two waste products filtered by the kidneys (BUN and creatinine). It provides a rapid overview of fluid and electrolyte balance, kidney function, blood sugar control, and blood acid-base status.
\synonyms
BMP; Chem-7; Chemistry Screen; Electrolyte and Kidney Panel

\medterm{Basic Reproduction Number} The average number of people one infected person would infect in a population where everyone is susceptible, under specified conditions. It helps describe spread but changes with behavior, immunity, treatment, and the environment.

\medterm{Basidiomycota} A major fungal group that reproduces by forming spores on club-shaped structures called basidia. It includes mushrooms, rusts, smuts, and some fungi that can cause human disease.

\medterm{Basilar} Relating to the base or lowest part of a body structure. For example, a basilar finding may be near the bottom of the skull, lung, or another organ.

\medterm{Basilar Artery} The artery formed by the union of the vertebral arteries that supplies the brainstem, cerebellum, and posterior brain.

\medterm{Basilar Artery Aneurysm} A weak, widened area in the artery supplying the brainstem and back of the brain. It may press on nearby nerves, cause double vision or weakness, or rupture and produce life-threatening bleeding around the brain.

\medterm{Basilar Membrane} A flexible membrane in the cochlea that supports the organ of Corti and vibrates in response to sound. Different parts respond best to different sound frequencies.

\medterm{Basilic Vein} A superficial vein of the upper limb that ascends along the medial side of the arm and joins the brachial veins to form the axillary vein. It is used for venous access in selected settings.

\medterm{Basiliximab} A monoclonal antibody that blocks the interleukin-2 receptor alpha chain on activated T cells. It is used as induction immunosuppression, especially in kidney transplantation, and increases susceptibility to infection.

\medterm{Basis Pontis} The basis pontis is the front portion of the pons containing descending motor fibers and pontine nuclei that relay signals to the cerebellum.

\medterm{Basket Cell} A basket cell is an inhibitory interneuron that forms synaptic contacts around the cell body or axon initial segment of another neuron.

\medterm{Basophil} A rare type of white blood cell involved in allergic and inflammatory reactions. It releases substances such as histamine and can increase in number with some allergies, infections, and blood disorders.

\medterm{Basophil Adenoma} A usually noncancerous tumour of hormone-producing cells in the pituitary gland. It may release too much hormone or press on nearby structures, causing headache, vision problems, or other hormone-related symptoms.

\medterm{Basophilia} An unusually high number of basophils in the blood. It may occur with allergies, inflammation, infection, an underactive thyroid, or disorders in which the bone marrow makes too many blood cells.

\medterm{Basophilic Focus} An area in a tissue sample that stains blue or purple with a laboratory dye. The finding describes appearance under the microscope and must be interpreted with the surrounding tissue and other tests.

\medterm{Basophilic Stippling} The abnormal presence of multiple fine or coarse, dark-blue granules scattered throughout erythrocytes on a Romanowsky-stained peripheral blood smear. The granules represent precipitated ribosomal ribonucleic acid (RNA) and are characteristically seen in lead poisoning, thalassemias, sideroblastic anemia, and myelodysplastic syndromes.

\synonyms
Punctate Basophilia

\medterm{Basosquamous Carcinoma} A skin cancer with features of both basal-cell and squamous-cell carcinoma. It usually develops on sun-damaged skin and may behave more aggressively than ordinary basal-cell cancer, so complete treatment and follow-up are important.

\medterm{Bassen-Kornzweig Syndrome} A rare inherited disorder in which the body cannot make certain particles that carry fat in the blood. It causes poor absorption of dietary fat, deficiencies of fat-soluble vitamins, unusual red cells, and progressive nerve or eye problems.

\medterm{Bateman Purpura} Alternate name; see \textit{Actinic Purpura}.

\medterm{Bath Salts Abuse And Addiction} Harmful use of synthetic stimulant drugs sold as “bath salts.” They can cause agitation, paranoia, hallucinations, seizures, dangerous overheating, heart problems, dependence, or death; severe symptoms require emergency care.

\medterm{Bather's Itch} An itchy rash caused by tiny parasite larvae released by infected snails in fresh or salt water. The larvae cannot mature in people, so the rash usually settles on its own but can be intensely uncomfortable.

\medterm{Bathing A Patient In Bed} Washing a person who cannot safely get out of bed. Care includes keeping water comfortably warm, protecting privacy and skin, preventing falls, checking for pressure injuries, and helping only as much as needed.

\medterm{Bathing An Infant} Washing a baby safely in shallow water while keeping one hand on the baby at all times. The water is comfortably warm, the room warm, and all supplies are ready before bathing; a baby is never left alone near water.

\medterm{Batrachotoxin} A very powerful poison found in some frogs and birds. It keeps nerve and heart sodium channels open, causing abnormal nerve activity, dangerous heart rhythms, paralysis, and potentially death.

\medterm{Batroxobin} An enzyme derived from snake venom that changes fibrinogen, a blood protein, into fibrin, the material that strengthens a blood clot. It has been studied or used in selected bleeding and clotting settings and can alter bleeding risk.

\medterm{Batson Plexus} A valveless network of epidural and paravertebral veins that connects the deep pelvic and thoracic veins to the internal vertebral venous system. Because of the lack of valves and low intravascular pressures, it provides an alternate route for the retrograde hematogenous spread of pelvic infections and metastases (especially prostate and breast carcinoma) to the spine and central nervous system.

\synonyms
Batson's Venous Plexus; Vertebral Venous Plexus

\medterm{Batten Disease} A group of inherited disorders in which waste material builds up in nerve cells. They usually begin in childhood and cause worsening seizures, vision loss, movement problems, learning difficulties, and loss of abilities.

\medterm{Battered Child Syndrome} An older term for repeated physical abuse of a child by a caregiver. Possible warning signs include injuries of different ages, unexplained bruises or fractures, head injury, inconsistent explanations, or delayed medical care; suspected abuse requires immediate safeguarding action.

\medterm{Battey Bacillus} An older name for the group of nontuberculous mycobacteria now called Mycobacterium avium complex. These organisms can cause lung disease, swollen lymph nodes, or widespread infection, especially when immunity is severely weakened.

\medterm{Battle Sign} Bruising behind the ear over the bone, which is a warning sign of a break in the base of the skull.

\medterm{Bauhin Valve} The valve between the last part of the small intestine and the first part of the large intestine. It helps regulate passage of intestinal contents and limits backward movement from the colon.

\medterm{Bauxite-Workers' Lung} A lung disorder caused by long-term workplace inhalation of bauxite dust. It may produce scarring or inflammation in the lungs and can cause cough, breathlessness, or reduced lung function.

\medterm{Bax Protein} A protein that helps start programmed cell death when a cell is damaged or no longer needed. Excessive or insufficient Bax activity can affect cell survival and has been studied in cancer and other diseases.

\medterm{Bayes Theorem} A rule for updating the likelihood of a diagnosis after new evidence. It combines the chance of the disease before the test with how accurately the test result supports or argues against that disease.

\medterm{Bazin Disease} A chronic, nodular, and often ulcerating form of lobular panniculitis predominantly affecting the posterior lower legs of young or middle-aged women. It represents a delayed-type hypersensitivity reaction associated with Mycobacterium tuberculosis or other systemic antigens.

\synonyms
Erythema Induratum of Bazin; Erythema Induratum

\medterm{Bazin's Disease} An older name for tender, inflamed lumps in the fatty tissue under the skin of the lower legs. It is sometimes linked with tuberculosis but can have other causes and may ulcerate as it heals.

\medterm{Bcgosis} An infection caused by the weakened mycobacteria used in BCG treatment or vaccination. It may remain local or spread to organs, especially in people with weakened immunity, and can cause fever, cough, or inflammation.

\medterm{Beach Chair Position} A semi-sitting position used during some operations, especially shoulder surgery. The head, neck, and pressure points must be supported, and blood pressure and brain blood flow require careful monitoring.

\medterm{Beacopp Regimen} A combination chemotherapy regimen used mainly to treat Hodgkin lymphoma; it includes bleomycin, etoposide, doxorubicin, cyclophosphamide, vincristine, procarbazine, and prednisone, with dosing and supportive care determined by the oncology team.

\medterm{Beals Syndrome} A rare inherited connective-tissue disorder present from birth. It may cause permanently bent joints, long slender fingers and toes, a curved spine, and ears with a crumpled appearance; severity varies.

\medterm{Beam Hardening Artifact} A false dark band or streak on a CT image caused when bone or metal removes some parts of the X-ray beam before it reaches the detector. It can obscure nearby body structures.

\medterm{Beat Number} The number assigned to a particular heartbeat or rhythmic event when beats are counted or analysed. It can help describe timing, rate, or changes in a recorded heart rhythm.

\medterm{Beats, Cardiac} Alternate index label; see \textit{Heartbeat}.

\medterm{Beau Lines} Transverse depressions or grooves across the nail plate caused by temporary interruption
of nail-matrix growth. They may follow severe illness, fever, trauma, malnutrition,
chemotherapy, or systemic inflammatory disease.

\medterm{Beau's Lines} Transverse grooves or depressions across a fingernail or toenail caused by temporary interruption of
nail-matrix growth. They may follow severe illness, injury, high fever, nutritional stress, or
local trauma; the timing of the groove can help estimate when the event occurred.

\medterm{Becaplermin} Becaplermin is a recombinant platelet-derived growth factor gel used in selected diabetic foot ulcers with appropriate wound care.

\medterm{Bechterew Disease} An older name for ankylosing spondylitis, a long-term inflammatory disease affecting the joints between the spine and pelvis and often the spine itself. It causes back pain and stiffness that may improve with movement.

\synonyms
Ankylosing Spondylitis; Marie-Strümpell Disease

\medterm{Beck Triad} Three findings that may occur when fluid or blood compresses the heart: low blood pressure, swollen neck veins, and quiet heart sounds. The triad may be absent, so suspected compression needs urgent examination and imaging.

\medterm{Becker Muscular Dystrophy} An inherited muscle disease caused by reduced or abnormal dystrophin, a protein that supports muscle fibres. It usually causes slowly progressive weakness beginning in childhood or adulthood and may also affect the heart.

\medterm{Becker Myotonia} An inherited muscle condition causing stiffness because muscles relax slowly after contraction. It often begins in childhood, may improve after repeated movement, and can cause muscle enlargement, difficulty starting movement, or brief weakness.

\synonyms
Autosomal Recessive Myotonia Congenita; Recessive Generalized Myotonia

\medterm{Becker Nevus} A usually harmless, irregularly shaped patch of darker skin that often becomes noticeable during adolescence and may develop increased hair growth. It commonly affects one shoulder or the upper trunk and is not usually cancerous.

\medterm{Becker's Muscular Dystrophy} An inherited muscle disease causing progressive weakness, usually beginning later and advancing more slowly than Duchenne muscular dystrophy. It can also weaken the heart and requires regular heart assessment.

\medterm{Becker’S Nevus} A usually harmless darkened patch of skin, often with extra hair, that commonly appears on one shoulder or the upper trunk during adolescence. It is not cancer, but changes or uncertainty require assessment.

\medterm{Beckwith-Wiedemann Syndrome} An overgrowth condition present from birth that may cause a large body or tongue, abdominal-wall defects, enlarged organs, low blood sugar in newborns, and increased risk of certain childhood cancers. Regular specialist monitoring is needed.

\medterm{Beclomethasone} Beclomethasone is a corticosteroid used by inhalation, nasal spray, or topical routes to reduce inflammation in asthma, rhinitis, or skin disease.

\medterm{Becquerel} The unit used to measure radioactivity. One becquerel means that one radioactive atom changes each second.

\medterm{Becquerel} The International System (SI) derived unit of radionuclide activity, defined as one nuclear decay or disintegration per second ($1\text{ Bq} = 1\text{ s}^{-1}$). Clinical diagnostic and therapeutic radioisotope doses are typically quantified in megabecquerels (MBq) or gigabecquerels (GBq).

\synonyms
Bq; SI Unit of Radioactivity

\medterm{Bed Rest} Staying in bed or remaining mostly reclined to limit activity. It is used only for selected medical reasons because prolonged inactivity can cause muscle loss, blood clots, constipation, pressure injuries, and reduced independence.

\medterm{Bed Rest During Pregnancy} Restricting activity or staying in bed during pregnancy for a specific medical reason. Routine bed rest is generally avoided because it has not prevented many pregnancy complications and can increase muscle loss and blood-clot risk.

\medterm{Bedbug} A small, flat, blood-feeding insect that hides in bedding, furniture, and cracks. Its bites can cause itchy skin bumps or allergic reactions, but bedbugs are not known to routinely spread human disease.

\medterm{Bedbug Bites} Itchy skin reactions caused by bedbugs feeding at night. Bites often appear as red bumps in clusters or lines on exposed skin, sometimes days after exposure; scratching can cause skin infection.

\medterm{Bedbug Bites} Alternate name; see \textit{Bed Bugs}.

\medterm{Bedbugs} Small blood-feeding insects that cause itchy bites and skin reactions and can spread through bedding, clothing, and furniture.

\medterm{Bednar Aphthae} Shallow sores on the back of a baby's palate caused by pressure or rubbing, often during feeding. They may make feeding painful and usually improve when the source of irritation is corrected.

\medterm{Bedpan} A container used by a person who cannot safely get out of bed to urinate or pass stool. It is positioned securely, privacy is maintained, and the skin is cleaned and dried afterward.

\medterm{Bedsore} An area of skin and deeper-tissue damage caused by prolonged pressure, usually over a bony area such as the heel, hip, or tailbone. Limited movement, moisture, poor nutrition, and reduced sensation increase risk.

\medterm{Bedsores} Localized ischemic injury to the skin and underlying subcutaneous tissue, muscle, or bone resulting from prolonged unrelieved pressure, shear forces, and friction over bony prominences in immobilized patients.

\synonyms
Pressure Ulcers; Decubitus Ulcers; Pressure Sores; Decubiti

\medterm{Bedtime Habits For Infants And Children} Regular bedtime routines that help children sleep safely and consistently. They may include a predictable wind-down period, suitable sleep conditions, and age-appropriate limits on screens; persistent sleep problems need assessment.

\medterm{Bedwetting} Alternate name; see \textit{Enuresis}.

\medterm{Bee And Wasp Sting} A sting from a bee or wasp that injects venom and usually causes immediate pain, redness, swelling, and itching. Trouble breathing, throat swelling, widespread hives, faintness, or vomiting may indicate anaphylaxis and requires emergency treatment.

\medterm{Bee Sting} A venom-injecting injury from a bee that usually causes local pain, redness, and swelling. A severe allergic reaction can cause widespread hives, breathing difficulty, faintness, or swelling of the mouth and throat.

\synonyms
Apis Envenomation; Honeybee Sting

\medterm{Bee Stings} Injuries caused by bees injecting venom into the skin. They may cause local pain and swelling, a large reaction extending beyond the sting, or life-threatening anaphylaxis in an allergic person.

\medterm{Beef or Pork Tapeworm Infection} Intestinal infection with Taenia saginata or Taenia solium acquired from undercooked beef or pork; T. solium eggs can also cause cysticercosis in tissues.
\medterm{Beer-Drinker's Heart} Alternate name; see \textit{Alcoholic Cardiomyopathy}.

\medterm{Beers Criteria} A regularly updated list of medicines that may cause extra harm in some older adults, such as confusion, falls, or bleeding. It supports medication review but does not replace individualized clinical judgment.

\medterm{Beeswax Poisoning} Illness after ingesting or inhaling a large amount of beeswax. Beeswax is generally of low toxicity, but large ingestions may cause gastrointestinal blockage or aspiration, especially in children.

\medterm{Behavior} The observable actions, responses, and activities of an organism, influenced by genetic, neurobiological, psychological, and environmental factors. In medicine, behavioral assessment helps diagnose psychiatric, developmental, and neurological conditions.

\medterm{Behavior Disorder} A persistent pattern of actions or emotional responses that causes significant distress, interferes with daily life, or harms relationships. Causes may be developmental, psychological, medical, or environmental.

\medterm{Behavior Test} A structured assessment of behavior, mood, attention, thinking, or social functioning. It may use observation, interviews, questionnaires, or standardized tasks to help identify developmental, psychiatric, or neurologic concerns.

\medterm{Behavior Therapy} A form of psychotherapy based on learning principles that uses reinforcement, extinction, and other behavioral techniques to modify maladaptive behaviors. It is effective for anxiety disorders, phobias, obsessive-compulsive disorder, and habit disorders.

\medterm{Behavioral Abnormality} A deviation from typical patterns of behavior that may indicate a psychiatric, neurodevelopmental, or neurologic condition. It is assessed by history, observation, standardized measures, and developmental norms to determine clinical significance.

\medterm{Behavioral Disorder} A persistent pattern of behavior that causes clinically significant distress or impairment in daily functioning, relationships, learning, or safety. The term is broad; interpretation considers age, development, mental health, medical conditions, environment, trauma, and substance exposure.

\medterm{Behavioral Genetics} The study of how genes and environmental factors influence behavior and mental traits. Most behaviors reflect many genes interacting with life experiences, not one gene alone.

\medterm{Behavioral Rehabilitation} Structured therapy that helps a person change behaviors, build coping skills, and regain daily function after illness, injury, addiction, or disability. It may include counseling, goal setting, and practice of useful routines.

\medterm{Behcet Disease} A long-term inflammatory illness that repeatedly causes painful mouth sores and genital sores. It may also inflame the eyes, skin, joints, blood vessels, brain, or digestive tract and can threaten vision or other organs.

\medterm{Behcet Syndrome} Another name for Behcet disease, an inflammatory illness that causes repeated mouth and genital sores and may affect the eyes, skin, joints, blood vessels, brain, or digestive tract.

\medterm{Behcet's Syndrome} Behcet syndrome is a relapsing inflammatory disorder that can cause recurrent mouth and genital ulcers, eye inflammation, skin lesions, and vascular, neurologic, or gastrointestinal disease. Treatment is tailored to the organs involved and may include immunosuppression.

\synonyms
Behçet Disease; Silk Road Disease; Adamantiades-Behçet Syndrome

\medterm{Behr Syndrome} A rare inherited neurological disorder causing progressive optic-nerve damage, poor vision, muscle stiffness, unsteady movement, and sometimes weakness or bladder problems.

\medterm{Behçet Disease} Standard spelling of Behcet disease, a relapsing inflammatory disorder that can affect the mouth, genitals, eyes, skin, joints, blood vessels, nervous system, or digestive tract.

\medterm{Behçet's Syndrome} Behçet's syndrome is a relapsing inflammatory disorder causing recurrent oral and genital ulcers, eye inflammation, skin lesions, and sometimes vascular, neurologic, or gastrointestinal disease.

\medterm{Beijerinck, Martinus W.} A Dutch microbiologist who helped establish the concept of viruses as infectious agents that reproduce only in living cells and contributed to the development of enrichment culture methods.

\medterm{Bejel} Bejel is endemic treponematosis, a chronic infection caused by Treponema pallidum pertenue or related treponemes and spread mainly by close nonsexual contact in endemic areas.

\medterm{Bekam} A Malay name for cupping, a practice in which cups create suction on the skin; wet cupping may also involve superficial skin incisions.

\medterm{Belching} The release of swallowed air or gas from the stomach through the mouth. Frequent
belching may occur with aerophagia, gastroesophageal reflux, functional dyspepsia, or
carbonated-drink consumption.

\medterm{Belinostat} A cancer medicine that blocks enzymes involved in turning genes on and off. It is used for some peripheral T-cell lymphomas that have returned or resisted earlier treatment.

\medterm{Bell Palsy} Sudden weakness or paralysis of one side of the face caused by dysfunction of the facial nerve. The eyelid may not close fully, and taste or hearing may change. Most cases improve, but urgent assessment is needed to distinguish it from stroke and to protect the eye.

\synonyms
Bell's Palsy

\medterm{Bell's Palsy} Acute peripheral facial-nerve weakness of uncertain cause, usually affecting one side of the face. Eye protection and timely clinical assessment are important; treatment is individualized.

\synonyms
Idiopathic Facial Palsy; Facial Nerve Palsy

\medterm{Bell's Phenomenon} Upward and outward movement of the eyes when a person closes the eyelids forcefully. This normal protective reflex helps shield the cornea and may become visible when facial-nerve weakness prevents complete eyelid closure.

\medterm{Bell-Clapper Deformity} A birth-related difference in how a testicle is attached inside the scrotum, allowing it to move and twist more freely than usual. It is often present on both sides and increases the risk of testicular torsion, a sudden emergency that cuts off blood flow.

\synonyms
Bell-Clapper Anomaly; High Tunica Vaginalis Insertion

\medterm{Belladonna} Atropa belladonna, the deadly nightshade plant, whose tropane alkaloids such as atropine and scopolamine can cause anticholinergic poisoning.

\medterm{Bells Palsy} Alternate name; see \textit{Bell's Palsy}.

\medterm{Belly} The abdomen, the region of the body between the chest and pelvis containing much of the digestive, urinary, and reproductive organs.

\medterm{Below-Knee Amputation} Surgery to remove the lower leg while preserving the knee joint. It may be needed for severe injury, infection, poor blood flow, or cancer; rehabilitation focuses on healing, strength, mobility, and a prosthesis when suitable.

\medterm{Bence Jones Protein} An abnormal antibody fragment made by some plasma-cell disorders. It may appear in blood or urine and help identify conditions such as multiple myeloma, but it can also damage the kidneys.

\medterm{Bence-Jones Protein} An abnormal antibody fragment made by a single group of plasma cells. Finding it in urine or blood may support evaluation for multiple myeloma or another plasma-cell disorder and possible kidney injury.

\medterm{Benchmarking} Comparing healthcare results, safety, costs, or processes with a standard or similar service. The comparison can identify areas for improvement but does not by itself prove that one service causes better outcomes.

\medterm{Bendamustine} A cancer medicine that damages tumour-cell DNA and is used for selected blood cancers. It can also damage healthy bone-marrow cells, causing low blood counts, infection, tiredness, nausea, or other side effects.

\medterm{Bendamustine Hydrochloride} The hydrochloride form of bendamustine, a chemotherapy medicine used for selected blood cancers. It damages DNA and may cause low blood counts, infection, infusion reactions, nausea, tiredness, or other side effects.

\medterm{Bendiocarb} A carbamate insecticide that can overstimulate nerves after significant exposure. Poisoning may cause sweating, salivation, vomiting, diarrhoea, small pupils, muscle twitching, breathing difficulty, or weakness and requires urgent medical treatment.

\medterm{Bendroflumethiazide} A diuretic medicine that helps the kidneys remove salt and water. It is used for high blood pressure or fluid retention and may lower potassium or sodium and increase uric acid or blood sugar.

\medterm{Bends} A form of decompression sickness caused by gas bubbles forming in the body when pressure falls too quickly, usually during or after diving. It commonly causes deep joint or muscle pain but can also injure the brain, spinal cord, lungs, or circulation.

\medterm{Benedikt Syndrome} A brainstem vascular stroke syndrome resulting from a paramedian lesion of the midbrain tegmentum involving the oculomotor nerve fascicles, red nucleus, and substantia nigra. It is clinically characterized by ipsilateral third cranial nerve palsy (ptosis, dilated pupil, downward-and-outward eye) and contralateral involuntary choreoathetosis, intention tremor, and ataxia.

\synonyms
Paramedian Midbrain Syndrome

\medterm{Benedikt’s Syndrome} A rare pattern of neurological problems caused by damage in the midbrain. It may cause a drooping eyelid or double vision on one side and tremor, poor coordination, or involuntary movements on the opposite side.

\medterm{Beneficence And Non-Maleficence} Two principles of medical ethics: beneficence means acting to benefit the patient, while non-maleficence means avoiding unnecessary harm. Good decisions balance expected benefits, risks, patient values, and available alternatives.

\medterm{Benefits And Risks Of Taking Amino Acid Supplement} Amino-acid supplements may help when a documented deficiency or specific medical need exists, but most people obtain enough from food. Unnecessary or excessive use can cause stomach symptoms, interactions, or metabolic, kidney, or liver problems.

\medterm{Benfotiamine} A fat-soluble form of thiamine, or vitamin B1, sold in some supplements and studied for diabetic nerve damage. Evidence of benefit is limited and it does not replace proven treatment for diabetes or neuropathy.

\medterm{Benfuracarb} A carbamate insecticide that can overstimulate nerves after significant exposure. It may cause sweating, salivation, vomiting, diarrhoea, small pupils, muscle twitching, breathing difficulty, or weakness and requires urgent medical treatment.

\medterm{Benign} Not cancerous. A benign growth does not invade nearby tissues or spread to other parts of the body, but it can still cause problems because of its size, location, or hormone production.

\synonyms
Nonmalignant; Noncancerous

\medterm{Benign Adrenal Tumors} Noncancerous growths in an adrenal gland. Some are found by chance and cause no problems, while others release excess hormones or become large enough to cause pressure symptoms.

\medterm{Benign Ameloblastoma} A usually slow-growing tumour arising from cells involved in tooth development, most often in the jaw. It is not cancer but can expand, damage bone, and displace teeth, so treatment usually requires specialist surgery.

\medterm{Benign Brain Tumor} A noncancerous growth in or around the brain. Although it does not spread like cancer, it can cause headache, seizures, vision or movement problems, or fluid blockage by pressing on nearby structures.

\medterm{Benign Chondroblastoma} A noncancerous cartilage-producing bone tumour that usually develops near the end of a long bone in an adolescent or young adult. It often causes pain, swelling, or reduced movement near a joint.

\medterm{Benign CNS Cyst} A noncancerous, fluid-filled sac in or around the brain or spinal cord. Many cause no symptoms, but a large or strategically located cyst can press on nerves or block the normal flow of brain fluid.

\medterm{Benign Ear Cyst Or Tumor} A noncancerous lump or fluid-filled sac in the outer, middle, or inner ear. It may cause no symptoms or lead to pain, discharge, hearing loss, dizziness, or pressure depending on its location and size.

\medterm{Benign Epithelioma} A noncancerous growth arising from cells that line the skin or an organ. It does not invade or spread, but removal or testing may be needed if it bleeds, grows, causes symptoms, or is difficult to distinguish from cancer.

\medterm{Benign Esophageal Stricture} A noncancerous narrowing of the esophageal lumen, most commonly resulting from chronic acid reflux injury (peptic stricture), corrosive ingestion, radiation, or prior surgery, causing progressive dysphagia.

\synonyms
Peptic Stricture; Esophageal Stenosis, Benign; Esophageal Stricture

\medterm{Benign Hemangiopericytoma} An older term for a usually slow-growing tumour with branching blood vessels, now often classified as a solitary fibrous tumour. Behaviour varies, so tissue examination and follow-up are important.

\medterm{Benign Intracranial Hypertension} Historical and misleading name; see \textit{Idiopathic Intracranial Hypertension}.

\medterm{Benign Liver Tumors} Noncancerous growths in the liver, including hemangiomas, focal nodular hyperplasia, and hepatocellular adenomas. Most are found by chance on ultrasound, CT, or MRI and do not require treatment. Some may cause pain, liver enlargement, bleeding, or rarely change into cancer.

\medterm{Benign Melanocytic Lesion} An older or imprecise term for a benign melanocytic growth, such as a melanocytic nevus or mole; unlike malignant melanoma, it is not cancerous.

\medterm{Benign Meningioma} A usually slow-growing, noncancerous tumour arising from the membranes around the brain or spinal cord. It may cause seizures, headache, weakness, or vision problems by pressing on nearby structures.

\medterm{Benign Mesenchymoma} A rare noncancerous tumour containing more than one type of tissue such as fat, cartilage, or bone. Its symptoms depend on the site and size, and diagnosis requires examination of tissue.

\medterm{Benign Migratory Glossitis} Alternate name; see \textit{Geographic Tongue}.

\medterm{Benign Mixed Tumor} A usually noncancerous salivary-gland tumour containing more than one tissue type. It often presents as a slow-growing, painless lump, but long-standing or recurrent growths require assessment because rare malignant change can occur.

\medterm{Benign Mullerian Papilloma} A rare noncancerous, wart-like growth in the vulva or vagina formed from Müllerian-type lining cells. It often appears as a small polyp and is diagnosed by examining a tissue sample.

\medterm{Benign Myoepithelioma} A noncancerous tumour of myoepithelial cells, often arising in a salivary gland or other glandular tissue. It usually grows slowly and may cause a painless lump; diagnosis depends on tissue examination.

\medterm{Benign Neoplasm} A noncancerous growth of cells that remains at its original site and does not spread to distant organs. It may still cause harm by pressing on nearby structures, blocking a passage, bleeding, or producing hormones.

\medterm{Benign Orgasmic Headache} A sudden severe headache during sexual activity or orgasm after bleeding, stroke, vessel disease, and other dangerous causes have been excluded. A first or unusually severe headache needs urgent medical assessment.

\synonyms
Primary Headache Associated with Sexual Activity; Sexual-Activity Headache

\medterm{Benign Paraganglioma} A usually slow-growing tumour of specialized nerve-related cells. It may release adrenaline-like hormones or press on nearby nerves and blood vessels; symptoms and treatment depend on its site and activity.

\medterm{Benign Paroxysmal Positional Vertigo} A common inner-ear disorder causing brief spinning attacks when the head changes position. Tiny calcium particles move into a balance canal and trigger dizziness, often lasting less than a minute.

\medterm{Benign Peripheral Nerve Tumor} A noncancerous tumour arising from or beside a nerve outside the brain and spinal cord. It may cause a lump, pain, tingling, numbness, weakness, or no symptoms at all.

\medterm{Benign Positional Vertigo} Alternate name; see \textit{Benign Paroxysmal Positional Vertigo}.

\medterm{Benign Prostatic Hyperplasia} Noncancerous enlargement of the prostate that can narrow the urine channel. It may cause a weak stream, difficulty starting, urgency, frequent night-time urination, dribbling, or incomplete emptying.

\medterm{Benign Prostatic Obstruction} Blockage at the bladder outlet caused by an enlarged prostate or increased muscle tone in the prostate. It may cause a weak or interrupted stream, straining, incomplete emptying, urinary retention, or recurrent infection.

\medterm{Benign Rolandic Epilepsy} A childhood seizure disorder, now often called self-limited epilepsy with centrotemporal spikes. Brief seizures commonly affect the face or mouth during sleep and may cause twitching, tingling, drooling, or temporary trouble speaking; it usually ends by adolescence.

\synonyms
Self-Limited Epilepsy with Centrotemporal Spikes; SeLECTS; Benign Childhood Epilepsy with Centrotemporal Spikes; Rolandic Epilepsy

\medterm{Benign Tumor} A noncancerous growth that remains localized and does not invade nearby tissue or spread to distant organs. It may still cause problems by pressing on structures, blocking a passage, bleeding, or making hormones.

\medterm{Benign Urethral Lesions} Noncancerous growths or structural changes in the urethra, such as polyps, caruncles, papillomas, or inflammatory lesions. They may cause bleeding, discharge, pain, narrowing, or difficulty passing urine, but some cause no symptoms.

\synonyms
Urethral Polyps; Urethral Caruncle; Benign Urethral Tumors

\medterm{Bennett Fracture} An intra-articular fracture-subluxation of the base of the first metacarpal bone. The smaller medial articular fragment remains attached to the trapezium by the anterior oblique ligament, while the main metacarpal shaft is displaced proximally, laterally, and dorsally by the pulling force of the abductor pollicis longus tendon.

\medterm{Benserazide} A medicine combined with levodopa in some Parkinson treatments. It blocks levodopa from being changed into dopamine outside the brain, allowing more to reach the brain and reducing some unwanted effects.

\medterm{Bent Penis} Curvature of the penis that may be present from birth or develop later. It can be painless or cause pain, difficulty with erection or intercourse, and sometimes problems passing urine.

\medterm{Bent Spine Syndrome} Severe forward bending of the trunk that improves when lying down. It results from weakness or abnormal contraction of back muscles and can occur with muscle, nerve, or movement disorders.

\synonyms
Camptocormia

\medterm{Bentall Procedure} Open-heart surgery that replaces a diseased aortic valve, the first part of the aorta, and the aortic root with an artificial valve and tube graft. The coronary arteries are reattached to the graft; lifelong monitoring and sometimes blood-thinning treatment are needed.

\synonyms
Bentall Operation; Composite Aortic Root Replacement

\medterm{Benzaldehyde} A chemical with an almond-like smell used in flavourings, fragrances, and manufacturing. Concentrated exposure can irritate the eyes, skin, or lungs; it is not a routine medical treatment.

\medterm{Benzazepines} A family of chemical structures made from joined benzene and seven-membered nitrogen-containing rings. Individual compounds may have different biological or medicinal effects, so the family name alone does not identify a treatment.

\medterm{Benzene} An industrial chemical and solvent found in some fuels and smoke. High short-term exposure can affect the brain and heart; repeated exposure can damage bone marrow and increase the risk of leukaemia.

\medterm{Benzene Poisoning} Illness caused by breathing, swallowing, or absorbing benzene. Short-term exposure may cause dizziness, drowsiness, or irregular heartbeat; repeated exposure can damage bone marrow and increase the risk of leukaemia.

\medterm{Benzhexol} Benzhexol, or trihexyphenidyl, is an anticholinergic drug used mainly to reduce tremor and rigidity in Parkinson disease or drug-induced extrapyramidal symptoms.

\medterm{Benzo[a]pyrene} A polycyclic aromatic hydrocarbon (PAH) formed during the incomplete combustion of organic matter, including tobacco smoke, vehicle exhaust, and grilled meats. Its active metabolite, benzo[a]pyrene-7,8-dihydrodiol-9,10-epoxide (BPDE), forms mutagenic DNA adducts that drive chemical carcinogenesis.

\synonyms
BaP; 3,4-Benzopyrene; Benzo(a)pyrene

\medterm{Benzoates} Salts or chemical forms of benzoic acid used mainly as food preservatives or ingredients in some products. Their effects and safety depend on the specific compound and amount; some people develop irritation or sensitivity.

\medterm{Benzocaine} A local anaesthetic applied to skin or moist body surfaces to reduce sensation. Rarely, especially after excessive use, it can cause methemoglobinaemia, a condition in which blood cannot carry oxygen normally.

\medterm{Benzodiazepine} A medicine that enhances the effect of the inhibitory chemical messenger GABA in the brain. It may reduce anxiety, induce sleep, relax muscles, or control seizures, but can cause sedation, impaired coordination, dependence, and dangerous breathing suppression when combined with opioids or other sedatives.

\medterm{Benzodiazepine Use Disorder} A pattern of benzodiazepine use that causes loss of control, harm, or problems at work, home, or in relationships. Regular use can cause dependence, and sudden stopping can produce life-threatening withdrawal.

\medterm{Benzodiazepine Withdrawal} Symptoms after reducing or stopping regular benzodiazepine use. Anxiety, insomnia, shaking, sweating, and sensitivity to light may occur; severe withdrawal can cause seizures, hallucinations, or delirium and requires supervised tapering.

\medterm{Benzopyran} A chemical ring structure formed by joining benzene and pyran rings. It occurs in some natural substances and medicines; the term describes chemistry rather than one specific disease or treatment.

\medterm{Benzopyrene} A group of chemicals produced by incomplete burning and found in tobacco smoke, vehicle exhaust, and heavily charred food. Some forms can damage DNA and increase cancer risk after the body changes them chemically.

\medterm{Benzopyrenes} Chemicals with a benzopyrene structure that form during incomplete combustion. Several occur in smoke and pollution, can damage DNA, and are linked with increased cancer risk.

\medterm{Benzoyl Peroxide} A topical medicine that kills acne-causing bacteria and helps remove blocked skin cells. It can cause dryness, peeling, irritation, or bleaching of hair and fabrics.

\medterm{Benzphetamine} A sympathomimetic amine and stimulant formerly used in some countries as a short-term adjunct for weight reduction; it can cause dependence, cardiovascular effects, and insomnia and is restricted or unavailable in many jurisdictions.

\medterm{Benztropine} An anticholinergic medicine used mainly for movement problems caused by some medicines and for acute muscle spasms called dystonia. It can cause dry mouth, blurred vision, constipation, urinary retention, confusion, or worsen glaucoma.

\medterm{Benzydamine} A topical anti-inflammatory pain-relieving medicine used for painful inflammation in the mouth, throat, or skin. It may cause temporary numbness, irritation, or an allergic reaction.

\medterm{Benzyl Benzoate} A topical chemical used to kill mites causing scabies and, in some settings, lice. It can irritate the skin and is applied as directed, with care to avoid the eyes, mouth, and broken skin.

\medterm{Benzylamines} Chemical compounds containing a benzyl group joined to an amino group. Some are used to make medicines or other chemicals, while their effects depend on the individual compound.

\medterm{Benzylpenicillin} A penicillin antibiotic used for serious infections caused by susceptible bacteria, including some streptococcal infections, meningitis, and syphilis. Allergic reactions can be severe.

\medterm{BEP Regimen} A chemotherapy combination of bleomycin, etoposide, and cisplatin, commonly used for some germ-cell cancers. It can cause low blood counts, infection, kidney damage, hearing changes, nerve symptoms, or lung injury and requires monitoring.

\medterm{Berberine} A plant compound sold in some supplements and studied for blood-sugar, cholesterol, and antimicrobial effects. Evidence varies, and it can cause stomach symptoms, low blood sugar, drug interactions, or harm during pregnancy.

\medterm{Bereavement} The emotional response to the death of someone important. Grief may include sadness, yearning, anger, numbness, sleep or appetite changes, and waves of distress; it varies widely and can sometimes lead to prolonged grief or depression.

\medterm{Bereavement Reaction} Emotional and physical responses after someone dies, such as sadness, disbelief, anger, poor sleep, or loss of appetite. These reactions are common and do not automatically indicate a mental disorder.

\medterm{Berger Disease} An older name for IgA nephropathy, a kidney disease in which an immune protein builds up in kidney filters. It may cause episodes of blood in the urine after infections, protein in the urine, high blood pressure, and kidney failure.

\medterm{Berger's Disease} Alternate name; see \textit{IgA Nephropathy}.

\medterm{Beriberi} Illness caused by severe vitamin B1 deficiency. It may damage nerves and cause weakness, numbness, or difficulty walking, or affect the heart and cause swelling, fast heartbeat, and heart failure; urgent vitamin replacement is needed.

\medterm{Bernard-Soulier Disease} A rare inherited bleeding disorder in which platelets are unusually large and do not stick properly to damaged blood vessels. It causes easy bruising, nosebleeds, heavy periods, and prolonged bleeding after injury or procedures.

\medterm{Bernard–Soulier Syndrome} A rare inherited bleeding disorder in which platelets are unusually large and cannot attach normally to injured blood vessels. It causes easy bruising, nosebleeds, heavy periods, and prolonged bleeding after injury or procedures.

\synonyms
Giant Platelet Disorder; Platelet-Type Bleeding Disorder 1

\medterm{Berry Aneurysm} A small, rounded balloon-like bulge at a branching point of an artery at the base of the brain. If it bursts, it can cause life-threatening bleeding around the brain.

\medterm{Berylliosis} A long-term lung disease caused by sensitivity to inhaled beryllium dust or fumes. It can cause cough, breathlessness, fatigue, and scarring or inflammation that resembles sarcoidosis.

\medterm{Beryllium} A light metal used in some industrial products. Breathing its dust or fumes can cause allergic lung inflammation, long-term lung disease, and increased cancer risk, especially after repeated workplace exposure.

\medterm{Beta 2-Microglobulin} A small protein released by most cells. Blood or urine levels may rise with inflammation, some blood cancers, or reduced kidney removal; the result is interpreted with symptoms and other tests.

\medterm{Beta Amyloid} Small protein fragments made from amyloid precursor protein. They can build up into plaques in the brain and are associated with Alzheimer disease, although plaques alone do not explain every symptom.

\synonyms
A$\beta$ Peptide; Amyloid-Beta Protein

\medterm{Beta Carotene} A yellow-orange plant pigment that the body can convert into vitamin A. It is found in foods such as carrots and leafy vegetables; excessive supplements can turn the skin yellow-orange.

\synonyms
Provitamin A Carotenoid; Carotenoid Provitamin

\medterm{Beta Cell} A cell in the pancreas that makes and releases insulin when blood glucose rises. Insulin helps body cells use glucose and keeps blood glucose from becoming too high.

\medterm{Beta Cells} Pancreatic cells that produce insulin and amylin. Insulin lowers blood glucose by helping cells take up glucose, while amylin slows the rise after meals; loss or failure of these cells contributes to diabetes.

\medterm{Beta Globulin} A group of blood proteins identified together during a protein blood test. They include proteins that carry iron or fats and support immunity; abnormal levels can occur with inflammation, liver disease, or immune disorders.

\medterm{Beta Glucans} Complex carbohydrates found in oats, barley, yeast, fungi, and some bacteria. Oat and barley forms can modestly lower cholesterol when eaten regularly; effects on immunity depend on the source and preparation.

\medterm{Beta Interferons} Immune-modifying medicines used to reduce relapses and new brain or spinal-cord lesions in some forms of multiple sclerosis. Common effects include flu-like symptoms and injection-site reactions; monitoring may detect liver or blood-count abnormalities.

\medterm{Beta Particle} A fast-moving electron or positively charged equivalent released by some radioactive atoms. It can damage nearby tissue and is therefore used carefully in selected radiation treatments and laboratory applications.

\medterm{Beta Radiation} Radiation made of fast electrons or their positively charged counterparts released by unstable atoms. It can injure living tissue but can also be directed at selected tissues for treatment or used in medical testing.

\medterm{Beta Thalassemia} An inherited blood disorder in which the body makes too little beta-globin, a part of haemoglobin. It can cause mild or severe anaemia, tiredness, jaundice, bone changes, or an enlarged spleen.

\medterm{Beta Tocopherol} One naturally occurring form of vitamin E. It acts as a fat-soluble antioxidant, but the body generally uses alpha-tocopherol more effectively; its specific clinical importance is still being studied.

\medterm{Beta-Adrenoceptor-Blocking Drugs} Medicines that block adrenaline-like signals at beta receptors. They can slow the heart, lower blood pressure, reduce tremor, and protect the heart in selected conditions, but may cause fatigue, low pulse, or breathing problems in susceptible people.

\medterm{Beta-Alanine} An amino acid that helps muscles make carnosine, which buffers acid during intense activity. Supplements may modestly improve some short, high-intensity exercise but commonly cause temporary tingling of the skin.

\medterm{Beta-Blockers} Medicines that block adrenaline-like effects, usually slowing the heartbeat and lowering blood pressure. They are used for several heart conditions, migraine prevention, tremor, glaucoma, and other problems; some can worsen asthma or mask low blood sugar.

\medterm{Beta-Blockers Overdose} Poisoning from taking too much of a beta-blocker. It can cause a very slow heartbeat, low blood pressure, wheezing, low blood sugar, confusion, seizures, or shock and requires emergency assessment.

\medterm{Beta-Carotene} A plant pigment that the body can convert into vitamin A. It occurs in foods such as carrots and leafy greens; high-dose supplements may increase lung-cancer risk in smokers and are not a substitute for a balanced diet.

\medterm{Beta-Carotene Blood Test} A blood test measuring beta-carotene, a plant pigment that can become vitamin A. A low result may suggest poor nutrition or fat absorption; a high result may cause harmless yellow-orange skin.

\medterm{Beta-Chemokines} Small immune signals that attract white blood cells to areas of infection or inflammation. They help coordinate defense, but excessive activity can contribute to long-lasting inflammation and tissue damage.

\medterm{Beta-Endorphin} A natural pain-relieving substance made in the brain and pituitary gland. It acts on opioid receptors and also influences stress responses, mood, and some hormone functions.

\medterm{Beta-Glucans} Complex carbohydrates found in foods such as oats, barley, mushrooms, and yeast. Their effects vary by source; some provide fibre and may modestly affect cholesterol, blood glucose, or immune activity.

\medterm{Beta-Haemolytic Streptococcus} A streptococcus that completely breaks down red blood cells around its colonies in a laboratory culture. Important group A and group B species can cause throat, skin, bloodstream, newborn, or other serious infections.

\medterm{Beta-Human Chorionic Gonadotrophin} The measurable part of the pregnancy hormone hCG. Blood or urine testing helps detect pregnancy and can help assess some pregnancy problems or selected tumours when interpreted with other findings.

\medterm{Beta-Lactam Antibiotics} A large group of antibiotics that stop susceptible bacteria building their protective cell wall. They include penicillins, cephalosporins, carbapenems, and related medicines; allergic reactions and bacterial resistance can limit their use.

\medterm{Beta-Lactamase} An enzyme made by some bacteria that breaks down certain penicillin-like antibiotics and makes them ineffective. Laboratory testing can identify this resistance, and some medicines block selected beta-lactamases.

\medterm{Beta-Mannosidosis} A rare inherited disorder in which cells cannot break down certain complex sugars. The stored material can cause developmental delay, hearing loss, repeated infections, seizures, or skeletal and skin problems.

\medterm{Beta-Oxidation} Alternate form; see \textit{Fatty Acid Oxidation}.

\medterm{Beta-Sheet} A common folded shape in proteins in which extended strands lie beside one another and are held by chemical bonds. Excessive beta-sheet folding can contribute to abnormal protein deposits in some diseases.

\medterm{Beta-Thalassemia} An inherited disorder in which the body makes too little or no beta-globin, a part of haemoglobin. Severity ranges from mild carrier status to severe anaemia requiring regular transfusions and monitoring for iron overload.

\medterm{Beta-Thalassemia Intermedia} A moderately severe inherited form of beta-thalassemia causing ongoing anaemia, enlarged spleen, jaundice, or bone changes. People may need occasional transfusions, but usually do not require regular lifelong transfusions.

\medterm{Beta-Thalassemia Minor} The carrier form of beta-thalassemia, caused by one altered beta-globin gene. It usually causes no symptoms or mild small-cell anaemia and is not treated with iron unless iron deficiency is also confirmed.

\synonyms
Beta-Thalassemia Trait; Cooley Trait

\medterm{Beta-Thalassemia Trait} The heterozygous carrier state for a beta-globin gene variant. It is usually asymptomatic or causes mild microcytosis with mild or no anemia; hemoglobin electrophoresis often shows increased HbA2. It is also called beta-thalassemia minor.

\synonyms
Beta-Thalassemia Minor; Thalassemia Minor; Cooley Trait

\medterm{Betacoronavirus} A group of coronaviruses that includes viruses causing respiratory illnesses in people and animals. Individual viruses differ in spread and severity; some cause mild colds, while others can cause severe pneumonia.

\medterm{Betaherpesvirinae} A subgroup of herpes viruses that includes cytomegalovirus and human herpesviruses 6 and 7. They remain in the body after infection and can reactivate, especially when immunity is weakened.

\medterm{Betaine} A natural substance involved in processing homocysteine and helping cells maintain their water balance. It is also sold as a supplement, but high doses can cause side effects or interact with medicines.

\medterm{Betamethasone} A strong corticosteroid medicine that reduces inflammation and immune activity. It may be applied to skin or given in other forms; prolonged or extensive use can thin the skin, worsen infections, or suppress the body's steroid production.

\medterm{Betamethasone Valerate} A corticosteroid applied to the skin to reduce inflammation and itching in conditions such as eczema or psoriasis. Excessive, prolonged, or covered application can thin the skin and increase absorption into the body.

\medterm{Betaxolol} A relatively beta-1-selective blocker used orally for some cardiovascular conditions and topically to lower intraocular pressure in glaucoma. It may cause bradycardia, hypotension, or bronchospasm.

\medterm{Bethanechol} Bethanechol is a muscarinic cholinergic agonist that stimulates bladder and gastrointestinal smooth muscle, used selectively for urinary retention or impaired motility.

\medterm{Bethesda System For Reporting Thyroid Cytopathology} A six-tier standardized system for reporting thyroid fine-needle aspiration results, from nondiagnostic and benign categories to follicular-pattern, suspicious, and malignant categories; each tier guides risk assessment and management.

\medterm{Bethlem Myopathy} A usually slowly progressive inherited muscle disease caused by changes in collagen VI. It causes weakness, often beginning in childhood, together with tight joints or shortened muscles, especially in the fingers, elbows, ankles, or spine; many people remain able to walk into adulthood.

\synonyms
Collagen VI-Related Myopathy (Mild Form); Benign Congenital Myopathy with Contractures

\medterm{Bevacizumab} A monoclonal antibody that blocks vascular endothelial growth factor and inhibits new blood-vessel formation. It is used for selected cancers and some eye diseases; important risks include hypertension, bleeding, thrombosis, proteinuria, impaired wound healing, and gastrointestinal perforation.

\medterm{Bexarotene} An oral retinoid X-receptor agonist used mainly for cutaneous T-cell lymphoma; it commonly causes hypertriglyceridaemia, hypothyroidism, and liver abnormalities and is teratogenic.

\medterm{Bezafibrate} Bezafibrate is a fibrate lipid-lowering drug that activates PPAR-alpha and lowers triglycerides; muscle and liver effects require monitoring.

\medterm{Bezoar} A compact mass of material that cannot be digested and becomes trapped in the stomach or intestine. It may contain hair, plant matter, medicines, or other substances and can cause pain, vomiting, blockage, or bleeding.

\medterm{Bezold-Jarisch Reflex} An inhibitory autonomic cardiovascular reflex characterized by the triad of profound bradycardia, peripheral vasodilation with severe hypotension, and transient apnea. It is triggered by the stimulation of non-myelinated vagal sensory C-fibers in the inferoposterior wall of the left ventricle by mechanical deformation, chemical agents, or acute ischemia during inferior myocardial infarction.

\medterm{Biacetyl} Biacetyl, or diacetyl, is a diketone compound with a buttery odor used in flavoring and industry; inhalational exposure can injure the small airways.

\medterm{Bias} A systematic deviation that distorts a measurement, study result, clinical judgment, or decision away from the truth or from an intended fair comparison.

\medterm{Bible Cyst} Alternate historical name; see \textit{Ganglion Cyst}.

\medterm{Bibliotherapy} The use of reading materials as a therapeutic intervention to support mental health, educate patients about conditions, or promote behavioral change. It may be used alone or alongside other treatments for depression, anxiety, or chronic disease self-management.

\medterm{Bicalutamide} A non-steroidal androgen-receptor antagonist used in prostate cancer, usually with androgen-deprivation therapy; adverse effects include hot flushes, breast enlargement, sexual dysfunction, and possible liver injury.

\medterm{Bicarbonate} A substance in the blood that helps keep the body’s acid level balanced. The lungs and kidneys regulate it, and a blood bicarbonate result can help identify acid–base or metabolic problems.

\medterm{Biceps} The biceps is a two-headed muscle, especially the biceps brachii of the upper arm, that flexes the elbow and supinates the forearm.

\medterm{Biceps Brachii Muscle} The two-part muscle on the front of the upper arm. It bends the elbow, turns the palm upward, and helps lift the arm. Injury or tendon rupture can cause weakness, pain, bruising, or a bulging “Popeye” appearance.

\medterm{Biceps Reflex} Bending of the elbow after the biceps tendon is gently tapped. This reflex tests the nerve supplying the biceps and the related spinal-cord pathways, mainly at the fifth and sixth neck levels.

\medterm{Bicipital} Bicipital means relating to a two-headed muscle, especially the biceps tendon, groove, or muscle of the upper arm.

\medterm{Bicipital Groove} The intertubercular groove of the humerus that transmits the tendon of the long head of
the biceps brachii, with the transverse humeral ligament and adjacent rotator-cuff
structures contributing to tendon stability.

\medterm{Bicipital Tuberosity} The radial tuberosity, a roughened prominence on the proximal radius where the biceps
brachii tendon inserts and produces flexion and supination of the forearm.

\medterm{Bickerstaff Brainstem Encephalitis} A rare autoimmune post-infectious central nervous system disorder closely related to the Miller Fisher variant of Guillain-Barré syndrome. It is characterized by the clinical triad of acute progressive ophthalmoplegia, severe ataxia, and central nervous system disturbance (altered consciousness or hyperreflexia), typically accompanied by serum anti-GQ1b IgG autoantibodies.

\medterm{Bicolor} Having two distinct colours. In medicine, the term may describe a skin lesion, eye, specimen, or imaging finding; its significance depends on the structure and clinical context.

\medterm{Biconcave} Curved inward on both sides. Normal red blood cells have this disc shape, which makes them flexible and provides a large surface area for carrying oxygen.

\medterm{Biconvex} Curved outward on both sides. The term may describe the shape of a lens, implant, anatomical structure, or finding on a scan.

\medterm{Bicornuate} Having two horn-like cavities or projections. A bicornuate uterus is a congenital difference in which the upper uterus forms two partly separated spaces.

\medterm{Bicornuate Uterus} A uterus with two upper cavities caused by incomplete joining during development before birth. Many people have no symptoms, but it can increase the risk of miscarriage, preterm birth, or abnormal baby position.

\medterm{Bicuspid} Having two points or flaps. The term describes some teeth, which have two main points, and some heart valves, which have two flaps instead of the usual three.

\medterm{Bicuspid Aortic Valve} A congenital aortic valve with two leaflets instead of the usual three. It can lead to
aortic stenosis, aortic regurgitation, enlargement of the aorta, or other congenital heart abnormalities.

\medterm{Bicuspid Valve} A heart valve with two leaflets rather than the usual three. The term most often refers
to congenital bicuspid aortic valve, which can cause aortic stenosis, regurgitation,
infective endocarditis, or ascending-aortic disease.

\medterm{Bid Protein} A protein that helps connect signals from outside a cell with programmed cell death inside it. When activated, it can help mitochondria release signals that lead to destruction of a damaged or unwanted cell.

\medterm{Bier Spots} Small, pale macules surrounded by an erythematous background on the extremities, typically elicited when the limbs are placed in a dependent position. They represent a benign physiological vascular anomaly caused by exaggerated local vasoconstriction in response to increased venous pressure and resolve upon limb elevation.

\medterm{Biermer's Anemia} Another name for pernicious anaemia, in which the body cannot absorb enough vitamin B12 because the stomach does not provide the needed absorption factor. It causes anaemia and can damage nerves if untreated.

\medterm{Bietti's Crystalline Dystrophy} A rare inherited eye disease in which tiny crystals and progressive damage develop in the retina and sometimes the cornea. It may cause difficulty seeing at night, loss of side vision, and gradual loss of central vision.

\medterm{Bifascicular Block} A cardiac intraventricular conduction defect in which two of the three main conduction fascicles below the bundle of His are blocked, characteristically presenting as complete right bundle branch block (RBBB) combined with either left anterior fascicular block (LAFB) or left posterior fascicular block (LPFB).

\medterm{Bifocal} Having two points of focus. Bifocal glasses contain one area for distance vision and another for near vision, allowing both prescriptions in one lens.

\medterm{Bifurcation} The point where one body structure, blood vessel, airway, nerve, or pathway divides into two branches.

\medterm{Big Toe Sign} Upward movement of the big toe, sometimes with spreading of the other toes, when the sole is stroked. In an adult, this may suggest a problem in pathways connecting the brain and spinal cord.

\medterm{Bigeminy} A heart-rhythm pattern in which each normal heartbeat is followed by an early extra beat. The extra beat often begins in the lower heart chambers and may cause palpitations or no symptoms.

\medterm{Biglycan} A protein-and-sugar component of the tissue framework around cells. It helps organize collagen and may influence inflammation, wound healing, and blood-vessel changes.

\medterm{Biguanide} A group of blood-glucose-lowering medicines that reduce glucose production by the liver and improve the body's response to insulin. Kidney function, side effects, and other health conditions guide their use.

\medterm{Biguanides} A group of medicines that lower blood glucose mainly by reducing glucose production in the liver and improving insulin sensitivity. They are used chiefly in type 2 diabetes.

\synonyms
Biguanide Drugs; Biguanide Antidiabetics

\medterm{Bilateral} Affecting both sides of the body or both members of a paired structure, such as both eyes, ears, knees, lungs, or kidneys.

\medterm{Bilateral Acoustic Neurofibromatosis} Alternate name; see \textit{Neurofibromatosis Type 2}.

\medterm{Bilateral Ankle Swelling} Alternate form; see \textit{Ankle Swelling}.

\medterm{Bilateral Gynecomastia} Enlargement of glandular breast tissue on both sides in a male. It may result from normal hormone changes, medicines, obesity, liver or kidney disease, or rarely cancer.

\medterm{Bilateral Hydronephrosis} Enlargement of the urine-collecting spaces in both kidneys because urine cannot drain normally. Causes include blockage, reflux, stones, or bladder problems; prolonged obstruction can damage kidney function.

\medterm{Bilateral Microphthalmos} A birth condition in which both eyes are unusually small. Vision may be reduced, and other eye or developmental differences may occur; assessment by an eye specialist is needed.

\medterm{Bilateral Oophorectomy} Surgery to remove both ovaries, sometimes performed for cancer, severe endometriosis, or high inherited cancer risk. Before natural menopause it causes immediate loss of ovarian hormones and infertility, so long-term effects need discussion.

\medterm{Bilateral Peritonsillar Abscesses} Collections of pus beside both tonsils, usually after a throat infection. They can cause severe sore throat, difficulty swallowing, trouble opening the mouth, muffled speech, and airway narrowing, requiring urgent hospital treatment.

\medterm{Bilateral Pneumonia} Infection or inflammation affecting both lungs. It may cause fever, cough, chest discomfort, breathlessness, or low oxygen; severity depends on the cause, extent, and person's health.

\medterm{Bilateral Tonic-Clonic Seizure} A seizure causing loss of awareness followed by stiffening and rhythmic jerking of both sides of the body. First aid aims to protect the person from injury without restraining them, and emergency help is needed if the seizure lasts five minutes or more.

\medterm{Bile} A fluid made by the liver and stored in the gallbladder that helps digest fats. It carries bile salts, cholesterol, pigments, and waste products into the intestine.

\medterm{Bile Acid Sequestrants For Cholesterol} Medicines that trap bile acids in the intestine so they leave in stool. The liver then uses cholesterol to make more bile acids, lowering LDL cholesterol; constipation and medicine interactions are common concerns.

\medterm{Bile Acids} Substances made by the liver and released into the intestine to help digest and absorb fats and vitamins A, D, E, and K. Abnormal blood levels can indicate liver or bile-flow problems.

\medterm{Bile Culture} A laboratory test that grows microorganisms from a bile sample to identify infection in the gallbladder or bile ducts and help select effective treatment.

\medterm{Bile Duct} A tube that carries bile from the liver and gallbladder into the first part of the small intestine. Blockage or inflammation can prevent bile flow and cause pain, jaundice, infection, or pancreatitis.

\medterm{Bile Duct Cancer} Cancer that begins in a bile duct. It can block bile flow and cause yellow skin, dark urine, pale stools, itching, abdominal pain, or weight loss; treatment depends on its location and spread.

\synonyms
Cholangiocarcinoma; Biliary Tract Cancer

\medterm{Bile Duct Carcinoma} Alternate name; see \textit{Bile Duct Cancer}.

\medterm{Bile Duct Obstruction} A blockage that prevents bile moving from the liver or gallbladder into the intestine. It may cause jaundice, dark urine, pale stools, itching, abdominal pain, infection, or pancreatitis.

\medterm{Bile Duct Stricture} Narrowing of a bile duct that slows or blocks bile flow. It may result from surgery, inflammation, stones, or cancer and can cause jaundice, pain, infection, or abnormal liver tests.

\medterm{Bile Ducts} Tubes that carry bile from the liver and gallbladder to the small intestine. A blockage or narrowing can cause jaundice, abdominal pain, infection, or inflammation of the pancreas.

\medterm{Bile Pigment} A coloured substance, mainly bilirubin, made when haemoglobin from old red blood cells is broken down. The liver processes it and excretes it in bile; build-up causes yellowing of the skin and eyes.

\medterm{Bile Reflux} Backward flow of bile and other intestinal contents into the stomach and sometimes the swallowing tube. It may cause upper-abdominal pain, nausea, bitter fluid, or inflammation and can occur after surgery or when normal stomach movement is impaired.

\synonyms
Duodenogastric Reflux; Biliary Gastritis

\medterm{Bile Salt Emulsification} The process in which bile salts break dietary fat into tiny droplets. This increases the surface area available to digestive enzymes and helps the intestine absorb fat and fat-soluble vitamins.

\medterm{Bile Salt Malabsorption} Failure to reabsorb enough bile salts in the lower small intestine. Excess bile salts enter the colon and draw in water, causing chronic watery diarrhoea, urgency, and abdominal discomfort.

\medterm{Bile Salts} Substances made from cholesterol in the liver and released in bile. They help break fats into small droplets so the intestine can digest and absorb fats and vitamins A, D, E, and K.

\medterm{Bilevel Positive Airway Pressure} Breathing support delivered through a mask that gives higher pressure when breathing in and lower pressure when breathing out. It can improve oxygen and carbon-dioxide levels in selected illnesses and requires monitoring.

\medterm{Bilharzia} An infection caused by Schistosoma parasites acquired when contaminated freshwater contacts the skin. It may affect the intestine, liver, urinary tract, or other organs and can cause long-term inflammation and scarring.

\medterm{Bilharzia} Alternate name; see \textit{Schistosomiasis}.

\medterm{Bili Lights} Special blue lights used to treat jaundice in newborns. The light changes bilirubin in the skin into forms the body can remove more easily, helping prevent bilirubin from damaging the brain.

\medterm{Biliary} Relating to bile, the gallbladder, or the tubes that carry bile from the liver to the intestine.

\medterm{Biliary Atresia} A serious disease in newborns in which the bile ducts outside the liver are blocked or absent. Bile builds up in the liver, causing jaundice and progressive liver damage unless treated promptly.

\synonyms
Congenital Biliary Atresia
Extrahepatic Biliary Atresia

\medterm{Biliary Calculi} Stones formed from hardened bile in the gallbladder or bile ducts. They may cause no symptoms or block bile flow, producing pain, jaundice, infection, or inflammation of the gallbladder or pancreas.

\medterm{Biliary Cirrhosis} Permanent liver scarring caused by long-term injury or blockage of the bile ducts. The term has also been used for primary biliary cholangitis, an immune disease affecting small bile ducts.

\medterm{Biliary Cirrhosis, Primary} Alternate name; see \textit{Primary Biliary Cholangitis}.

\medterm{Biliary Colic} Steady pain in the upper right abdomen or upper middle abdomen caused by a gallstone temporarily blocking the gallbladder outlet. It often follows a meal and lasts minutes to hours; fever or persistent pain suggests a complication.

\medterm{Biliary Decompression} A procedure that drains blocked bile ducts to restore bile flow. A doctor may place a stent or drainage tube through an endoscope or through the skin, depending on the cause and location of the blockage.

\medterm{Biliary Dyskinesia} Gallbladder pain without gallstones or another obvious structural problem, often associated with abnormal emptying. Diagnosis requires typical symptoms and exclusion of other causes; a scan may help but is not decisive alone.

\medterm{Biliary Fever} Fever caused by infection or inflammation in the gallbladder or bile ducts. Fever with upper-abdominal pain or jaundice may indicate a rapidly worsening infection and needs urgent medical care.

\medterm{Biliary Fistula} An abnormal passage allowing bile to leak from the gallbladder or bile ducts into another organ, the abdominal cavity, or the skin. It may cause pain, infection, bile drainage, or skin irritation.

\medterm{Biliary Sand} Fine crystals or thick sludge suspended in bile, usually containing cholesterol, calcium salts, or pigment. It may disappear, form gallstones, block a duct, or cause gallbladder inflammation.

\medterm{Biliary Scintigraphy} A specialized nuclear medicine scan that tracks the production and flow of bile from the liver into the gallbladder, bile ducts, and small intestine using an injected radioactive tracer (such as technetium-99m labeled iminodiacetic acid). It is the most accurate imaging test for diagnosing acute cholecystitis (indicated by failure of the gallbladder to fill due to cystic duct blockage). It is also used to detect post-surgical bile leaks, evaluate biliary dyskinesia by measuring gallbladder emptying after cholecystokinin administration, and differentiate biliary atresia from neonatal hepatitis in jaundiced infants.

\synonyms
Cholescintigraphy; HIDA Scan; Hepatobiliary Scintigraphy; Hepatobiliary Iminodiacetic Acid Scan

\medterm{Biliary Stent} A small tube placed inside a bile duct to keep it open and allow bile to drain when it is narrowed or blocked. Stents can move, clog, or cause infection or pancreatitis and require follow-up.

\medterm{Biliary System} The liver, gallbladder, and bile ducts that make, store, and carry bile into the small intestine to help digest fats.

\medterm{Biliary Tract} The network of ducts carrying bile from the liver and gallbladder to the small intestine. Blockage or inflammation can cause upper-abdominal pain, jaundice, infection, or pancreatitis.

\medterm{Biliary Tract Biopsy} The acquisition of tissue specimens or exfoliated cells from the extrahepatic or intrahepatic bile ducts or gallbladder for histopathologic examination, commonly performed via endoscopic retrograde cholangiopancreatography (ERCP), intraductal brush cytology, or image-guided percutaneous needle biopsy.

\synonyms
Bile Duct Biopsy; Gallbladder Biopsy; Biliary Cytology

\medterm{Biliary Tract Injuries} Damage to the gallbladder or bile ducts, often after abdominal surgery or trauma. Bile may leak or drainage may become blocked, causing abdominal pain, jaundice, infection, or later narrowing.

\medterm{Bilious} Containing or relating to bile. Bilious vomit is usually yellow or green because bile has entered the stomach contents.

\medterm{Bilious Vomiting} Vomiting green or yellow fluid containing bile. In a newborn or infant, clearly green vomit can indicate bowel blockage and requires immediate medical assessment.

\medterm{Bilirubin} A yellow substance made mainly when old red blood cells are broken down. The liver changes it and sends it into bile; too much bilirubin in the blood causes yellowing of the skin and eyes.

\medterm{Bilirubin Blood Test} A blood test measuring total and direct bilirubin. High results can occur when red blood cells break down too quickly, the liver cannot process bilirubin, or bile flow is blocked; newborn levels help assess risk of brain injury.

\medterm{Bilirubin Encephalopathy} Brain injury caused by very high levels of unconjugated bilirubin, especially in newborns. Early signs may include poor feeding, sleepiness, weak or stiff muscles, and seizures; permanent injury is called kernicterus.

\medterm{Bilirubin Metabolism} The process by which the body makes bilirubin from old red blood cells, carries it to the liver, changes it into a water-soluble form, and removes it in bile and stool.

\medterm{Bilirubin Test, Urine} A clinical urinalysis that detects conjugated (water-soluble) bilirubin in urine. A positive result indicates hepatobiliary disease, hepatocellular injury, or biliary tract obstruction; unconjugated bilirubin is bound to albumin and does not filter into urine.

\synonyms
Urine Bilirubin; Urinary Bilirubin Test

\medterm{Bilirubinuria} The presence of water-soluble bilirubin in urine. It may occur with liver disease or blocked bile flow and can make urine appear unusually dark.

\medterm{Biliverdin} A green pigment formed when haemoglobin is broken down. The body normally converts it into bilirubin, which the liver processes and removes in bile.

\medterm{Billroth I} A surgical procedure for partial gastrectomy in which the distal stomach (antrum and pylorus) is resected and the remaining proximal gastric remnant is anastomosed directly to the duodenum (gastroduodenostomy).

\synonyms
Billroth I Reconstruction; Gastroduodenostomy

\medterm{Billroth II} A surgical procedure for partial gastrectomy in which the distal stomach is resected, the duodenal stump is closed, and the remaining gastric remnant is anastomosed to a loop of the proximal jejunum (gastrojejunostomy), leaving a blind duodenal limb for biliopancreatic drainage.

\synonyms
Billroth II Reconstruction; Gastrojejunostomy

\medterm{Billroth's Operation} An older type of partial stomach-removal surgery. In Billroth I the remaining stomach is joined to the duodenum; in Billroth II it is joined to the jejunum, a later part of the small intestine.

\medterm{Bilophila} A group of bacteria that can live without oxygen in the human intestine. Most are harmless there, but Bilophila wadsworthia has been found in some appendicitis, abdominal, wound, and bloodstream infections.

\medterm{Bilophila Wadsworthia} A bacterium that tolerates bile and normally occurs in small numbers in the colon. If intestinal bacteria enter damaged or normally sterile tissue, it may contribute to appendicitis, abdominal infection, or wound infection.

\medterm{Bimanual} Performed using both hands. A bimanual pelvic examination uses one hand inside the vagina and the other on the abdomen to assess pelvic organs.

\medterm{Bimatoprost} Bimatoprost is a prostaglandin analogue used as eye drops to lower intraocular pressure and, in some formulations, promote eyelash growth.

\medterm{Bimedial} Relating to or located on the inner side of two structures. The exact structures involved must be specified because the term alone does not identify a disease or diagnosis.

\medterm{Bimellar} Made of two thin layers or plates. The term describes the structure of a membrane, tissue, or other object and has no disease meaning by itself.

\medterm{Binary Fission} A form of reproduction in which one cell divides into two genetically similar cells. It is common among bacteria and many single-celled organisms.

\medterm{Binder Excipient} An inactive ingredient added during tablet manufacture to help powder particles stick together. It has no intended treatment effect but affects the tablet's strength and breakdown.

\medterm{Binding Protein} A protein that attaches specifically to another molecule, such as a hormone, medicine, nutrient, DNA segment, or cell-surface signal. Binding can transport, store, activate, or block the other molecule.

\medterm{Binding Site} The particular area of a molecule where another molecule attaches. The interaction can change activity, transmit a signal, allow transport, or enable a laboratory test to detect the target.

\medterm{Binet-Simon Scale} An early intelligence test for children that compared performance with tasks expected at different ages. It influenced modern cognitive testing but is no longer the standard tool for clinical assessment.

\medterm{Binge Drinking} A pattern of consuming enough alcohol over a short period to produce a high blood alcohol concentration and increased risk of injury or acute toxicity.

\synonyms
Episodic Heavy Drinking; Heavy Episodic Drinking

\medterm{Binge Eating Disorder} Alternate spelling; see \textit{Binge-Eating Disorder}.

\medterm{Binge-Eating Disorder} A mental-health condition involving repeated episodes of eating large amounts of food with loss of control and distress, without regular behaviors intended to compensate. Weight alone does not establish the diagnosis.

\synonyms
BED; Compulsive Overeating Disorder

\medterm{Binocular Vision} The use of both eyes together to produce one image and judge depth. Misalignment or poor coordination between the eyes can cause double vision, eyestrain, or reduced depth perception.
\synonyms
Two-Eyed Vision; Binocular Visual Function

\medterm{Binocular Visual Loss} Reduced vision affecting both eyes. Causes may involve the eyes, optic nerves, brain, or visual pathways; sudden loss, weakness, speech trouble, or severe headache requires emergency assessment.

\medterm{Binswanger Disease} A form of small-vessel disease in the brain that damages deep white matter. It can cause gradually worsening thinking, walking difficulty, mood changes, and urinary problems.

\medterm{Binswanger's Dementia} A form of vascular cognitive impairment caused by diffuse small-vessel disease affecting the brain's deep white matter.

\synonyms
Subcortical Ischemic Vascular Disease; Subcortical Arteriosclerotic Encephalopathy

\medterm{Bio-Informatics} The use of computers and statistics to store, compare, and interpret biological information such as DNA sequences, genes, proteins, and medical data.

\medterm{Bioabsorbable Vascular Stent} A temporary endovascular scaffold constructed from biodegradable polymers (such as poly-L-lactic acid) or bioresorbable metallic alloys (such as magnesium) engineered to provide vessel patency during vascular remodeling before undergoing hydrolytic degradation and systemic resorption.

\synonyms
Bioabsorbable Vascular Scaffold; BVS; Bioresorbable Stent; BRS

\medterm{Bioaccumulation} The gradual build-up of a substance in an organism when absorption is faster than removal. Pollutants, medicines, and metals can accumulate in tissues and may cause harm when levels become high.

\medterm{Bioactive} Having an effect on living cells, tissues, or organisms, such as a bioactive drug, hormone, nutrient, or signaling molecule.

\medterm{Bioassay} A test that measures the effect or activity of a substance by observing its action on living cells, tissues, organisms, or a biological system. It can be used to estimate potency, toxicity, or biologic activity.

\medterm{Bioavailability} The amount of a medicine that reaches the bloodstream in an active form and how quickly it gets there. It depends on how the medicine is given, absorbed, and broken down; intravenous medicines have complete availability.

\medterm{Bioburden} The number and types of living microorganisms present on a surface, medical device, tissue, or medicine before sterilization. Testing it helps confirm that cleaning and sterilization methods are adequate.

\medterm{Biochanin A} A naturally occurring isoflavone found mainly in legumes and other plants. It is studied for biologic effects but is not established as a treatment for disease.

\synonyms
Biochanin a Isoflavone; 5,7-Dihydroxy-4'-Methoxyisoflavone

\medterm{Biochemical} Relating to the chemical substances and reactions that occur in living organisms, including metabolism and cellular signaling.

\medterm{Biochemical Assessment} Evaluation using laboratory measurements such as blood proteins, salts, glucose, fats, vitamins, and minerals. Results can help assess nutrition, metabolism, organ function, or disease but must be interpreted with the clinical picture.

\medterm{Biochemical Reaction} A chemical change occurring in a living cell or body, usually controlled by enzymes. It may build or break down molecules and supports energy production, growth, repair, and normal organ function.

\medterm{Biochemical Recurrence} A rise in a laboratory marker after treatment that suggests disease remains or has returned, even though symptoms or scans may not yet show it. The meaning depends on the disease and the test used.

\synonyms
Biochemical Relapse; Biochemical Failure

\medterm{Biochemistry} The study of the chemical substances and reactions occurring in living organisms. Clinical biochemistry uses laboratory measurements to assess health and disease.

\medterm{Biocompatibility} The ability of a material or device to perform its intended function without causing an unacceptable harmful reaction in body tissues.

\medterm{Biodegradable Implant} An implant designed to provide support temporarily and then gradually break down into substances the body can process or remove. Its safety depends on how quickly it breaks down and how tissues respond.

\medterm{Biodistribution} The way a medicine, tracer, or other substance spreads through the blood and different organs over time. It affects where treatment acts, how well an image is produced, and radiation or toxicity to healthy tissue.

\medterm{Bioengineering} The application of engineering principles to biology and medicine, including the design of medical devices, biomaterials, tissue-engineering systems, prostheses, and diagnostic technologies.

\medterm{Bioequivalence} A finding that two versions of a medicine enter the bloodstream at a similar rate and amount when given at the same dose. It supports the expectation of comparable effects and safety for approved uses.

\medterm{Bioethics} The study of moral questions in medicine, biology, research, and healthcare technology. It considers consent, privacy, fairness, benefit, harm, quality of life, and respect for patients and communities.

\medterm{Bioethics, Geriatric} The systematic exploration of moral and ethical considerations in the medical care of elderly patients, encompassing decisional capacity, informed consent, advanced healthcare directives, surrogate decision-making, palliative care transitions, and protection of vulnerable older adults.

\synonyms
Geriatric Bioethics; Ethics in Geriatric Medicine; Ethics of Aging

\medterm{Biofeedback} A technique that shows information about a body function, such as muscle tension, heart rate, or breathing, so a person can learn to control it. It may support treatment of selected pain, stress, or movement problems.

\medterm{Biofilm} A community of microorganisms attached to a surface and surrounded by a protective material they produce. Biofilms can form on teeth, wounds, catheters, and implants and may be difficult to remove.

\medterm{Biofilm-Associated Infections} Infections caused by microorganisms living in a protective layer on tissue, wounds, teeth, or medical devices. They can persist despite antibiotics and may require removal or replacement of an infected device.

\medterm{Biofluid} A fluid from a living body, such as blood, urine, saliva, spinal fluid, or joint fluid. Biofluids can be tested for cells, chemicals, microorganisms, medicines, or signs of disease.

\medterm{Biologic Therapy} Treatment using a medicine made from living cells or biological substances to target a specific immune or disease pathway. It is used for selected inflammatory, autoimmune, cancer, and other conditions.

\medterm{Biologic Therapy In Dermatology} Treatment with targeted biological agents, usually monoclonal antibodies or receptor
proteins, that inhibit selected immune pathways. Biologics are used for diseases such as
psoriasis, atopic dermatitis, hidradenitis suppurativa, and urticaria; infection risk,
vaccination status, screening, and monitoring must be considered.

\medterm{Biologic Variation} Normal change in a body measurement within the same person over time or between different people. It must be considered when deciding whether a laboratory change represents real disease or treatment-related change.

\medterm{Biological} Relating to living organisms, their parts, functions, processes, or products. In medicine, a biological product is made from or resembles a substance from a living source.

\medterm{Biological Availability} The amount and speed at which a medicine or other substance becomes available to produce an effect in the body. It depends on the route, absorption, breakdown, and removal of the substance.

\medterm{Biological Clock} The body's internal timing system that helps coordinate daily sleep, alertness, hormone release, body temperature, and other repeating functions in response to light and other environmental signals.

\medterm{Biological Half-Life} The time needed for the body to remove half of a medicine or other substance through metabolism and excretion. It differs from the physical half-life of a radioactive substance and helps guide dosing or safety decisions.

\medterm{Biological Model} A cell, tissue, organism, computational system, or other representation used to study biological mechanisms, disease, or treatment responses.

\medterm{Biological Process} A coordinated series of molecular, cellular, or physiological events in a living organism, such as metabolism, cell division, inflammation, or tissue repair.

\medterm{Biological Relative} A person who shares genetic ancestry with someone. Family medical history uses biological relatives to assess inherited risks, while relatives by marriage or legal relationship do not share that genetic information.

\medterm{Biological Therapy} Treatment using biological substances, cells, vaccines, antibodies, or similar products to change immune activity or target a specific disease pathway. It is used in selected cancers, immune disorders, and other conditions.

\medterm{Biological Transport} The movement of substances such as water, salts, nutrients, medicines, or cells within the body or across cell membranes. It may occur by diffusion, channels, pumps, filtration, or transport vesicles.

\medterm{Biological Variability} Natural differences in a body measurement between people or changes in the same person over time, separate from laboratory error. It can affect how test results are interpreted.

\synonyms
Biologic Variation; Physiologic Variability

\medterm{Biological Warfare} The deliberate use of microorganisms, toxins, or other biological agents to cause illness, death, or disruption in people, animals, or plants.

\medterm{Bioluminescence} Light produced by a living organism through a chemical reaction. It is used by some organisms for communication or defense and is used in laboratories to detect biological activity.

\medterm{Biomarker} A measurable feature that gives information about a normal body process, disease, likely outcome, or response to treatment. It may be measured in blood, tissue, genes, images, or other samples.

\medterm{Biomarker Testing} Laboratory testing of genes, proteins, metabolites, or other measurable features to help identify disease, estimate likely outcome, choose treatment, or monitor response.

\medterm{Biomaterial} A natural or manufactured material designed to contact or interact with the body. Examples include implant surfaces, artificial joints, wound dressings, prostheses, and materials that release medicines.

\medterm{Biomechanical Compliance} How easily a body structure changes size or shape when pressure or force is applied. Lung and blood-vessel compliance affect breathing and circulation, while joint compliance affects movement and stability.

\medterm{Biomechanics} The study of how forces and movement affect the body. It examines walking, posture, joints, muscles, and loads to understand injury, improve rehabilitation, or design supportive devices.

\medterm{Biomedical Cement} A medical material used to fill gaps, stabilize broken bone, anchor an implant, or repair tissue. Its use and strength depend on the procedure and the tissue being treated.

\medterm{Biomedical Device} An instrument, implant, software system, or other product intended to diagnose, monitor, prevent, or treat disease or injury. Its safety and usefulness depend on proper design, testing, maintenance, and correct use.

\medterm{Biomedical Engineering} The use of engineering, materials, computing, and design to create or improve medical devices, diagnostic tests, implants, prostheses, imaging systems, and other healthcare technologies.

\medterm{Biomedical Tape} Adhesive tape designed for medical use, such as securing dressings, tubes, sensors, or small devices. Different types vary in strength, flexibility, water resistance, and risk of skin irritation.

\medterm{Biomedicine} Medicine grounded in the biological sciences, including the study of disease mechanisms, diagnosis, prevention, and treatment through laboratory and clinical research.

\medterm{Biometry} Measurement of body structures or functions. In eye care, optical biometry measures the eye to help calculate the strength of an artificial lens before cataract surgery.

\medterm{Biomicroscopy} Examination of living tissue with a microscope. In eye care, a slit lamp provides magnified views of the cornea, front chamber, lens, and other eye structures.

\medterm{Biomolecule} A molecule made by or found in living organisms, such as a protein, fat, sugar, DNA, or RNA. Biomolecules provide structure, energy, information, and control of cell functions.

\medterm{Biomonitoring} Repeatedly measuring substances or body changes over time using blood, urine, breath, hair, or other samples. It can assess exposure to poisons or medicines, treatment levels, hormones, or disease activity.

\medterm{Biomphalaria} A group of freshwater snails that can carry the parasite Schistosoma mansoni. The parasite develops in the snails and can later infect people who contact contaminated freshwater.

\medterm{Bionics} The use of electronic or mechanical devices to replace or support body functions, such as hearing, movement, or sensation. Examples include cochlear implants, powered prostheses, and other assistive systems.

\medterm{Biophotonics} The use of light to study, diagnose, monitor, or treat living tissues and biological processes. Examples include optical imaging, fluorescence testing, laser surgery, and photodynamic therapy.

\medterm{Biophysical Profile} A pregnancy assessment combining fetal-heart-rate monitoring with ultrasound. It checks fetal movement, breathing movements, muscle tone, and amniotic-fluid volume to help assess fetal well-being.

\medterm{Biopsy} Removal of a small sample of cells or tissue so it can be examined under a microscope to diagnose disease, including cancer. Types include fine-needle aspiration (cells or fluid), core-needle (a tissue cylinder), incisional (part of an area), excisional (the whole area), punch (a circular piece of skin), and endoscopic biopsy (through a viewing tube). Possible complications include pain, bleeding, and infection; risks vary with the body site and procedure.

\synonyms
Biopsies

\medterm{Biopsy Lung} Removal of a small lung-tissue sample using a needle, bronchoscope, or operation. Examination can investigate infection, inflammation, scarring, or cancer; the method depends on the lesion and procedural risks.

\medterm{Biopsy Needle} A needle used to remove cells or tissue for laboratory examination. A fine needle usually collects loose cells, while a core needle removes a small cylinder that preserves the tissue's arrangement.

\medterm{Biopsy Site} The body location from which cells or tissue are removed for testing. It is chosen to sample the abnormal area while limiting injury to nearby organs, nerves, and blood vessels.

\medterm{Biopsy Specimen} Cells or tissue removed from the body for laboratory examination. Microscopy, special stains, immune tests, or genetic tests may be used to identify cancer, infection, inflammation, or another condition.

\medterm{Biosafety} Practices, equipment, and containment methods used to prevent accidental exposure to or release of infectious organisms, biological materials, or toxins.

\medterm{Biosafety Level} A classification describing the laboratory facilities, equipment, work practices, and protective measures needed to handle biological agents safely according to their risk.

\medterm{Biosecurity} Measures that prevent biological agents or toxins from being accidentally or deliberately released, stolen, or misused. It includes secure facilities, controlled access, staff training, risk assessment, and responsible handling of biological materials.

\medterm{Biosensor} A device that uses a biological material to recognize a substance and a detector to measure it. Medical examples include sensors for glucose, hormones, infectious organisms, or heart-related markers.

\medterm{Biosimilar} A biological medicine highly similar to an already approved biological medicine, with no meaningful differences in safety or effectiveness for approved uses. Because these medicines are complex, similarity is shown through comparative testing.

\medterm{Biospecimen} A sample of human tissue, cells, blood, urine, or another body fluid collected for testing or research. Correct collection, labelling, storage, and handling help keep the sample reliable.

\medterm{Biostatistics} The use of statistics to plan medical studies, analyse health data, measure uncertainty, and interpret evidence about disease, diagnosis, prevention, and treatment.

\medterm{Biosynthesis} The process by which living cells use enzymes and energy to build complex molecules from simpler substances, such as making proteins from amino acids.

\medterm{Biot Breathing} An irregular breathing pattern with groups of normal-sized breaths interrupted by unpredictable pauses in breathing. It may occur with serious damage affecting the lower brainstem and requires urgent assessment.

\medterm{Biot's Breathing} An abnormal pattern of irregular breaths interrupted by unpredictable pauses in breathing. It suggests disruption of the brain's breathing control and may occur with severe neurological illness.

\medterm{Biotechnology} The use of living cells, organisms, or biological molecules to develop products and processes such as vaccines, biological medicines, genetic tests, and gene therapies.

\medterm{Biotelemetry} Measuring and sending body information from a person or animal to a distant receiver. It can be used for heart-rhythm, temperature, breathing, glucose, or intensive-care monitoring.

\medterm{Bioterrorism} The deliberate release of infectious organisms, toxins, or other biological agents to cause illness, death, fear, or disruption. Response involves rapid detection, public-health coordination, protection, testing, isolation when needed, and treatment or prevention.

\medterm{Bioterrorism Agents} Infectious organisms or toxins that could be deliberately released to cause widespread illness or panic. Public-health authorities prepare for agents causing anthrax, smallpox, botulism, plague, tularemia, or severe viral haemorrhagic fever.

\medterm{Biotherapy} Treatment using living organisms, cells, genes, or biological substances to prevent or treat disease. Immunotherapy, some cell therapies, and selected biological medicines are forms of biotherapy.

\medterm{Biotin} A water-soluble B vitamin needed for enzymes involved in processing fats, sugars, and amino acids. Deficiency is uncommon but can cause a skin rash, hair loss, weakness, or neurological symptoms.

\medterm{Biotin Vitamin B7} Alternate name; see \textit{Biotin}.
\medterm{Biotinidase Deficiency} A rare inherited disorder in which the body cannot recycle biotin efficiently. Without treatment it can cause seizures, developmental problems, hearing or vision loss, skin rash, and weakness; lifelong biotin usually prevents symptoms.

\medterm{Biotinylation} The attachment of biotin to another molecule, usually a protein, so it can be found or separated in a laboratory test. The biotin acts as a tag that binds to special substances, making the target easier to detect or collect. Biotinylation can occur naturally or be done artificially.

\medterm{Biotransformation} The natural process by which the body changes drugs and other foreign substances, mainly through enzymes in the liver but also in other tissues. It often makes them less active and easier to remove. Phase I reactions modify the substance; phase II reactions may attach another molecule to help eliminate it. These steps are not always sequential, and the process can produce an active or harmful substance instead.

\medterm{Biotrauma} Inflammation caused by mechanical injury to the lungs during breathing-machine treatment. Inflammatory signals can spread through the body and contribute to organ injury; careful, lung-protective settings reduce the risk.

\medterm{BiPAP} Bilevel Positive Airway Pressure (BiPAP), a form of noninvasive breathing support delivered through a mask. It provides higher pressure during inhalation and lower pressure during exhalation. It may be used for conditions such as hypercapnic COPD exacerbation, sleep-related breathing disorders, and cardiogenic pulmonary edema.

\synonyms
Bilevel Positive Airway Pressure

\medterm{BiPAP Machine} A device that delivers two levels of breathing pressure through a mask—higher during inhalation and lower during exhalation. Settings are prescribed and adjusted according to the person's breathing problem and response.

\medterm{Biparental} Involving both biological parents. In genetics, biparental inheritance means that a child receives relevant genetic material from both the mother and the father.

\medterm{Biparietal Diameter} The width of a baby's head measured by ultrasound from one side to the other. It helps estimate pregnancy age and monitor growth, but is interpreted with other measurements because head shape can affect the result.

\medterm{Bipartite Ossification} Formation or persistence of two separate ossification centers in a bone. If they do not fuse, a bipartite bone, such as a bipartite patella, may result. This is usually an incidental developmental variant but can become painful after injury or repetitive stress.

\medterm{Biphenyl Compounds} Chemicals made of two linked benzene rings. Polychlorinated biphenyls (PCBs) are long-lasting industrial pollutants; exposure can harm the liver, immune system, and hormone system and may increase cancer risk.

\medterm{Bipolar} Alternate name; see \textit{Bipolar Disorder}.

\medterm{Bipolar Affective Disorder} Alternate name; see \textit{Bipolar Disorder}.

\medterm{Bipolar Disorder} A mental health condition involving episodes of mania, hypomania, and/or major depression. Bipolar I is defined by at least one manic episode; depressive episodes are common but not required. Bipolar II includes at least one hypomanic episode and at least one major depressive episode, without a history of mania. Episodes can affect mood, sleep, energy, judgment, and daily functioning.

\synonyms
Manic-Depressive Illness

\medterm{Bipolar I Disorder} A mood disorder defined by at least one manic episode lasting at least seven days or
requiring hospitalization. A manic episode is a distinct period of abnormally elevated,
expansive, or irritable mood and increased goal-directed activity or energy.

\medterm{Bipolar II Disorder} A mood disorder characterized by at least one hypomanic episode and at least one major depressive episode, with no history of a full manic episode. The combination of hypomania and depression can cause substantial distress or impairment even though hypomania is less severe than mania.

\medterm{Bipolar Neuron} A bipolar neuron has one dendrite and one axon extending from opposite sides of its cell body, as found in the retina and sensory pathways.

\medterm{Bipolaris} A genus of fungi found mainly in soil and plants. Some species can rarely cause eye, sinus, skin, or invasive infections, particularly in people with reduced immunity. Identification requires laboratory culture or molecular testing because treatment depends on the species and site of infection.

\medterm{Bird Fancier's Lung} A form of hypersensitivity pneumonitis caused by breathing proteins from bird feathers or droppings. It can cause cough and shortness of breath; repeated exposure may lead to permanent lung scarring.

\medterm{Bird Flu} An infectious zoonotic illness caused by avian influenza type A viruses (notably H5N1, H5N6, H7N9) that naturally circulate among wild waterfowl and domestic poultry. Transmission to humans occurs via direct contact with infected birds or contaminated environments, producing acute respiratory distress syndrome (ARDS), multi-organ dysfunction, and high case-fatality rates.

\synonyms
Avian Influenza; Avian Flu; Bird Flu; H5N1 Infection

\medterm{Birt-Hogg-Dubé Syndrome} An inherited condition caused by a change in the FLCN gene. It can cause small, usually harmless skin growths, cysts in the lungs with repeated collapsed lung, and an increased risk of kidney tumors. It is autosomal dominant, meaning one altered copy of the gene can cause the condition.

\medterm{Birth} The process by which a baby leaves the uterus, including labour, delivery of the baby, and delivery of the placenta.

\medterm{Birth Asphyxia} Too little oxygen or inadequate breathing before, during, or shortly after birth. It can injure the brain and other organs and requires immediate newborn assessment, breathing support, and treatment.

\medterm{Birth Control} Methods used to prevent pregnancy, including condoms, pills, implants, injections, intrauterine devices, fertility-awareness methods, and permanent procedures. Choice depends on effectiveness, health, side effects, future plans, access, and personal preference.

\medterm{Birth Control And Family Planning} Planning whether and when to have children using contraception, fertility care, pregnancy advice, and reproductive health services. Care respects personal choices and considers health conditions, safety, goals, and informed consent.

\medterm{Birth Control Methods} Birth-control methods include condoms, pills, implants, injections, intrauterine devices, patches, rings, fertility-awareness approaches, and permanent contraception. Effectiveness, sexually transmitted infection protection, side effects, reversibility, and medical eligibility vary by method.

\medterm{Birth Control Pill Overdose} Illness after taking too many contraceptive tablets. Most exposures cause little or no serious harm, but nausea, vomiting, drowsiness, or other symptoms may occur, and the product and amount taken help guide management.

\medterm{Birth Control Pills} Tablets taken on a schedule to prevent pregnancy, usually by stopping ovulation, thickening cervical mucus, or changing the uterine lining. They do not protect against sexually transmitted infections and work best when taken correctly.

\medterm{Birth Control, Extended-Release} Alternate name; see \textit{Long-Acting Reversible Contraception}.

\medterm{Birth Defect} A structural, functional, or metabolic abnormality present at birth resulting from chromosomal aberrations, monogenic mutations, prenatal teratogenic exposures (e.g., alcohol, retinoic acid, infections), or complex multifactorial interactions during embryogenesis and organogenesis.

\synonyms
Congenital Anomaly; Congenital Malformation; Birth Defects; Congenital Disorder

\medterm{Birth Defects} Structural or functional problems present at birth that may affect development or health. Causes include genetic or chromosomal changes, infections, medicines, nutritional deficiencies, and environmental exposures; sometimes no cause is found.

\medterm{Birth Injury} Physical harm to a newborn during labour or delivery. It may affect the skin, bones, nerves, muscles, or brain and can result from difficult delivery, the baby's position, or medical complications.

\medterm{Birth Order} The sequence in which siblings are born within a family. While birth order has been studied for psychological effects, its direct medical significance is limited to obstetric considerations such as risks associated with grand multiparity.

\medterm{Birth Pains} Birth pains are painful uterine contractions and related pressure during labor as the cervix dilates and the baby moves through the birth canal.

\medterm{Birth Rate} The number of live births in a population during a specified period, usually expressed per 1,000 people per year.

\medterm{Birth Spacing} The practice of planning the time between giving birth and becoming pregnant again. Health guidelines generally recommend waiting at least 18 to 24 months between pregnancies to allow the mother's body to recover and restore essential nutrients, reducing the risks of premature birth, low birth weight, and delivery complications.
\synonyms Interpregnancy interval; Child spacing; Pregnancy spacing.

\medterm{Birth Weight} A newborn's weight measured soon after birth. Weight below 2,500 grams is considered low and is linked with increased health risks; weight below 1,500 grams is very low and usually requires specialized newborn care.

\medterm{Birthmark} A localized mark on the skin or in blood vessels that is present at birth or appears soon afterward. Examples include pigmented moles, haemangiomas, and port-wine stains.

\medterm{Birthmarks} Marks on the skin present at birth or appearing early in life. Most are harmless, but rapid growth, colour or shape change, bleeding, ulceration, pain, or a mark affecting vision or breathing needs medical assessment.

\medterm{Birthmarks And Other Skin Pigmentation Problems} Birthmarks or changes in skin colour may result from blood vessels, pigment cells, inflammation, medicines, hormones, or inherited conditions. A rapidly changing, bleeding, ulcerated, painful, or distinctly different mark requires examination.

\medterm{Bis In Die} Latin for “twice daily,” traditionally abbreviated b.i.d. Clear prescriptions state when doses are to be taken rather than relying only on this abbreviation.

\medterm{Bisacodyl} A stimulant laxative that increases bowel movement and helps empty the bowel. It is used short term for constipation or before some procedures; cramps, diarrhoea, and dehydration can occur.

\medterm{Bisexual} Bisexual describes sexual or romantic attraction to more than one gender; it is an identity or orientation, not a disease or disorder.

\medterm{Bisexuality} Bisexuality is a sexual orientation involving attraction to more than one gender; it is not a disease or disorder.

\medterm{Bishop Score} A score used before inducing labour to describe how ready the cervix is for birth. It considers opening, thinning, softness, position, and the baby's head position; a higher score generally predicts an easier induction.

\medterm{Bismuth} A metal used in some medicines and industrial products. Certain bismuth compounds can soothe diarrhoea or stomach symptoms, but excessive or prolonged exposure can cause nervous-system or kidney toxicity.

\medterm{Bisoprolol} A beta-1-selective blocker used to treat hypertension, angina, and selected patients with heart failure. It lowers heart rate and blood pressure and may cause bradycardia, hypotension, or fatigue.

\medterm{Bispecific Antibody Therapy} Treatment using an engineered antibody that binds two different targets, such as a cancer-cell marker and an immune-cell marker, to bring cells together or block more than one disease pathway. It is used for selected cancers and can cause immune reactions or other treatment-specific effects.

\medterm{Bisphenol A} An industrial chemical used in some plastics and resins. It can act like a weak hormone, and possible health effects of exposure are studied; reducing unnecessary contact is especially reasonable for infants and children.

\medterm{Bisphenol A} An endocrine-disrupting industrial chemical used predominantly in the synthesis of polycarbonate plastics and epoxy resins lining food and beverage containers. It exhibits weak estrogenic and antiandrogenic activity and has been implicated in metabolic and reproductive disturbances.

\synonyms
BPA; 4,4'-(Propane-2,2-diyl)diphenol

\medterm{Bisphosphonate} A medicine that slows the removal of old bone. It is used mainly for osteoporosis, Paget disease, and some cancer-related bone problems. Rare but important risks include jaw-bone damage and unusual thigh-bone fractures.

\medterm{Bisphosphonates} Medicines that slow the breakdown of bone. They are used mainly to prevent fractures in osteoporosis and to treat Paget disease and some bone problems caused by cancer.

\synonyms
Bisphosphonate

\medterm{Bite, Chigger} Alternate index label; see \textit{Chiggers}.

\medterm{Bite, Flea} Alternate index label; see \textit{Flea Bites in Humans}.

\medterm{Bite, Snake} Alternate index label; see \textit{Snake Bite}.

\medterm{Bitemporal Hemianopia} Loss of side vision on the outer side of both eyes. It usually occurs when a growth or other problem presses on the place where the two optic nerves cross, often near the pituitary gland.

\medterm{Bitewing} A dental X-ray showing the crowns of upper and lower teeth and the bone around them in one image. It is especially useful for finding decay between teeth and early loss of supporting bone.

\medterm{Bitewing Radiograph} An X-ray taken inside the mouth that shows the crowns of back teeth and the nearby gum-supporting bone. It helps detect decay between teeth and bone loss from gum disease.

\medterm{Bitot's Spots} Distinctive, triangular, foamy, silvery-white or pearly-gray patches on the white part of the eye (bulbar conjunctiva), most commonly found on the outer (temporal) side near the cornea. They are a classic sign of vitamin A deficiency (Stage X1B of xerophthalmia) caused by the buildup of dried keratin and gas-forming bacteria on the dried eye surface. Bitot's spots commonly occur alongside night blindness and ocular dryness. Prompt treatment with high-dose vitamin A supplements usually resolves the spots within weeks and prevents progression to permanent corneal damage and blindness.
\synonyms{Bitot Spots; Conjunctival Foamy Patches; Keratinized Conjunctival Plaques}

\medterm{Bituminosis} A work-related breathing disorder caused by repeated exposure to bitumen or asphalt fumes and dust. It may irritate the airways and contribute to chronic cough, breathlessness, or other lung disease.

\medterm{Biuret Test} A laboratory test that detects protein by turning violet when protein reacts with a copper-containing solution. It is mainly used in teaching, research, and some chemical analyses rather than routine patient diagnosis.

\medterm{Bivalve} Bivalve means having two shells or valves; bivalve mollusks may be relevant to food allergy, poisoning, or parasitic exposure.

\medterm{Bizarre Parosteal Osteochondromatous Proliferation} A rare benign but locally recurrent bone-surface lesion, also called Nora Lesion, composed of cartilage, bone, and spindle cells; it most often affects the hands or feet.

\medterm{BK Virus} A common virus that usually remains inactive after childhood infection. It can reactivate when immunity is suppressed, especially after organ transplantation, and damage the transplanted kidney or urinary tract.

\medterm{BK Virus Nephropathy} Kidney inflammation and damage caused by reactivation of BK virus, usually in a kidney-transplant recipient. It can reduce transplant function and is managed mainly by carefully reducing immune-suppressing treatment and monitoring viral levels.

\medterm{Black Death} The major plague pandemic of the fourteenth century, caused mainly by the bacterium Yersinia pestis. It caused widespread severe illness and death and is distinct from the individual forms of plague seen today.

\medterm{Black Eye} Bruising and swelling around the eye, usually caused by blunt trauma. Vision may remain normal, but eye pain, double
vision, reduced vision, an irregular pupil, severe headache, vomiting, or loss of consciousness can indicate an eye
or orbital injury and requires urgent assessment.

\synonyms
Periorbital Hematoma; Shiner

\medterm{Black Fly} A small biting fly that can cause painful or itchy skin reactions. In some tropical regions, black flies spread the parasite that causes river blindness.

\medterm{Black Hair, Hairy Tongue} A harmless condition in which the tiny surface projections of the tongue become long and dark, trapping debris or pigment. Smoking, antibiotics, dry mouth, and poor oral care can contribute.

\medterm{Black Hairy Tongue} A usually harmless dark, furry appearance of the tongue caused by overgrown surface projections trapping debris or pigment. It may improve with tongue cleaning, good oral hygiene, and addressing smoking or dry mouth.

\medterm{Black Lung} Lung disease caused by long-term inhalation of coal dust, usually in miners. It can cause cough and breathlessness and may progress from small areas of dust-related damage to extensive permanent scarring.

\medterm{Black Lung Disease} A dust-related lung disease caused by prolonged coal-dust exposure. It may cause no symptoms at first but can lead to cough, breathlessness, reduced lung function, and severe permanent scarring.

\medterm{Black Nightshade Poisoning} Poisoning after eating black nightshade or its berries, which may contain toxic plant chemicals. It can cause nausea, vomiting, abdominal pain, confusion, enlarged pupils, seizures, or abnormal heart rhythms.

\medterm{Black Or Tarry Stools} Black or tarry stools may indicate upper gastrointestinal bleeding, but can also result from iron or bismuth. New or persistent findings require medical assessment.

\medterm{Black Stain} Dark discoloration on the tooth surface, often near the gumline, caused by external deposits or chromogenic bacteria. Dental assessment
determines the cause and appropriate removal.

\medterm{Black Triangle} A triangular gap between neighboring teeth near the gumline, usually caused by loss or shrinkage of the gum papilla after gum disease, injury, or dental treatment. It may trap food and can be treated cosmetically or orthodontically.

\medterm{Black Widow Spider} The black widow spider is a venomous spider whose bite can cause severe muscle pain, abdominal cramping, sweating, hypertension, and autonomic symptoms.

\medterm{Blackflesh} A nontechnical description of tissue that looks dark or black. Causes include bruising, severe loss of blood supply, dead tissue, infection, or staining. New, painful, cold, or spreading discoloration needs urgent medical assessment.

\medterm{Blackhead} An open blocked hair follicle filled with oil and dead skin cells. The material looks black because it reacts with air, not because the skin is dirty.

\medterm{Blackheads} Open comedones formed when a hair follicle becomes blocked with sebum and keratin and
the surface material oxidizes. They are a common non-inflammatory lesion of acne,
especially on the face, chest, and back.

\medterm{Blackout} A period for which a person cannot remember events, often after heavy alcohol use. Memory loss can also occur with drugs, seizures, head injury, fainting, or other causes of altered consciousness.

\synonyms
Amnestic Episode; Alcohol-Related Amnesia When Caused by Alcohol

\medterm{Blackwater Fever} A severe complication of malaria in which red blood cells break down inside blood vessels. It can cause dark urine, severe anaemia, jaundice, kidney failure, and shock and requires emergency treatment.

\medterm{Bladder} A muscular hollow organ that stores urine from the kidneys and contracts to release it through the urethra during urination.

\medterm{Bladder Augmentation} Surgery that enlarges the bladder, usually by adding a section of bowel or another tissue. It can increase urine storage when the bladder is too small or has high pressure.

\medterm{Bladder Biopsy} Removal of a small piece of bladder lining, usually during cystoscopy, for laboratory examination. It can investigate a visible growth, unexplained blood in the urine, inflammation, or suspected bladder cancer.

\medterm{Bladder Calculi} Alternate name; see \textit{Bladder Stones}.

\medterm{Bladder Cancer} A malignant tumor arising in the urinary bladder, most often from the urothelial lining. Blood in the urine is a common presenting
feature, although symptoms vary.

\synonyms
Urinary Bladder Neoplasm; Urothelial Carcinoma

\medterm{Bladder Carcinoma} Cancer that begins in the lining of the urinary bladder, most often in cells called urothelial cells. Blood in the urine is common; smoking and some workplace chemical exposures increase risk.

\medterm{Bladder Carcinoma} Alternate name; see \textit{Bladder Cancer}.

\medterm{Bladder Control Problems In Children} Difficulty controlling urine in a child who has passed the usual stage of toilet training. It may involve daytime leakage, bedwetting, urgency, frequent urination, or difficulty emptying and may be related to development, constipation, infection, nerve problems, or urinary-tract differences.

\medterm{Bladder Control, Loss of} Alternate index label; see \textit{Urinary Incontinence}.

\medterm{Bladder Diary} A record of drinks, urination times, urine amounts, urgency, leakage, and nighttime voiding over several days. It helps describe bladder habits and patterns of urinary symptoms.

\medterm{Bladder Diseases} Disorders of the bladder, including infection, stones, tumors, neurogenic dysfunction, inflammation, obstruction, and impaired emptying.
\medterm{Bladder Disorders} A broad spectrum of pathological conditions affecting the urinary bladder, encompassing infectious cystitis, interstitial cystitis (bladder pain syndrome), urothelial carcinoma, neurogenic bladder dysfunction, calculi, and bladder outlet obstruction.

\synonyms
Bladder Diseases; Urinary Bladder Diseases; Diseases of the Bladder

\medterm{Bladder Diverticula} Pockets that push outward from the bladder wall. They may be present from birth or develop when urine repeatedly faces resistance leaving the bladder, and can trap urine, leading to infection or stones.

\medterm{Bladder Diverticulum} A pouch extending from the bladder wall. It may cause no symptoms or trap urine and lead to infection, stones, or incomplete emptying; treatment addresses the cause and symptoms.

\medterm{Bladder Exstrophy} Bladder Exstrophy is a rare birth defect in which the bladder is open and exposed outside the body through a defect in the lower abdominal wall. It is part of the exstrophy–epispadias complex and may also affect the urethra, pelvic bones, pelvic floor, abdominal muscles, and genitals. It can cause difficulty storing urine, urinary leakage, urine-control problems, and risks to the kidneys.

\medterm{Bladder Exstrophy Repair} Reconstructive surgery for bladder exstrophy that places the bladder inside the body and repairs the lower abdomen and genitals. Treatment is often staged and aims to protect the kidneys and improve urine control.

\medterm{Bladder Fistula} An abnormal passage between the bladder and another organ or the skin, often the vagina, bowel, or uterus. It may cause constant urine leakage, repeated infection, gas or stool in urine, or unusual vaginal wetness.

\medterm{Bladder Health And Disorders} Healthy bladder function stores urine and releases it at an appropriate time. Common problems include infection, urgency, leakage, inability to empty, stones, inflammation, and nerve-related bladder dysfunction.

\medterm{Bladder Infection} An infection of the bladder, usually caused by bacteria. It may cause burning urination, frequent or urgent urination, lower-abdominal discomfort, cloudy or bloody urine, and sometimes fever.

\medterm{Bladder Infection} Alternate name; see \textit{Cystitis}.

\medterm{Bladder Inflammation} Alternate name; see \textit{Interstitial Cystitis}.

\medterm{Bladder Neck} The narrow lower part of the bladder where it joins the urethra. Its muscles help control the release of urine and can contribute to blockage if they scar or fail to relax.

\medterm{Bladder Neck Contracture} Scar-related narrowing of the bladder outlet, often after prostate surgery or other treatment. It can restrict urine flow and cause a weak stream, difficulty emptying, or urinary retention.

\medterm{Bladder Neck Obstruction} Narrowing or blockage where the bladder joins the urethra, making it difficult for urine to leave the bladder. It may cause a weak stream, straining, incomplete emptying, urinary retention, or recurrent infection.

\medterm{Bladder Neoplasm} An abnormal growth arising in the bladder. It may be noncancerous or cancerous; diagnosis requires examination, imaging, cystoscopy, and often a tissue sample to determine its type and extent.

\medterm{Bladder Outlet Obstruction} A blockage where urine leaves the bladder, caused by an enlarged prostate, scarred urethra, stone, tumour, or pelvic-organ prolapse. It may cause a weak stream, difficulty starting, retention, infection, or kidney damage.

\medterm{Bladder Pain} Pain or pressure felt in the lower abdomen or pelvis that may relate to infection, stones, inflammation, bladder pain syndrome, injury, or cancer. Blood in urine, fever, severe pain, or inability to urinate needs prompt assessment.

\medterm{Bladder Pain Syndrome} Long-lasting bladder or pelvic pain, pressure, or discomfort with urinary urgency or frequency when infection and other causes have been excluded. Symptoms often flare and settle and require individualized treatment.

\medterm{Bladder Perforation} A tear or hole in the bladder wall caused by trauma, pelvic surgery, instrumentation,
severe inflammation, or obstruction. Urine may leak into surrounding tissues, causing abdominal pain,
distention, blood in the urine, or infection. It requires urgent medical assessment and may need catheter
drainage or surgical repair.

\medterm{Bladder Problems} Conditions that interfere with storing or passing urine, such as leakage, urgency, repeated infection, stones, or inability to empty. Causes include infection, weak pelvic muscles, nerve disease, prostate enlargement, or blockage.

\medterm{Bladder Prolapse} Alternate name; see \textit{Anterior Vaginal Prolapse}.

\medterm{Bladder Retraining} A planned urination schedule that gradually increases time between bathroom visits to improve bladder control and reduce urgency.

\medterm{Bladder Spasm} An involuntary contraction of the bladder muscle that can cause sudden urgency, cramping,
urine leakage, or pain. It may occur with infection, catheter use, bladder irritation, neurological
conditions, or after surgery. Treatment addresses the cause and may include bladder training or medicines.

\medterm{Bladder Stones} Hard mineral deposits in the bladder that may cause pain, blood in urine, frequent urination, interrupted flow, or difficulty emptying. Some cause no symptoms and are found during tests for another problem.

\synonyms
Bladder Calculi

\medterm{Bladder Trauma} Injury to the bladder caused by a blow, pelvic fracture, penetrating wound, surgery, or a medical instrument. It may cause blood in the urine, lower-abdominal pain, difficulty urinating, or leakage of urine into nearby tissues.

\medterm{Bladder, Urinary} Alternate index label; see \textit{Urinary Bladder}.

\medterm{Blalock-Taussig Operation} Surgery that creates a connection between a major artery and a lung artery to increase blood flow to the lungs in selected blue-baby heart defects. Modern versions usually use a small artificial tube.

\medterm{Bland Diet} A temporary eating pattern using soft, mildly seasoned foods and avoiding items that worsen a person's symptoms. It may help some stomach or bowel problems but does not treat their underlying cause.

\medterm{Blaschko Lines} Invisible patterns formed by the movement of skin cells before birth. Some birthmarks and skin diseases follow these lines as stripes or swirls; they usually do not follow nerves or blood vessels.

\medterm{BLAST} An immature cell that normally develops into a blood cell. A high number of blasts in blood or bone marrow may indicate acute leukaemia or another serious marrow disorder.

\medterm{Blast Cell} An immature precursor cell, especially one in the bone marrow that will normally develop into a blood cell. Too many blasts can indicate acute leukaemia or another marrow disorder.

\medterm{Blast Crisis} An advanced, rapidly worsening phase of some blood cancers in which immature blast cells fill the bone marrow or blood and behave like acute leukaemia. It requires urgent specialist treatment.

\medterm{Blast Injury} Harm caused by an explosion. The pressure wave can injure lungs, ears, or organs; flying objects can cause wounds; being thrown can cause fractures; burns, smoke, toxins, and crushing can add further injuries.

\medterm{Blast Phase} An advanced phase of some blood cancers in which many immature blast cells appear in the blood or marrow. The disease becomes more aggressive and may resemble acute leukaemia.

\medterm{Blastic Plasmacytoid Dendritic-Cell Neoplasm} A rare, fast-growing blood cancer arising from immature immune cells. It often affects the skin, bone marrow, lymph nodes, or nervous system and requires urgent specialist treatment.

\medterm{Blastocyst} An early embryo formed about five to six days after fertilization. It is a hollow group of cells with an inner group that can become the fetus and an outer layer that helps form the placenta.

\medterm{Blastocyst Transfer} An assisted-reproduction procedure in which an embryo grown in the laboratory for about five or six days is placed inside the uterus. The decision to transfer one or more embryos considers success rates, patient factors, and the risk of multiple pregnancy.

\medterm{Blastocystis} A common, single-celled microscopic parasite that lives in the human digestive tract, spread by swallowing food or water contaminated with the organism. While many individuals carry it harmlessly without symptoms, it can cause intestinal illness (blastocystosis) leading to watery diarrhea, abdominal cramps, bloating, gas, nausea, and occasionally skin hives. Diagnosis is confirmed by detecting the parasite during a stool examination under a microscope or with stool DNA tests. When symptomatic and persistent, it is treated with antiparasitic medications such as metronidazole.

\synonyms
Blastocystis Hominis; Blastocystis spp.; Blastocystosis; Intestinal Blastocystis

\medterm{Blastocystis Hominis} Historical name; see \textit{Blastocystis}.

\medterm{Blastocystosis} Alternate name; see \textit{Blastocystis}.

\medterm{Blastoma} A cancer that develops from immature or embryonic cells. The specific name identifies the organ and tumour type, such as a kidney, nerve, liver, or eye blastoma.

\medterm{Blastomyces Dermatitidis} A fungus found in soil and damp organic material that can cause blastomycosis when its spores are inhaled. Infection usually affects the lungs but can spread to skin, bones, joints, or other organs.

\medterm{Blastomycosis} A fungal infection acquired by inhaling Blastomyces spores, usually from soil. It commonly causes pneumonia and may spread to the skin, bones, joints, or other organs; symptoms can develop gradually.

\medterm{Blasts} Immature cells in blood-forming tissue that normally develop into mature blood cells. An excess in blood or bone marrow may indicate acute leukaemia or another serious marrow disorder.

\medterm{Blautia Producta} A bacterium that normally lives in the human intestine and grows without oxygen. It is usually harmless but has occasionally been found in infections, especially in people with serious illness.

\medterm{Bleb} A small thin-walled pocket containing air or fluid. In the lung, a bleb can rupture and cause a collapsed lung; on the skin, the term may describe a blister-like lesion.

\medterm{Bleeder's Disease} Alternate name; see \textit{Hemophilia}.

\medterm{Bleeding} The escape of blood from blood vessels. It may be visible outside the body or occur inside a tissue, organ, or body cavity; severe, unexplained, or persistent bleeding needs medical attention.

\medterm{Bleeding Disorder} A condition causing unusually heavy, prolonged, or spontaneous bleeding because of problems with clotting proteins, platelets, blood vessels, or medicines. Examples include hemophilia, von Willebrand disease, platelet disorders, and some forms of disseminated intravascular coagulation.

\medterm{Bleeding Disorders} Alternate name; see \textit{Bleeding Disorder}.

\medterm{Bleeding During Pregnancy} Vaginal bleeding at any time during pregnancy. Causes range from minor cervical changes to miscarriage, placental problems, or labour; heavy bleeding, pain, dizziness, or reduced fetal movement requires urgent assessment.

\medterm{Bleeding Esophageal Varices} Severe bleeding from enlarged, fragile veins in the lower swallowing tube. These veins usually develop when liver disease raises pressure in veins entering the liver. Vomiting blood or passing black stools is an emergency.

\medterm{Bleeding Gums} Bleeding from the gum tissue, often caused by plaque-related gingivitis, periodontitis, trauma, or vigorous brushing.
It can also occur with blood disorders or medicines that affect clotting. Persistent bleeding, swelling,
loose teeth, or spontaneous bleeding requires dental or medical evaluation.

\medterm{Bleeding Into A Thyroid Nodule} Alternate name; see \textit{Hemorrhagic Thyroid Nodule}.

\medterm{Bleeding Tendency} An increased susceptibility to bleeding, which may result from platelet disorders, coagulation factor deficiencies, vascular abnormalities, or anticoagulant medications. Evaluation includes personal and family history, complete blood count, and coagulation studies.

\medterm{Bleeding Time} An older test measuring how long a small skin puncture takes to stop bleeding. It mainly reflected platelet function and small blood vessels and has largely been replaced by more reliable laboratory tests.

\medterm{Bleeding Varices} Enlarged fragile veins, usually in the swallowing tube or stomach, that have ruptured and are bleeding. Vomiting blood, black stools, weakness, or faintness indicates a life-threatening emergency.

\medterm{Blennorrhea} Thick discharge containing mucus and pus from a moist body surface, especially the urethra or eye. Causes include infection and inflammation, so the affected site and other symptoms guide testing and treatment.

\medterm{Bleomycin} A cancer medicine that damages tumour-cell DNA and is used for selected cancers. Its most important serious risk is lung injury; fever, skin changes, low blood counts, and allergic reactions can also occur.

\medterm{Bleomycin Sulfate} The sulphate form of bleomycin, a cancer medicine that damages DNA. It can cause potentially permanent lung injury, so breathing symptoms and lung function may need monitoring during treatment.

\medterm{Blepharitis} A common, usually chronic inflammation of the eyelid margins that can cause itchy or burning eyelids, redness, tearing, dry-eye symptoms, and crusting at the base of the eyelashes. Anterior blepharitis affects the skin, eyelash follicles, and lash bases; posterior blepharitis affects the meibomian oil glands. It is associated with bacterial overgrowth, dandruff, rosacea, allergies, or gland dysfunction and is not usually contagious.

\synonyms
Eyelid Margin Disease; Meibomian Gland Dysfunction

\medterm{Blepharoconjunctivitis} Inflammation of the eyelid edges and the surface covering the eye. It can cause crusting, itching, redness, tearing, and irritation due to infection, allergy, dry eye, or blocked eyelid glands.

\medterm{Blepharophimosis} Abnormally short horizontal openings between the eyelids. It may occur with drooping eyelids and other facial or developmental features and can interfere with vision.

\medterm{Blepharoplasty} Surgery to reshape or repair an eyelid, sometimes removing or repositioning excess skin, muscle, or fat. It may improve vision blocked by heavy lids, restore an injured lid, or be performed for appearance.

\medterm{Blepharoptosis} Drooping of the upper eyelid because of a weak muscle, damaged nerve, stretched tendon, birth difference, ageing, or a heavy eyelid. It can obstruct vision and may signal a neurological emergency if sudden.

\medterm{Blepharospasm} Involuntary tightening of the muscles around the eyes, causing repeated blinking or forceful eyelid closure. It may occur alone or with eye irritation, medicines, or a movement disorder.

\medterm{Blind} Having severe or complete loss of vision that substantially limits the ability to see. The cause and remaining level of vision are described because blindness varies in degree and may affect one or both eyes.

\medterm{Blind Loop Syndrome} Poor absorption caused by excess bacteria growing in a slow-moving or bypassed part of the small intestine. It may cause bloating, diarrhoea, weight loss, and vitamin deficiencies and can follow surgery or bowel disease.

\medterm{Blind Spot} An area within the field of vision where sight is absent. Everyone has a normal small blind spot where the optic nerve leaves the retina, but a new or enlarged blind spot may indicate eye or nerve disease.

\medterm{Blinded} In a research study, describing a participant, clinician, assessor, or analyst who does not know the assigned treatment or relevant result. This helps reduce expectations from influencing care or measurements.

\medterm{Blinding} Keeping treatment assignments or study information hidden from selected participants, clinicians, assessors, or analysts to reduce expectation and measurement bias. A study states exactly who is unaware and what is concealed.

\medterm{Blindness} Severe or complete loss of vision in one or both eyes caused by disease, injury, or a developmental condition affecting the eye, optic nerve, or brain. Sudden vision loss is an emergency.

\medterm{Blindness And Vision Loss} Reduced or absent ability to see caused by problems in the eye, optic nerve, brain, injury, or development. It may be sudden or gradual, partial or complete, and temporary or permanent.

\medterm{Blinking} Repeated brief closing and reopening of the eyelids that spreads tears, protects the cornea, and clears particles. It may increase with irritation, stress, tics, or vision problems and decrease during prolonged screen use.

\medterm{Blister} A raised pocket of fluid within or beneath the outer layer of skin. It can result from rubbing, burns, heat, infection, allergy, or a skin disease and commonly affects hands and feet.

\synonyms
Bulla; Vesicle When Small

\medterm{Bloating} A feeling of abdominal fullness, tightness, or visible swelling. Gas, constipation, food intolerance, or digestive disorders are common causes; persistent bloating with weight loss, vomiting, pain, or early fullness needs assessment.

\medterm{Block} An obstruction or interruption of flow, movement, electrical signals, or communication. The meaning depends on the structure involved, such as an airway, nerve, blood vessel, or heart.

\medterm{Blockage Of Upper Airway} Obstruction of airflow through the nose, mouth, throat, or voice box. It can rapidly cause suffocation; noisy breathing, inability to speak, blue skin, or severe distress requires immediate emergency help.

\medterm{Blocked Tear Duct} Partial or complete blockage of the passage that drains tears from the eye into the nose. It can cause constant watering, discharge, or repeated infection and is common in babies.

\medterm{Blocked Ureter} Alternate name; see \textit{Ureteral Obstruction}.

\medterm{Blocking Agent} A substance that prevents or reduces a biological process by blocking a receptor, channel, enzyme, or other target. Its effects and risks depend on the target and the tissue affected.

\synonyms
Blocker; Antagonist When It Acts by Receptor Antagonism

\medterm{Blocking Agents} Substances that reduce a body process by blocking a receptor, enzyme, ion channel, nerve signal, or other pathway. Their effects and risks vary according to the target and medicine.

\medterm{Blocking Antibody} An immune protein that prevents another molecule from attaching to its target or interferes with a biological interaction. It may be used as treatment, protect against disease, or contribute to an autoimmune disorder.

\medterm{Blood} A body fluid made of plasma, red blood cells, white blood cells, and platelets. It carries oxygen, nutrients, hormones, and waste and helps fight infection, stop bleeding, and regulate temperature.

\medterm{Blood Alcohol Concentration} The amount of alcohol in a person's blood, usually expressed as a percentage or mass per volume. Effects depend on the level, timing, tolerance, other substances, and health; one result does not fully predict impairment.

\medterm{Blood Alcohol Level} A measurement of alcohol concentration in blood. Its effects depend on the amount, timing, tolerance, other substances, and individual health, so the result must be interpreted with the person's condition.

\medterm{Blood Bank} A specialized medical facility and laboratory service that collects, tests, separates, and safely stores donated blood and its components (such as red blood cells, platelets, and plasma). It performs vital compatibility testing, blood typing, and antibody screening (crossmatching) before transfusions to prevent severe transfusion reactions and ensure patient safety.

\synonyms
Blood Banking; Transfusion Service; Immunohematology Laboratory

\medterm{Blood Banking} Alternate name; see \textit{Blood Bank}.

\medterm{Blood Biomarker} A measurable sign of a disease or condition that can be isolated or detected in a blood sample, such as a disease-associated antibody, hormone, protein, or other substance.

\medterm{Blood Blister} A raised pocket under the skin containing blood, usually caused by rubbing, pinching, pressure, or another injury that breaks small blood vessels.

\medterm{Blood Brain Barrier} A protective filter formed by specialized blood vessels and nearby brain cells. It controls what can pass from the bloodstream into the brain, limiting many medicines, chemicals, germs, and immune cells.

\medterm{Blood Cancer} A cancer affecting blood-forming tissue, blood cells, or the lymphatic system. Abnormal cells can crowd out normal cells, impair immunity, cause anaemia or bleeding, or spread to organs.

\medterm{Blood Cancer} Alternate name; see \textit{Leukemia}.

\medterm{Blood Cell} A cellular part of blood, including red blood cells that carry oxygen, white blood cells that help fight infection, and platelets that help stop bleeding. Most are made in the bone marrow.

\medterm{Blood Cholesterol} Cholesterol carried in the blood. The body needs it to build cell membranes and make hormones and bile acids. Too much LDL cholesterol can form fatty buildup in the arteries and increase the risk of heart disease and stroke. Blood cholesterol is usually checked as part of a lipid profile.

\medterm{Blood Circulation} The continuous movement of blood from the heart through the lungs and body and back again. It delivers oxygen and nutrients and carries away carbon dioxide and other waste.

\medterm{Blood Cleaner} A method used to reduce certain infectious organisms in donated blood or blood components while keeping them suitable for transfusion. It may use chemicals or light, but it does not remove every possible infection risk.

\medterm{Blood Clot} A mass of thickened blood formed during clotting. It is protective when it seals an injury, but a clot inside a blood vessel can block flow or travel to the lungs, brain, heart, or another organ.

\medterm{Blood Clot in Lung} Alternate name; see \textit{Pulmonary Embolism}.

\medterm{Blood Clots} Clumps of thickened blood. Clots can stop bleeding, but abnormal clots in veins or arteries may block blood flow and cause deep-vein thrombosis, pulmonary embolism, stroke, heart attack, or tissue damage.

\synonyms
Thrombi; Thrombosis; Blood Clot

\medterm{Blood Coagulation} The linked process that changes blood from a liquid to a firm clot after injury. Too little clotting causes bleeding, while excessive or misplaced clotting can cause dangerous vessel blockage.

\medterm{Blood Coagulation Disorder} A condition in which blood clots too little, too much, or at the wrong time. Bleeding disorders result from inadequate clot formation, whereas thrombophilic disorders cause excessive or misplaced clotting. Causes include genetic variants, medicines, liver disease, immune conditions, cancer, and other illness. Examples include hemophilia, von Willebrand disease, thrombocytopenia, thrombophilia, and disseminated intravascular coagulation.

\medterm{Blood Count} A blood test measuring the number and features of red cells, white cells, and platelets. It helps assess anaemia, infection, inflammation, bleeding risk, and disorders of the bone marrow.

\medterm{Blood Culture} A test in which blood is placed in special bottles to see whether bacteria or fungi grow. It helps investigate bloodstream infection, sepsis, or infection of the heart valves.

\medterm{Blood Culture Contamination} Accidental entry of skin or environmental organisms into a blood-culture bottle during collection or handling, rather than true infection in the blood. It can lead to unnecessary antibiotics, so results are interpreted with symptoms and other culture sets.

\medterm{Blood Differential Test} A blood test measuring the numbers and proportions of different white blood cells, including neutrophils, lymphocytes, monocytes, eosinophils, and basophils. It can help investigate infection, inflammation, allergy, leukaemia, or marrow disease.

\medterm{Blood Donation} Voluntary collection of whole blood or selected blood components from a screened donor for transfusion or manufacture of medical products. Eligibility checks, infection testing, sterile collection, and post-donation recovery protect donor and recipient.

\medterm{Blood Donation Before Surgery} Giving and storing a person's own blood before an operation for possible later transfusion. It is suitable only in selected situations and can cause anaemia or delays if the operation is postponed.

\medterm{Blood Donor} A person who gives blood or blood components for transfusion or medical use. Screening protects the donor and helps ensure that donated blood is safe for the recipient.

\medterm{Blood Draw} Collection of blood, usually with a needle placed in a vein, for testing, donation, or treatment. Temporary pain, bruising, dizziness, or rarely infection may occur.

\medterm{Blood Film} A thin layer of peripheral blood spread onto a glass microscope slide, stained with a Romanowsky stain (e.g., Wright-Giemsa), and examined under microscopy for quantitative and qualitative evaluation of erythrocyte morphology, leukocyte differentials, platelet adequacy, and identification of bloodborne parasites.

\synonyms
Peripheral Blood Smear; Blood Smear; Peripheral Smear

\medterm{Blood Flow} The movement of blood through vessels and organs. It depends mainly on the pressure difference and resistance in the vessels and delivers oxygen and nutrients while removing waste.

\medterm{Blood Gas Analysis} A test measuring oxygen, carbon dioxide, acidity, and related values in blood. It helps assess breathing, oxygen delivery, and acid-base balance in serious illness.

\medterm{Blood Gas, Venous} The clinical laboratory measurement of pH, partial pressure of carbon dioxide ($p\text{CO}_2$), bicarbonate ($\text{HCO}_3^-$), and base excess in a peripheral or central venous blood sample to assess systemic acid-base balance, ventilatory status, and tissue perfusion.

\synonyms
Venous Blood Gas; VBG; Central Venous Blood Gas

\medterm{Blood Gases} Measurements of oxygen, carbon dioxide, acidity, and related substances in blood. They help determine how well a person is breathing, exchanging gases, and maintaining acid-base balance.

\medterm{Blood Glucose} The amount of sugar called glucose in the blood. The body regulates it using food absorption, liver production, insulin, and other hormones; levels that are too low or high can cause illness.

\medterm{Blood Glucose Monitoring} The regular measurement of sugar (glucose) levels in the blood, primarily used to manage diabetes and detect dangerously high (hyperglycemia) or low (hypoglycemia) blood sugar. Testing is commonly performed using a handheld glucometer with a small drop of blood from a fingertip prick, or with a wearable continuous glucose monitor (CGM) that checks glucose in the fluid under the skin around the clock. It tracks how meals, physical activity, illness, and medications (such as insulin) affect glucose levels throughout the day.

\synonyms
Blood Sugar Monitoring; Self-Monitoring of Blood Glucose; SMBG; Continuous Glucose Monitoring; CGM

\medterm{Blood Group} A classification of blood based on specific inherited markers (antigens) on the surface of red blood cells. The two most vital systems are the ABO system (types A, B, AB, and O) and the Rhesus (Rh) system (positive or negative based on the RhD marker). Accurate blood grouping and crossmatching are essential before transfusions to prevent life-threatening immune reactions and during pregnancy to prevent Rh disease (hemolytic disease of the newborn).

\synonyms
Blood Type; Blood Groups; ABO and Rh Grouping; Blood Typing

\medterm{Blood Group, ABO} Alternate name; see \textit{ABO Blood Group System}.

\medterm{Blood Groups} Alternate name; see \textit{Blood Group}.

\medterm{Blood In Semen} Blood in the fluid released during ejaculation. It is often temporary and may follow infection, inflammation, injury, or a urinary-tract procedure. Persistent or recurrent blood, pain, fever, or urinary symptoms needs medical assessment.

\synonyms
Hematospermia; Hemospermia

\medterm{Blood in Stool} Alternate name; see \textit{Rectal Bleeding}.

\medterm{Blood In The Eye} Bleeding on the eye surface or inside the eye. A small surface bleed may be harmless, but bleeding inside the eye, eye pain, injury, or any vision change requires prompt eye assessment.

\medterm{Blood in Urine} The abnormal presence of red blood cells in the urine, detected either macroscopically as pink, red, or tea-colored urine (gross hematuria) or microscopically on urinalysis ($\ge 3$ erythrocytes per high-power field). Common etiologies include urinary tract infections, nephrolithiasis, glomerular nephritis, benign prostatic hyperplasia, and urothelial malignancies.

\synonyms
Hematuria; Gross Hematuria; Microscopic Hematuria; Blood in the Urine

\medterm{Blood Lancet} Alternate name; see \textit{Lancet}.

\medterm{Blood Lead Level Test} A quantitative diagnostic laboratory assay measuring venous or capillary whole-blood lead concentration in micrograms per deciliter ($\mu\text{g/dL}$). Utilized for pediatric screening, occupational exposure monitoring, and evaluating symptomatic plumbism.

\synonyms
BLL Assay; Venous Blood Lead Screen; Plumbism Blood Test
\medterm{Blood Loss} Loss of blood from the circulation because of injury, surgery, childbirth, internal bleeding, or disease. Severe loss can reduce oxygen delivery and cause shock, requiring urgent control of bleeding and replacement of blood or fluids.

\medterm{Blood Low Red Cell Count} Alternate name; see \textit{Anemia}.

\medterm{Blood PH} A measure of how acidic or alkaline the blood is. The lungs regulate carbon dioxide and the kidneys regulate acids and bicarbonate; normal arterial pH is about 7.35--7.45.

\medterm{Blood Plasma} The liquid part of blood, made mostly of water and containing proteins, salts, nutrients, hormones, gases, and waste products. It carries blood cells and dissolved substances throughout the body.

\medterm{Blood Poisoning} A non-technical, lay term historically used to describe a severe bloodstream bacterial infection (septicemia). In modern medicine, this condition is formally diagnosed and treated as \textit{Sepsis}---a life-threatening condition in which the body's overwhelming response to an infection causes organ dysfunction.

\synonyms
Septicemia; Sepsis; Bloodstream Infection; Bacteremia

\medterm{Blood Preservation} The controlled collection, processing, cooling, freezing, and storage of donated blood or its components so they remain safe and effective until transfusion.

\medterm{Blood Pressure} The force of circulating blood against artery walls, recorded as two numbers in millimetres of mercury. The systolic number is the pressure when the heart contracts; the diastolic number is the pressure between beats.

\medterm{Blood Pressure Chart: Reading By Age} A guide for interpreting blood-pressure readings in relation to age and clinical circumstances. Diagnosis usually requires repeated, correctly measured readings; pregnancy, illness, and treatment goals can change the appropriate range.

\medterm{Blood Pressure Cuff} The inflatable part of a blood-pressure measuring device that wraps around a limb and briefly compresses an artery. Cuffs come in different sizes and may be used with a manual sphygmomanometer or an automated monitor.

\medterm{Blood Pressure Measurement} Measurement of arterial pressure with an inflatable cuff or another approved device. The result records systolic pressure when the heart contracts and diastolic pressure when it relaxes; correct cuff size and positioning improve accuracy.

\medterm{Blood Pressure Monitors For Home} Devices that measure blood pressure outside a clinic, usually with an inflatable upper-arm cuff. A validated monitor, correctly sized cuff, proper positioning, and repeated readings give more reliable results.

\medterm{Blood Pressure, Low} A systolic blood pressure below $90\text{ mmHg}$ or diastolic blood pressure below $60\text{ mmHg}$, resulting in inadequate end-organ perfusion. Etiologies include hypovolemia, septic shock, cardiogenic impairment, adrenal insufficiency, orthostatic autonomic failure, and antihypertensive overdosage.

\synonyms
Low Blood Pressure; Hypotension; Orthostatic Hypotension

\medterm{Blood Product} A component made from donated blood for medical treatment, such as red blood cells, platelets, plasma, or cryoprecipitate. The product is selected according to the person's bleeding, anaemia, or clotting problem.

\medterm{Blood Products} Components prepared from donated blood for transfusion, including red blood cells, platelets, plasma, and cryoprecipitate. Each replaces a different blood function, and the choice depends on the person's symptoms, bleeding, and laboratory results.

\medterm{Blood Protein} A protein dissolved in blood plasma. These proteins help transport substances, maintain fluid balance, support immunity, and control clotting; abnormal amounts can occur with inflammation, liver or kidney disease, immune disorders, or some cancers.

\medterm{Blood Results} Measurements obtained from a blood sample, such as cell counts, glucose, salts, hormones, or organ-function markers. They must be interpreted using the laboratory's reference range, symptoms, medicines, and other findings.

\medterm{Blood Sedimentation} The settling of red blood cells in a tube of treated blood. The erythrocyte sedimentation rate is a nonspecific sign that may rise with inflammation, infection, anaemia, pregnancy, or other conditions.

\medterm{Blood Smear} A thin layer of blood spread on a glass slide, stained, and examined under a microscope. It shows the size, shape, and appearance of blood cells and can reveal abnormal cells, parasites, or cell fragments.

\medterm{Blood Spatter Analysis} The forensic examination of bloodstain size, shape, and distribution to help reconstruct events. It may distinguish drops caused by gravity, contact transfer, impact, or forceful expulsion, but conclusions have limits.

\medterm{Blood Substitutes} Fluids or artificial products intended to replace some functions of blood. Most restore circulating volume; some are designed to carry oxygen. They do not replace every function of blood and are used only in selected situations.

\medterm{Blood Sugar} The amount of glucose in the blood. Insulin, glucagon, the liver, food, activity, illness, and other hormones affect it. Persistently high or very low levels can cause symptoms and serious organ injury.

\medterm{Blood Sugar Test} A test measuring glucose in the blood. It may be done while fasting, at a random time, or after a glucose drink, and helps diagnose or monitor diabetes and other disorders of glucose control.

\medterm{Blood Test} Analysis of blood cells or substances dissolved in blood, including chemicals, proteins, hormones, antibodies, or genetic material. Results help investigate symptoms, diagnose disease, assess organ function, and monitor treatment.

\medterm{Blood Thinner} A common name for an anticoagulant medicine that lowers the blood's ability to form harmful clots. It does not actually make the blood thinner. These medicines help prevent or treat conditions such as deep vein thrombosis, pulmonary embolism, and stroke, but can increase bleeding risk.

\synonyms
Anticoagulant

\medterm{Blood Transfusion} The administration of whole blood or a specific blood component, such as red blood
cells, platelets, plasma, or cryoprecipitate, to treat blood loss, anaemia, thrombocytopenia,
or a clotting-factor deficiency. Red cells improve oxygen-carrying capacity; platelets,
plasma, and cryoprecipitate support haemostasis. ABO/Rh typing, antibody screening,
and compatibility testing help reduce transfusion reactions. Risks include febrile or
allergic reactions, acute haemolysis, circulatory overload, and rare transfusion-transmitted
infection; the decision depends on the patient's symptoms, laboratory results, bleeding,
and overall clinical context.

\medterm{Blood Typing} A laboratory test identifying markers on red blood cells, especially the ABO and Rh systems. It helps select compatible blood for transfusion and identifies some pregnancy-related risks from blood-group incompatibility.

\medterm{Blood Urea Nitrogen} The amount of nitrogen from urea in the blood. Urea is made when the body breaks down protein and is removed mainly by the kidneys; the result is interpreted with creatinine, hydration, and the clinical situation.

\medterm{Blood Urea Nitrogen Test} A blood test measuring urea nitrogen, a waste product made during protein breakdown. It helps assess kidney function and hydration, but abnormal results can also reflect diet, bleeding, medicines, or other illness.

\medterm{Blood Urea Nitrogen Test} A clinical serum chemistry test measuring the amount of nitrogen derived from urea, the principal end product of protein catabolism formed in the liver. Elevated serum concentrations (azotemia) occur in acute or chronic renal impairment, prerenal hypoperfusion, severe dehydration, and upper gastrointestinal hemorrhage.

\synonyms
BUN Test; Serum Blood Urea Nitrogen; Urea Nitrogen Test

\medterm{Blood Vessel} A tube that carries blood through the body. Arteries carry blood away from the heart, veins return it, and tiny capillaries allow oxygen, nutrients, and waste products to pass between blood and tissues.

\medterm{Blood Vessels} The arteries, arterioles, capillaries, venules, and veins that carry blood through the body. They deliver oxygen and nutrients, remove waste, and help control blood pressure, temperature, and fluid exchange.

\medterm{Blood Viscosity} How thick blood is and how much it resists flowing. It is influenced mainly by the number of red blood cells, but also by blood proteins and temperature. Excessively thick blood can slow circulation and increase clot risk.

\medterm{Blood Volume} The total amount of blood circulating in the body, including plasma and blood cells. Too little can cause low blood pressure and shock; too much can strain the heart and lungs.

\medterm{Blood-Borne Disease} An infection that can spread through infected blood or certain body fluids, for example through shared needles, needlestick injuries, sexual exposure, or contaminated medical equipment. Examples include HIV and hepatitis B and C.

\medterm{Blood-Borne Pathogens} Germs that can spread through infected blood or certain body fluids. Examples include HIV and hepatitis B and C viruses; prevention includes vaccination where available, gloves, safe sharps handling, and appropriate exposure testing.

\medterm{Blood-Brain Barrier} A protective filter formed by specialized cells lining brain blood vessels and their supporting tissues. It controls which medicines, chemicals, germs, and immune cells can move from the blood into brain tissue.

\medterm{Bloodborne Pathogens} Infectious agents spread through infected blood or certain body fluids. Important examples include HIV and hepatitis B and C viruses; exposure may require urgent medical assessment, testing, and preventive treatment.

\medterm{Bloom Syndrome} A rare inherited disorder caused by changes in both copies of the BLM gene. It causes short stature, a sun-sensitive facial rash, frequent infections, and a greatly increased risk of several cancers.

\medterm{Blount Disease} A growth disorder in which the upper part of the shinbone grows unevenly, causing worsening bowing of one or both legs. It can affect young children or adolescents and may cause an abnormal gait, unequal leg length, or knee damage.

\medterm{Blue Baby} An infant whose skin, lips, or nail beds look blue because the blood carries too little oxygen or circulation is inadequate. This can signal heart, lung, airway, blood, or other serious disease and needs urgent assessment.

\medterm{Blue Bloater} Outdated descriptive label; use the specific diagnosis, such as \textit{Chronic Bronchitis} or \textit{Chronic Obstructive Pulmonary Disease}.

\medterm{Blue Bottle} A common name for a bluebottle jellyfish. Its tentacles can inject venom that causes immediate stinging pain, raised marks, redness, and sometimes muscle cramps, breathing difficulty, or a severe allergic reaction.

\medterm{Blue Discoloration Of The Skin} A blue or gray tint of the skin, lips, or nail beds caused by low oxygen, reduced blood flow, or abnormal haemoglobin. Sudden or widespread discoloration, especially with breathing difficulty or confusion, is an emergency.

\medterm{Blue Nevi} Usually harmless moles formed by pigment-producing cells deeper in the skin, giving them a blue, blue-gray, or blue-black color. A new, enlarging, or changing lesion requires examination because melanoma can sometimes look similar.

\medterm{Blue Nevus} A usually harmless blue-gray or blue-black mole caused by pigment-producing cells located deeper than those in most moles. A new, changing, bleeding, or painful lesion requires examination.

\medterm{Blue Nightshade Poisoning} Illness caused by eating blue nightshade or a related poisonous plant containing glycoalkaloids. Symptoms may include vomiting, diarrhoea, abdominal pain, confusion, weakness, seizures, or abnormal heart rhythm; it is an emergency requiring urgent medical treatment.

\medterm{Blue NæVus} A usually harmless blue-gray or blue-black mole formed by pigment-producing cells deep in the skin. It requires assessment if it changes in size, shape, colour, or symptoms.

\medterm{Blue Sclera} A blue-gray appearance of the normally white outer layer of the eye, usually caused by thinning that allows deeper pigment to show through. It can occur in inherited connective-tissue or bone disorders and needs clinical assessment.

\medterm{Blue Sclerae} A distinct bluish or grayish discoloration of the ocular sclera caused by abnormal thinness or translucency of the scleral collagen fibers, which allows the pigment of the underlying uveal tract and choroid to be visible. It is a hallmark physical finding of osteogenesis imperfecta and can also occur in Ehlers-Danlos syndrome, Marfan syndrome, and severe iron deficiency anemia.

\medterm{Blue-Green Algae} Blue-green algae are cyanobacteria that can grow rapidly in water. Some release toxins that cause skin irritation, vomiting, liver injury, breathing problems, or neurologic illness when people or animals are exposed.

\medterm{Blumberg Sign} Sudden sharp pain when a doctor presses down gently on the belly and then quickly lets go, showing that the lining of the abdomen is swollen and irritated.

\medterm{Blunt Abdominal Trauma} An injury to the abdomen without an object piercing the skin, often caused by a crash, fall, or blow. It can damage the spleen, liver, kidneys, bowel, or pancreas, sometimes without early visible bruising.

\medterm{Blunt Force Trauma} Injury caused by a blow from a blunt object or by striking a blunt surface, without penetration by a sharp edge. It can produce bruises, grazes, torn skin, internal bleeding, fractures, or organ damage.

\medterm{Blunted Affect} Reduced outward expression of emotion. A person may show less facial movement, voice variation, or body language than expected, although their internal feelings may be stronger. It can occur in mental or neurologic disorders.

\medterm{Blurred Vision} Vision that appears out of focus, hazy, or unclear. Causes include an incorrect glasses prescription, dry eye, eye disease, medicines, migraine, or neurologic illness. Sudden blurring, pain, weakness, or vision loss needs urgent assessment.

\medterm{Blush} Temporary facial redness or warmth caused by increased blood flow near the skin surface. It may follow emotion, heat, exercise, fever, alcohol, medicines, hormonal changes, or some medical conditions.

\medterm{Blushing} Temporary redness or warmth of the face caused by increased blood flow in the skin. Common triggers include emotion, heat, exercise, alcohol, hormonal changes, medicines, fever, and some illnesses.

\medterm{Board-Like Rigidity} Severe involuntary tightening of the abdominal muscles, making the abdomen feel hard. It can indicate irritation of the abdominal lining from serious inflammation, infection, bleeding, or perforation and requires urgent assessment.

\medterm{Boas' Sign} Pain or unusual skin sensitivity below the right shoulder blade, traditionally associated with gallbladder disease. It is not specific enough to establish a diagnosis and must be interpreted with other findings.

\medterm{Bobbing, Ocular} Alternate index label; see \textit{Ocular Bobbing}.

\medterm{Bobble-Head Doll Syndrome} A rare childhood movement disorder causing repeated forward-and-back head movements. It is often linked to a fluid-filled cyst or other lesion near the third ventricle of the brain and requires neurologic evaluation and imaging.

\medterm{Bocavirus} A group of small viruses that can infect humans and animals. Human bocavirus most often affects the airways of young children, although infection may cause illness in other tissues or people with weakened immunity.

\medterm{Bocavirus Infection} A viral infection, usually of the airways in young children. It may cause a runny nose, cough, fever, wheezing, bronchiolitis, or pneumonia, often together with another infection. Treatment is usually supportive; breathing difficulty or dehydration needs prompt care.

\medterm{Bochdalek Hernia} A birth defect in which abdominal organs move through an opening in the back and side of the diaphragm into the chest. It can prevent normal lung development and cause severe breathing problems in a newborn.

\medterm{Bock's Disease} An older name for sarcoidosis, an inflammatory disease in which clusters of immune cells form in organs. It most often affects the lungs and lymph nodes but can also affect the skin, eyes, heart, or nervous system.

\medterm{BODE Index} A scoring tool for chronic obstructive pulmonary disease that combines body weight, breathing-test results, breathlessness, and walking distance. A higher score suggests more severe disease and a greater risk of hospital admission or death; it does not predict an individual outcome with certainty.

\synonyms
BODE Score; COPD Multidimensional Index

\medterm{Bodily Secretion} A substance made and released by cells or glands, such as sweat, saliva, mucus, tears, digestive juices, or hormones. Secretions protect tissues, aid digestion, regulate body functions, or help remove waste.

\medterm{Body} The physical organism, including its cells, tissues, organs, and interacting body systems. Health depends on the structure and function of these parts and how they respond to internal and external conditions.

\medterm{Body Aches} Widespread muscle or body pain that can occur with infection, fever, exercise, inflammation, medicines, or other illness. Severe pain with weakness, dark urine, breathing difficulty, chest pain, or high fever needs prompt assessment.

\medterm{Body Cavity} An enclosed space inside the body that contains and protects organs. Major cavities include the skull, spinal canal, chest, abdomen, and pelvis; their linings may produce small amounts of lubricating fluid.

\medterm{Body Cell Mass} The living, metabolically active cells of the body, including cells in muscles, organs, blood, and other tissues. It excludes stored fat and water outside cells and is sometimes estimated in nutrition assessment.

\medterm{Body Composition} The relative amounts of fat, muscle and other lean tissue, bone mineral, and water in the body. It provides more information about health and nutrition than body weight alone, although estimates vary by method.

\medterm{Body Dysmorphia} Alternate name; see \textit{Body Dysmorphic Disorder}.

\medterm{Body Dysmorphic Disorder} A mental-health disorder involving persistent distress about one or more perceived flaws in appearance that are minor or not noticeable to others. Repetitive checking, grooming, comparison, or avoidance interferes with daily life.

\synonyms
BDD

\medterm{Body Fat} Adipose tissue that stores energy, cushions organs, insulates the body, and releases hormones and other signals. Too much or too little body fat can affect metabolism, movement, fertility, immunity, and other aspects of health.

\medterm{Body Fluid} The water and dissolved substances (such as electrolytes, nutrients, proteins, and gases) that make up approximately 60\% of human body weight. Body fluids are organized into two major compartments: intracellular fluid (fluid inside cells, making up about two-thirds) and extracellular fluid (fluid outside cells, including interstitial fluid between tissues and blood plasma), along with specialized fluids such as cerebrospinal fluid, synovial joint fluid, lymph, and digestive secretions. They are essential for transporting nutrients and waste, regulating body temperature, and lubricating tissues and joints.

\synonyms
Body Fluids; Bodily Fluids; Total Body Water; Body Fluid Compartments

\medterm{Body Fluids} Alternate name; see \textit{Body Fluid}.

\medterm{Body Habitus} A person's general build, including body size, shape, muscle, fat distribution, and proportions. Clinicians may record it during an examination, but it does not by itself diagnose disease or predict health.

\medterm{Body Height} The vertical distance from the soles of the feet to the top of the head, measured while standing when possible. It helps assess growth and nutrition and is used in medication, kidney-function, and other calculations.

\medterm{Body Identification} The process of establishing who a deceased person is. It may use fingerprints, dental records, medical information, DNA, personal effects, or visual recognition, with the method chosen according to the circumstances.

\medterm{Body Image} A person's thoughts, feelings, and beliefs about their body and appearance. Distress or inaccurate beliefs about appearance can occur in body dysmorphic disorder, eating disorders, depression, or other mental-health conditions.

\medterm{Body Lice} Infestation with human body lice, which live mainly in clothing and move to the skin to feed. Bites cause intense itching and scratching; treatment includes washing or heat-treating clothing and bedding and prescribed insecticide when needed.

\synonyms
Pediculosis Corporis; Body Louse Infestation

\medterm{Body Mass Index} A screening measure calculated by dividing weight in kilograms by height in meters squared. It helps assess weight in relation to height, but it does not directly measure body fat or provide a complete assessment of health. Results are interpreted with age and other clinical factors; children and teens are assessed using age- and sex-specific percentiles.

\synonyms
BMI; Quetelet Index; BMI Calculation

\medterm{Body Odor} Smell produced when skin secretions, especially sweat, are broken down by skin microorganisms. A sudden or unusual change may be related to diet, medicines, infection, or a metabolic disorder.

\medterm{Body Odor and Sweating} Alternate index label; see \textit{Sweating and Body Odor}.

\medterm{Body Odour} Smell from the breakdown of skin secretions by microorganisms. The smell can change with sweating, hygiene, diet, medicines, infection, or some metabolic diseases.

\medterm{Body Piercing} Creation of an opening through skin or a mucous membrane to insert jewellery. Possible complications include pain, bleeding, infection, scarring, allergic reaction, and damage to teeth, nerves, or cartilage.

\medterm{Body Plethysmography} A breathing test performed inside a sealed cabin to measure lung volumes, including air that cannot be exhaled. It helps identify restricted lungs, over-expanded lungs, and some airway problems.

\medterm{Body Position} The arrangement of the body or a body part, such as standing, lying on the back, lying face down, or lying on the side. Position can affect breathing, circulation, examination results, and medical procedures.

\medterm{Body Region} A named area of the body used to describe the location of an organ, symptom, injury, or examination finding, such as the chest, abdomen, pelvis, or arm.

\medterm{Body Ringworm} A fungal infection of the skin on the body. It usually causes an expanding, ring-shaped, scaly rash with a more active edge and clearer centre; it spreads by contact and is treatable with antifungal medicine.

\synonyms
Tinea Corporis; Dermatophytosis of Glabrous Skin

\medterm{Body Space} A normal or potential compartment inside the body, such as the space around the lungs or abdominal organs. Air, blood, fluid, or infection can collect there and interfere with organ function.

\medterm{Body Substance} Material produced by or found in the body, including blood, saliva, urine, faeces, mucus, tissue, and other secretions. Some substances can carry infection and require appropriate handling precautions.

\medterm{Body Surface} The outside of the body, mainly the skin and accessible mucous membranes. Clinicians may estimate its area when assessing burns, calculating some medicine doses, or making physiological measurements.

\medterm{Body Surface Area} An estimate of the body's external area, usually calculated from height and weight. It is used in burn assessment, some medicine doses, and selected medical calculations; it is an estimate rather than a direct measurement.

\medterm{Body Temperature} The degree of heat in the body, determined by the balance between heat production and heat loss. It varies with age, activity, time of day, illness, and measurement site; fever and hypothermia require clinical interpretation.

\medterm{Body Temperature Norms} Expected temperature ranges vary with age, time of day, activity, illness, and whether temperature is measured by mouth, ear, armpit, skin, or rectum. The measurement method and symptoms are considered when interpreting fever or low temperature.

\medterm{Body Wall} The layers of skin, connective tissue, muscle, and sometimes bone that surround and protect a body compartment, such as the chest or abdomen. They also help movement, breathing, and support.

\medterm{Body Weight} The measured mass of a person, including fat, muscle, bone, water, and organs. It helps assess growth, nutrition, and fluid changes and is used with other information for some medical calculations.

\medterm{Boeck's Sarcoid} An older name for sarcoidosis, an inflammatory disease in which clusters of immune cells form in tissues. It commonly affects the lungs and nearby lymph nodes but may also involve the skin, eyes, heart, or nervous system.

\medterm{Boerhaave Syndrome} A sudden full-thickness tear of the food pipe, usually after forceful vomiting. Leakage into the chest can cause severe chest or upper-abdominal pain, infection, shock, and breathing problems; it is a surgical emergency.

\medterm{Boerhaave's Syndrome} A spontaneous full-thickness tear of the food pipe, most often following forceful vomiting. Contents can leak into the chest and cause severe infection, chest pain, breathing difficulty, and shock.

\medterm{Bogros' Space} A layer of tissue behind the groin region and in front of the abdominal lining. Surgeons use this natural plane when repairing some groin hernias.

\medterm{Bohn Nodules} Small, harmless cysts containing keratin that appear on a newborn's gums or the roof of the mouth. They usually disappear without treatment and do not interfere with feeding.

\medterm{Bohr Effect} The tendency of haemoglobin to release more oxygen when tissues become more acidic or contain more carbon dioxide. This helps active tissues receive oxygen where energy use and carbon-dioxide production are greatest.

\medterm{Bohring Syndrome} A rare genetic disorder associated with severe developmental delay, poor growth before birth, distinctive facial features, reduced movement, and frequent feeding or breathing difficulties.

\medterm{Boil} A painful, pus-filled skin abscess centered on an infected hair follicle and surrounding tissue, most commonly caused by \textit{Staphylococcus aureus} bacteria. It typically begins as a tender, red nodule that progressively enlarges, develops a central yellowish head of pus, and eventually ruptures to drain. Multiple or recurrent boils are termed furunculosis, whereas a deeper, connected cluster of several boils is classified as a \textit{Carbuncle}.

\synonyms
Furuncle; Furunculosis; Boils; Skin Abscess

\medterm{Boils} Alternate name; see \textit{Boil}.

\medterm{Boils and Carbuncles} Alternate name; see \textit{Boil} and \textit{Carbuncle}.

\medterm{Bolt Device} A bolt is a threaded fastener used to hold parts together. In medicine, the term may describe a device component or a specific product name, so its meaning depends on the clinical context.

\medterm{Bolus} A single measured amount of medicine or fluid given over a short time, usually by injection or infusion. In digestion, it means food that has been chewed and is ready to swallow.

\medterm{Bolus Gastrostomy Feeding} The administration of prescribed liquid enteral formula or hydration directly into the stomach through a gastrostomy tube over a short period using a syringe or gravity drip. Involves tube flushing before and after each feeding to maintain patency.

\synonyms
Bolus G-Tube Feeding; Bolus Enteral Nutrition; Syringe Feeding

\synonyms
G-Tube Bolus Feeding; Enteral Bolus Nutrition
\medterm{Bolus Infusion} Delivery of a measured amount of fluid or medicine into a vein over a short period. The dose and speed depend on the medicine, the person's condition, and the intended effect; errors can cause serious harm.

\medterm{Bombay Blood Group} A rare inherited blood type in which red blood cells lack the H marker needed to form A or B markers. It may look like type O but is usually compatible only with blood from another Bombay-type donor.

\medterm{Bombesin Receptor} A protein on or inside a cell that binds bombesin-like signalling substances and changes cell activity. These receptors help regulate processes such as hormone release, digestion, and cell growth and are studied in some cancers.

\medterm{Bone} A strong living tissue that forms the skeleton, supports the body, protects organs, enables movement with muscles, stores minerals, and contains marrow that makes blood cells. Bone is continually broken down and rebuilt.

\medterm{Bone-Anchored Hearing Implant} A hearing device fixed to the skull bone that transmits sound vibrations to the inner ear, bypassing the outer or middle ear. It may be used for selected conductive, mixed, or single-sided hearing loss.

\medterm{Bone And Mineral Disorders} Diseases affecting bone strength, bone formation, or the body's balance of calcium and phosphate. Examples include osteoporosis, rickets, osteomalacia, parathyroid disorders, and inherited bone diseases.

\medterm{Bone Callus} New tissue that forms around a broken bone during healing. It first consists of a soft blood clot and connective tissue, then cartilage and immature bone, which are gradually reshaped into stronger bone.

\medterm{Bone Cancer} Cancer that starts in bone, such as osteosarcoma or chondrosarcoma, or cancer that has spread there from another organ. Persistent focal pain, swelling, or a fracture after minor injury may occur; diagnosis uses imaging and tissue testing.

\synonyms
Bone Sarcoma; Primary Bone Malignancy

\medterm{Bone Cement} A material, usually polymethylmethacrylate, used to secure artificial joints or fill spaces in bone. It hardens chemically and may be used in joint replacement, vertebroplasty, kyphoplasty, or bone reconstruction. Possible complications include cement leakage, infection, loosening, and bone cement implantation syndrome.

\medterm{Bone Cement} A self-curing acrylic biomaterial composed of polymethylmethacrylate (PMMA) utilized in orthopedic and spine surgery to anchor prosthetic joint components, fill osseous voids, and stabilize vertebral compression fractures during vertebroplasty or kyphoplasty.

\synonyms
Polymethylmethacrylate; PMMA; Acrylic Bone Cement

\medterm{Bone Conduction} Hearing sound through vibrations carried by the skull to the inner ear, bypassing the outer and middle ear. Bone-conduction testing helps identify whether hearing loss comes from the outer or middle ear or from the inner ear and nerve.

\medterm{Bone Cyst} A fluid-filled or blood-filled space within bone. Many bone cysts are benign, but they can weaken the bone and cause pain, swelling, or fracture. Common types include simple bone cysts and aneurysmal bone cysts.

\medterm{Bone Densitometry} See \textit{Bone Density Screening}.

\medterm{Bone Density} The amount of mineral packed into a measured area of bone. It is usually estimated with a DXA scan; lower density is associated with greater fracture risk, but results must be interpreted with age and other risk factors.

\medterm{Bone Density Screening} A DXA scan that measures bone mineral at sites such as the hip and spine. It helps identify osteoporosis and estimate fracture risk, especially in people with age-related or other risk factors.

\medterm{Bone Development} The formation, growth, and shaping of bones from before birth through adulthood. Bones lengthen and widen as new tissue is made and old tissue is removed; hormones, nutrition, activity, and injury influence this process.

\medterm{Bone Disease} Any condition that weakens, deforms, inflames, infects, or otherwise damages bone. Examples include osteoporosis, fractures, infections, tumours, inherited skeletal disorders, and diseases caused by abnormal mineral or hormone levels.

\medterm{Bone Fracture} A partial or complete break in a bone caused by injury, repeated stress, or weakened bone. Pain, swelling, bruising, deformity, and loss of function may occur. An open fracture or loss of sensation or circulation needs emergency care.

\medterm{Bone Fractures} Alternate name; see \textit{Bone Fracture}.

\medterm{Bone Graft} Bone tissue or a substitute placed in a bone defect to support repair, join bones together, rebuild an area, or help a fracture heal. The material may come from the person, a donor, or a synthetic source.

\medterm{Bone Graft, Dental} An oral and maxillofacial surgical procedure involving the placement of autologous bone, allograft, xenograft, or synthetic osteoconductive matrix into the maxilla or mandible to reconstruct alveolar ridge deficiencies or augment the maxillary sinus prior to dental implant placement.

\synonyms
Dental Bone Graft; Alveolar Ridge Augmentation; Maxillofacial Bone Grafting

\medterm{Bone Health} The strength, structure, and ability of bones to grow, repair, and resist fractures. It is influenced by age, nutrition, hormones, physical activity, medicines, smoking, alcohol, and medical conditions.

\medterm{Bone Infection} Infection of bone, usually caused by bacteria. It may reach bone through the bloodstream, spread from nearby tissue, or enter through surgery or an open injury, causing pain, swelling, fever, or poor healing.

\medterm{Bone Lesion Biopsy} Removal of a small sample from an abnormal area of bone with a needle or operation for laboratory examination. Testing can distinguish infection, noncancerous growth, and cancer; the route must be planned with any later surgery in mind.

\medterm{Bone Marrow} Soft tissue inside many bones. Red marrow makes red cells, white cells, and platelets; yellow marrow stores fat and can produce more blood cells when the body needs them.

\medterm{Bone Marrow Aspiration} Removal of a small amount of liquid marrow, usually from the back of the hipbone, with a needle. It is examined to assess blood-cell development, abnormal cells, infection, or genetic changes.

\medterm{Bone Marrow Aspiration And Biopsy} Two ways of sampling marrow, usually from the back of the hipbone. Aspiration collects liquid marrow for cell and genetic testing; biopsy removes a small solid core to assess marrow structure, cell production, scarring, infection, or cancer.

\medterm{Bone Marrow Biopsy} A procedure that removes marrow, usually from the back of the hipbone, for examination. It may include liquid aspiration and a solid core to assess blood-cell production, abnormal cells, scarring, infection, or cancer.

\medterm{Bone Marrow Culture} A laboratory test in which marrow is placed in special conditions to grow and identify bacteria, fungi, or other organisms when an infection of the marrow or bloodstream is suspected.

\medterm{Bone Marrow Depression} Reduced production of blood cells by the marrow. It can cause anaemia, increased infection risk from too few white cells, bleeding from too few platelets, or a combination of these problems.

\medterm{Bone Marrow Disorder} Any disease that affects the marrow's production of red cells, white cells, or platelets. It may cause anaemia, frequent infections, bleeding, abnormal blood counts, or symptoms from abnormal marrow cells.

\medterm{Bone Marrow Donation} The harvesting of hematopoietic stem and progenitor cells (HSPCs) from an allogeneic or autologous donor via surgical iliac crest bone marrow aspiration or apheresis of granulocyte colony-stimulating factor (G-CSF)-mobilized peripheral blood stem cells for hematopoietic stem cell transplantation.

\synonyms
Hematopoietic Stem Cell Donation; Stem Cell Donation; Bone Marrow Harvest

\medterm{Bone Marrow Failure} A condition in which the marrow cannot make enough red cells, white cells, or platelets. Causes include inherited disorders, aplastic anaemia, medicines, toxins, infections, marrow cancers, and other diseases.

\synonyms
Marrow Failure
Bone Marrow Failure Syndrome

\medterm{Bone Marrow Hyperplasia} An increase in the number of marrow cells because blood-cell production has accelerated. It may be a response to ongoing red-cell destruction, blood loss, low oxygen, or another increased demand for blood cells.

\medterm{Bone Marrow Necrosis} Death of the blood-forming tissue inside bones. It can occur with severe sickle-cell disease, blood cancers, or infection and may cause severe bone pain, anaemia, low white-cell or platelet counts, and fever.

\medterm{Bone Marrow Transplantation} A treatment in which healthy blood-forming stem cells are infused to replace or restore bone-marrow cells damaged by disease or treatment. The cells may come from the patient or a donor. It is used for selected blood cancers, marrow failure, and some inherited disorders; risks include infection, graft failure, organ toxicity, and graft-versus-host disease after a donor transplant.

\synonyms
BMT; Bone Marrow Transplant

\medterm{Bone Mass} The amount of bone tissue in the skeleton or in a particular area. Greater bone mass generally provides more strength, while low bone mass can increase the risk of osteoporosis and fractures.

\synonyms
Bone Mineral Content; Skeletal Bone Mass

\medterm{Bone Matrix} The material surrounding bone cells that gives bone its shape and strength. It contains a flexible framework, mainly type I collagen, reinforced with mineral crystals made largely from calcium and phosphate.

\medterm{Bone Metastases} Areas of cancer that have spread to bone from another organ. They may cause bone pain, weakness or fracture, high blood calcium, or pressure on the spinal cord and nerves.

\medterm{Bone Metastasis} A cancer deposit that has spread to bone from another organ. It may cause pain, weakness, fracture, high blood calcium, or spinal-cord pressure; treatment depends on the original cancer, sites involved, and symptoms.

\medterm{Bone Mineral Density} The amount of mineral, mainly calcium, in a measured area of bone. A DXA scan estimates bone strength, helps diagnose osteoporosis, assesses fracture risk, and monitors treatment together with age and other risk factors.

\synonyms
BMD; Bone Density; Bone Mineral Content

\medterm{Bone Mineral Density Test} See \textit{Bone Density Screening}.

\medterm{Bone Necrosis} Death of bone tissue, usually because its blood supply is reduced or because of severe injury, infection, radiation, or prolonged corticosteroid use. It can cause pain, bone collapse, joint damage, and fracture.

\medterm{Bone Neoplasm} An abnormal growth of cells arising in bone. It may be noncancerous or cancerous; imaging and often a biopsy determine its type, behaviour, extent, and appropriate treatment.

\medterm{Bone Pain} Aching or other discomfort arising from bone. Causes include injury, fracture, infection, reduced blood supply, metabolic disease, overuse, or cancer. Persistent, severe, or unexplained pain requires medical assessment.

\medterm{Bone Pain Or Tenderness} Pain when a bone is used, pressed, or touched. It may result from injury, fracture, infection, inflammation, metabolic disease, or cancer; persistent or unexplained tenderness needs medical assessment.

\medterm{Bone Plate} A metal or other body-compatible plate fixed to bone with screws. It holds broken pieces in position, supports healing, or helps correct a deformity during orthopaedic surgery.

\medterm{Bone Regeneration} Formation of new bone to repair or replace bone that is damaged, missing, or diseased. It occurs naturally during fracture healing and may be supported by grafts, implants, or other treatments.

\medterm{Bone Remodeling} The lifelong replacement of old bone with new bone. Cells remove worn bone and other cells build replacement bone, helping maintain strength, shape, and the body's calcium balance.

\medterm{Bone Resorption} The removal and breakdown of bone tissue, which releases minerals into the blood. It is normal during bone renewal, but excessive resorption can weaken bones and contribute to osteoporosis.

\medterm{Bone Sarcoma} Alternate name; see \textit{Bone cancer}.

\medterm{Bone Saw} A surgical instrument used to cut bone during procedures such as amputation, joint replacement, fracture repair, or orthopaedic reconstruction.

\medterm{Bone Scan} A nuclear-medicine test that shows areas of increased or reduced bone activity after a small amount of tracer is given. It can help investigate fractures, infection, cancer spread, and some metabolic bone diseases, but abnormal uptake is not specific.

\medterm{Bone Scan} A nuclear medicine diagnostic imaging technique in which technetium-99m-labeled bisphosphonates (e.g., $^{99\text{m}}\text{Tc}$-MDP) are administered intravenously and imaged with a gamma camera to identify areas of osteoblastic activity, osteomyelitis, occult fractures, or skeletal metastases.

\synonyms
Bone Scintigraphy; Skeletal Scintigraphy; Radionuclide Bone Scan

\medterm{Bone Screws} Screws used during orthopaedic surgery to hold broken bone pieces together or attach plates, joint components, or other implants. Their size, material, and placement depend on the bone and the operation.

\medterm{Bone Spurs} Extra projections of bone that usually form near joints after long-term wear, arthritis, injury, or abnormal stress. They may cause no symptoms or may irritate nearby tissues and restrict movement.

\synonyms
Osteophytes

\medterm{Bone Transplantation} Transfer of bone tissue or a bone substitute to repair a defect, support fracture healing, or assist reconstruction of the spine, jaw, or a joint. The graft may come from the person, a donor, or a manufactured material.

\medterm{Bone Tumor} An abnormal mass that starts in bone. It may be noncancerous or cancerous; examination, scans, and often a biopsy determine its type, whether it is growing or spreading, and the best treatment.

\medterm{Bone Turnover Markers} Substances measured in blood or urine that reflect how quickly bone is being built or broken down. They may support assessment of some metabolic bone disorders or treatment response but do not replace bone-density testing.

\medterm{Bone Type} A way of grouping bones by shape. Long bones support movement, short bones provide stability, flat bones protect organs or provide broad muscle attachment, irregular bones have complex shapes, and sesamoid bones lie within tendons.

\medterm{Bone X-Ray} An X-ray picture of a bone or group of bones used to look for fractures, abnormal shape, arthritis, infection, bone loss, or some tumours.

\medterm{Bone-Marrow Transplantation} A treatment that replaces damaged or diseased blood-forming marrow with healthy stem cells from the patient or a donor. It is used for selected blood cancers, marrow failure, and some inherited disorders.

\medterm{Bones And Joints} Bones provide support and protection, while joints connect bones and allow movement. Joints may be fixed, slightly movable, or freely movable and are held together by tissues such as cartilage, ligaments, and joint fluid.

\medterm{Bones of the Skeleton} The body’s framework of bones and cartilage, providing support, protection, leverage for movement, mineral storage, and marrow-based blood-cell production.
\medterm{Bongkrek Poisoning} Severe poisoning caused by bongkrekic acid in contaminated fermented foods, especially coconut or corn products. It can cause vomiting, abdominal pain, weakness, liver failure, neurologic injury, and death; urgent treatment is required.

\medterm{Bonnevie-Ullrich Syndrome} An older name for the physical features sometimes associated with Turner syndrome, a condition involving an absent or altered X chromosome. Features may include short stature, a webbed neck, and reduced ovarian function.

\medterm{Bony Exostosis} A localized extra growth of bone. It may occur on the jaw or palate and is often harmless, but removal may be considered if it repeatedly causes injury, interferes with speech or dentures, or needs diagnosis.

\medterm{Bony Labyrinth} A system of channels inside the dense part of the skull bone that houses the inner ear. It contains the cochlea for hearing and the vestibule and semicircular canals for balance.

\medterm{Borborygmi} Rumbling or gurgling sounds made when gas and liquid move through the intestines. They are often normal, but major changes together with pain, vomiting, swelling, or inability to pass stool may signal bowel disease.

\medterm{Borborygmus} A single rumbling or gurgling intestinal sound caused by movement of gas and liquid through the bowel. Its significance depends on symptoms and the overall abdominal examination.

\medterm{Border Definition} The sharpness and shape of the edge of a skin lesion or abnormal growth. A poorly defined or irregular edge can be concerning, but the border alone cannot distinguish harmless from cancerous disease.

\medterm{Borderline} Describes a result or condition near the boundary between normal and abnormal or between diagnostic categories. It usually requires interpretation with symptoms, repeat testing, or other findings; the meaning depends on the clinical context.

\medterm{Borderline Personality Disorder} A long-term mental-health condition involving unstable emotions, self-image, relationships, and behavior. It may include intense fear of abandonment, impulsive actions, self-harm, rapidly changing moods, and difficulty managing feelings.

\synonyms
BPD; Emotionally Unstable Personality Disorder (EUPD)

\medterm{Bordetella} A group of small bacteria that can infect the airways of humans or animals. Bordetella pertussis causes whooping cough; other species usually affect animals but can rarely cause opportunistic human infection.

\medterm{Bordetella Avium} A bacterium that mainly causes respiratory disease in turkeys and other birds. Human infection is very rare and is usually reported in people with weakened immunity or serious underlying illness.

\medterm{Bordetella Bronchiseptica} A bacterium that commonly causes respiratory infection in dogs and other animals. It rarely infects humans, but may cause pneumonia or bloodstream infection, especially in people with weakened immunity.

\medterm{Bordetella Hinzii} A bacterium found mainly in poultry. It rarely causes respiratory or bloodstream infection in humans, usually in people with weakened immunity or significant underlying disease.

\medterm{Bordetella Parapertussis} A bacterium that can cause a whooping-cough-like respiratory infection. Illness is often milder than infection with Bordetella pertussis, but prolonged coughing can still occur and the infection can spread to others.

\medterm{Bordetella Pertussis} The bacterium that causes whooping cough, a highly contagious respiratory infection. It attaches to airway lining, produces toxins that disrupt airway function, and can cause repeated coughing, vomiting, breathing pauses, and serious illness in infants.

\medterm{Bordetella Petrii} An environmental bacterium that rarely causes human disease. When it is found in a specimen, clinicians decide whether it represents true infection by considering the body site, symptoms, and other test results.

\medterm{Boric Acid} A chemical used in some industrial products, insecticides, and limited topical preparations. Swallowing it or absorbing a large amount can poison the body; significant exposure requires prompt medical assessment.

\medterm{Boric Acid Poisoning} Illness caused by swallowing or absorbing too much boric acid. It may cause vomiting, diarrhoea, abdominal pain, skin changes, kidney injury, seizures, shock, or death and requires urgent poison or emergency care.

\medterm{Borna Disease} An uncommon viral infection of the nervous system caused by Borna disease virus. It is mainly recognized in animals; human infection is rare and may cause severe inflammation of the brain.

\medterm{Bornholm Disease} A contagious viral illness causing sudden attacks of severe chest or upper-abdominal muscle pain, often with fever, headache, or vomiting. It is usually caused by an enterovirus and generally resolves with supportive care.

\medterm{Boron} A chemical element found naturally in food, water, soil, and industrial materials. Small dietary amounts are generally tolerated, but excessive exposure can cause digestive symptoms, skin problems, reproductive effects, or developmental toxicity.

\medterm{Borrelia} A group of spiral-shaped bacteria spread mainly by ticks or lice. Different species cause Lyme disease, relapsing fever, or other infections, with symptoms and treatment depending on the species and exposure.

\medterm{Borrelia Burgdorferi} A spiral-shaped bacterium spread by infected blacklegged ticks and the main cause of Lyme disease in North America and parts of Europe. Infection may affect the skin, joints, nervous system, or heart.

\medterm{Borrelia Infection} Infection with a Borrelia bacterium. Depending on the species, it may cause Lyme disease or relapsing fever, with symptoms such as fever, headache, fatigue, rash, muscle pain, or joint and nerve problems.

\medterm{Borrelia Miyamotoi} A bacterium spread by some hard-bodied ticks that causes hard-tick relapsing fever. Illness commonly includes repeated episodes of fever, headache, fatigue, chills, and muscle or joint aches.

\medterm{Borreliosis} Infection caused by a Borrelia bacterium. Human forms include Lyme disease and relapsing fever; symptoms differ according to the species, stage of infection, route of spread, and organs involved.

\medterm{Bortezomib} A cancer medicine that blocks proteasomes, cell structures that remove unwanted proteins. This can kill susceptible cancer cells and is used mainly for multiple myeloma and some lymphomas; nerve damage, low blood counts, digestive symptoms, and infections can occur.

\medterm{Bosentan} A medicine that blocks endothelin, a substance that narrows blood vessels. It is used for pulmonary arterial hypertension, but can injure the liver and harm a developing fetus, so liver tests and pregnancy prevention are important.

\medterm{Bosniak Classification} A system for describing kidney cysts on contrast CT or MRI according to their walls, septations, mineral deposits, and solid parts. Categories help estimate cancer risk and guide monitoring, further imaging, or surgery.

\medterm{Bot Fly} A fly whose larvae can live in animal or human tissue, causing myiasis. In people, a larva may produce a painful boil-like skin lump with a small breathing opening; removal is performed by a clinician to avoid complications.

\medterm{Botfly} A fly whose larvae can infest skin or wounds in animals and occasionally people. The infestation is called myiasis and may cause a painful lump, discharge, itching, or tissue irritation.

\medterm{Bothriocephalus} An older name for a group of fish tapeworms. Related tapeworm infections may cause abdominal symptoms and, when prolonged, can reduce vitamin B12 absorption and lead to anaemia.

\medterm{Botryomycosis} A chronic bacterial infection that forms firm granules within lesions and can resemble a fungal infection. It most often affects the skin and tissue beneath it but may involve internal organs, especially in people with weakened immunity.

\medterm{Bottle Feeding} Feeding an infant expressed breast milk or infant formula from a bottle. Safe preparation, clean equipment, correct formula mixing, and responsive feeding help prevent infection, choking, overfeeding, and poor nutrition.

\medterm{Botulinum Toxin} A substance made by Clostridium botulinum that temporarily blocks nerve signals to muscles and some glands. Carefully measured injections treat selected muscle spasms, movement disorders, migraine, sweating, and cosmetic concerns; too much can cause weakness or breathing failure.

\medterm{Botulinum Toxin Injection, Laryngeal} Targeted electromyography- or laryngoscopy-guided injection of botulinum neurotoxin into specific intrinsic laryngeal muscles (such as the thyroarytenoid or posterior cricoarytenoid) to reduce involuntary spasms in focal laryngeal dystonias, predominantly adductor or abductor spasmodic dysphonia.

\synonyms
Laryngeal Botulinum Toxin Injection; Laryngeal Botox; Spasmodic Dysphonia Injection

\medterm{Botulism} A rare, life-threatening illness caused by botulinum toxin, which blocks nerve signals. It causes blurred or double vision, drooping eyelids, difficulty swallowing, progressive weakness, and sometimes breathing failure; urgent antitoxin and intensive care may be needed.

\medterm{Bouchard Nodes} Firm enlargements of the middle joints of the fingers, usually caused by osteoarthritis. They may be painless or may cause stiffness, pain, reduced movement, and finger deformity.

\medterm{Bouchard's Node} A hard enlargement of one middle finger joint, usually caused by osteoarthritis. It results from changes in the joint and surrounding bone and may cause pain, stiffness, or reduced movement.

\medterm{Bouchard's Nodes} Hard enlargements of the middle finger joints caused most often by osteoarthritis. They can limit movement or cause pain and stiffness, although some people have no symptoms.

\synonyms
Proximal Interphalangeal Osteophytes; PIP-Joint Nodes

\medterm{Bougie} A slender medical instrument used to examine or gradually widen a narrowed passage, such as the food pipe or urethra. It must be used carefully because it can cause pain, bleeding, perforation, or infection.

\medterm{Bougienage} Gradual widening or examination of a narrowed body passage with a bougie. It may relieve swallowing or urinary obstruction, but the cause of narrowing must be assessed and the procedure can cause bleeding or perforation.

\medterm{Bouin's Solution} A laboratory liquid containing picric acid, formaldehyde, and acetic acid that preserves tissue for microscopic examination. It is hazardous to handle and can interfere with some later genetic tests.

\medterm{Bounding Pulse} A hyperdynamic, high-amplitude arterial pulse characterized by a rapid, forceful upstroke and wide pulse pressure (Corrigan or water-hammer pulse), classically associated with aortic regurgitation, thyrotoxicosis, severe anemia, fever, patent ductus arteriosus, or septic shock.

\synonyms
Hyperdynamic Pulse; Water-Hammer Pulse; Corrigan Pulse

\medterm{Bourbon Virus} A rare tick-associated viral infection that can cause fever, tiredness, poor appetite, nausea, and low white-cell or platelet counts. No specific treatment or vaccine is established; care is supportive and other serious infections must be excluded.

\medterm{Boutonneuse Fever} A tick-borne infection caused usually by Rickettsia conorii. It commonly causes fever, headache, rash, and a dark sore at the bite site; prompt antibiotic treatment reduces the risk of serious illness.

\medterm{Boutonniere Deformity} A finger shape problem where the middle joint stays bent inward and the tip bends backward, often caused by an injury or rheumatoid arthritis.

\medterm{Bouveret Syndrome} A rare blockage near the stomach outlet caused by a gallstone that has entered the intestine through an abnormal connection from the gallbladder. It can cause persistent vomiting, upper-abdominal pain, bloating, and inability to eat.

\medterm{Bovine} Relating to cattle or to material obtained from cattle. In medicine it may describe an animal disease, a biological product, tissue used in treatment, or an infection that can pass between cattle and people.

\medterm{Bovine Anaplasmosis} A tick- or blood-spread infection of cattle caused by Anaplasma marginale. The bacterium infects red blood cells and destroys them, causing anaemia, weakness, fever, reduced milk production, or death.

\medterm{Bovine Papillomavirus} A group of papillomaviruses that infect cattle and can cause skin warts or other growths. Different strains vary in their effects; human disease is not the usual concern.

\medterm{Bovine Spongiform Encephalopathy} A fatal disease of cattle caused by misfolded proteins called prions. It progressively damages the brain and nervous system; people can rarely develop variant Creutzfeldt--Jakob disease after eating contaminated cattle tissue.

\medterm{Bow Legs} An angular deformational alignment of the lower extremity characterized by outward lateral bowing of the tibia and femur such that the knees remain abnormally separated when the medial malleoli of the ankles are approximated.

\synonyms
Genu Varum; Bowlegs; Bandy Legs

\medterm{Bowel} The small and large intestines. They receive partly digested food, absorb nutrients and water, move contents forward, and form and store stool before it leaves the body.

\medterm{Bowel Burn} Injury to the bowel lining caused by heat, chemicals, radiation, reduced blood flow, or another damaging exposure. It can cause pain, bleeding, perforation, infection, or narrowing and requires assessment according to the cause.

\medterm{Bowel Cancer} Cancer arising in the colon or rectum. It may cause a persistent change in bowel habit, blood in stool, anaemia, weight loss, abdominal pain, or no early symptoms; screening and timely assessment improve detection.

\medterm{Bowel Disorders} Pathological conditions affecting the small or large intestine, encompassing inflammatory bowel disease (Crohn disease, ulcerative colitis), infectious enterocolitis, ischemic bowel disease, bowel obstruction, and colorectal neoplasms.

\synonyms
Intestinal Diseases; Bowel Diseases; Intestinal Disorders

\medterm{Bowel Incontinence} Inability to control the passage of stool or liquid from the anus. It can result from diarrhoea, constipation, muscle or nerve injury, childbirth, bowel disease, or neurologic illness and may often be improved with treatment.

\synonyms
Fecal Incontinence

\medterm{Bowel Infarction} Death of part of the intestine because its blood supply is blocked or severely reduced. It can cause severe abdominal pain, bloody stool, vomiting, and abdominal tenderness and is a life-threatening emergency.

\medterm{Bowel Ischemia} Reduced blood flow to the intestine, causing injury from too little oxygen. It may cause sudden or worsening abdominal pain, vomiting, or bloody stool and can progress to bowel death if circulation is not restored.

\medterm{Bowel Obstruction} A blockage that prevents food, fluid, and gas from moving through the intestine. It can cause cramping pain, swelling, vomiting, and inability to pass stool or gas and may require urgent hospital treatment.

\medterm{Bowel Preparation} Emptying and cleansing the colon before a colonoscopy or other bowel procedure, usually with a special diet and a prescribed laxative solution. Following the instructions is important because poor preparation can hide lesions.

\synonyms
Bowel Prep; Colon Cleansing; Colonic Preparation

\medterm{Bowel Program} A planned routine using food, fluids, medicines, movement, and scheduled toilet visits to produce regular bowel movements and manage constipation or loss of bowel control.

\medterm{Bowel Retraining} A structured routine that uses regular toilet times, suitable food and fluids, physical activity, and helpful posture to establish predictable bowel movements and improve constipation or bowel-control problems.

\medterm{Bowel Sounds} Gurgling or clicking sounds made as the intestines move gas and liquid. Their presence and pattern can provide clues during an abdominal examination but cannot diagnose a disorder by themselves.

\medterm{Bowel Transit Time} The time food and other contents take to travel through the digestive tract. Testing it can help investigate unusually slow or fast movement when a bowel-movement or digestive disorder is suspected.

\medterm{Bowel Urgency} A sudden, hard-to-delay need to pass stool. It can occur with diarrhoea, infection, irritable bowel syndrome, inflammatory bowel disease, or disorders of the rectum and may lead to accidents.

\medterm{Bowen's Disease} An early form of skin cancer confined to the outermost skin layer. It usually appears as a slowly enlarging, persistent scaly red or brown patch and can invade deeper tissue if untreated.

\medterm{Bowenoid Papulosis} A condition linked to high-risk human papillomavirus that causes multiple reddish-brown or purple bumps on genital or nearby skin. The cells look cancerous under a microscope, but the lesions often behave less aggressively; treatment and follow-up reduce the risk of progression.

\medterm{Bowlegs} Alternate name; see \textit{Bow Legs}.

\medterm{Bowler's Thumb} Irritation or scarring of the digital nerve along the thumb caused by repeated pressure from a bowling ball. It may produce pain, numbness, tingling, or a tender lump and may improve with rest or equipment changes.

\medterm{Bowman Capsule} A cup-shaped structure at the beginning of each kidney filtering unit. It surrounds a knot of tiny blood vessels and collects the fluid filtered from blood before it passes through the rest of the kidney tubule.

\medterm{Bowman Membrane} A thin, strong layer near the front of the cornea that supports and protects the eye's clear surface. Significant injury can scar it because it does not normally grow back.

\medterm{Bowman's Capsule} The cup-shaped beginning of a kidney filtering unit. It surrounds a cluster of small blood vessels and receives the fluid filtered from blood.

\medterm{Bowman's Gland} A small gland high inside the nose that releases watery fluid onto the smell-sensing lining. The fluid helps dissolve scent molecules and clear the surface so odours can be detected.

\medterm{Boxed Warning} The strongest safety warning required in a medicine's prescribing information when there is a serious or life-threatening risk. It explains a major hazard but does not necessarily mean the medicine is unsuitable for every patient.

\medterm{Boxer Fracture} A transverse or comminuted fracture through the distal metaphysis or neck of the fifth metacarpal bone (and occasionally the fourth metacarpal), characteristically accompanied by volar angulation of the metacarpal head resulting from direct axial loading with a clenched fist.

\medterm{Boxer's Ear} A deformity of the outer ear caused by repeated injury or untreated blood collection around the ear cartilage. It can produce thickened, scarred tissue with a cauliflower-like appearance.

\medterm{Boxer's Fracture} A break near the knuckle end of the fifth metacarpal, the hand bone below the little finger. It commonly follows punching and causes pain, swelling, tenderness, or a changed knuckle shape.

\medterm{Braces and Retainers} Alternate name; see \textit{Dental Braces}.

\medterm{Brachial} Relating to the upper arm, including its bones, muscles, blood vessels, and nerves. The brachial artery supplies the arm, and the brachial plexus carries nerve signals to and from the arm and hand.

\medterm{Brachial Artery} The main artery of the upper arm. It carries blood from the armpit toward the elbow, where it divides into the radial and ulnar arteries, and it is commonly used to measure a pulse or blood pressure.

\medterm{Brachial Plexopathy} Damage or dysfunction of the nerve network supplying the shoulder, arm, and hand. It can cause pain, weakness, numbness, or paralysis after injury, pressure, inflammation, a tumour, or radiation.

\medterm{Brachial Plexus} A network of nerves from the lower neck and upper chest that supplies movement and sensation to the shoulder, arm, forearm, and hand. Its branches include the musculocutaneous, axillary, radial, median, and ulnar nerves.

\medterm{Brachial Plexus Cords} Three parts of the nerve network supplying the arm, named lateral, posterior, and medial because of their position around an artery in the armpit. They join to form nerves controlling movement and feeling in the upper limb.

\medterm{Brachial Plexus Injuries} Damage to the nerve network supplying the shoulder and upper limb, caused by stretching, pressure, or a penetrating injury. Symptoms range from pain or numbness to weakness or paralysis, depending on the nerves affected.

\medterm{Brachial Plexus Injury} Damage to the nerve network of the shoulder and arm, sometimes occurring during a difficult birth or after trauma. It may cause weakness, loss of movement, numbness, or paralysis; recovery depends on the severity and nerves involved.

\medterm{Brachial Plexus Injury In Newborns} Damage to the nerves supplying a newborn's shoulder, arm, and hand, usually after a difficult delivery. The affected arm may be weak or still, and recovery ranges from spontaneous improvement to requiring specialist treatment.

\medterm{Brachial Plexus Neuropathies} Disorders affecting the nerve network that supplies the shoulder, arm, and hand. They can cause pain, weakness, numbness, or loss of movement because of injury, pressure, inflammation, tumours, or radiation.

\medterm{Brachial Plexus Roots} The first portions of the nerve network that supplies the arm, arising from the lower neck and upper chest. Damage to different roots produces characteristic patterns of weakness, numbness, and pain.

\medterm{Brachial Plexus Trunks} Three sections of the brachial plexus—upper, middle, and lower—that carry nerve fibres toward the arm. Injury to the upper trunk can cause weakness of the shoulder and elbow, while lower-trunk injury can mainly affect the hand.

\medterm{Brachial Vein} Deep veins accompanying the brachial artery that drain the arm and join the basilic vein to form the axillary vein.

\medterm{Brachialgia} Pain that travels from the neck or shoulder down the arm. It often results from irritation of a neck nerve or the brachial plexus and may occur with numbness, tingling, or weakness.

\medterm{Brachialis Muscle} A strong muscle at the front of the upper arm, beneath the biceps. It is the main muscle that bends the elbow and works whether the palm faces up, down, or sideways.

\medterm{Brachiocephalic Trunk} The first large branch of the aorta in the chest. It carries blood toward the right side of the head, neck, and arm and divides into the right common carotid and right subclavian arteries.

\medterm{Brachiocephalic Veins} Two large veins formed by the joining of veins from the head, neck, and arm. They drain blood from these areas into the superior vena cava, which returns it to the heart.

\medterm{Brachioplasty} An operation to remove or reshape excess skin and tissue of the upper arm, sometimes called an arm lift. Possible complications include scars, altered sensation, bleeding, fluid collection, infection, and wound problems.

\medterm{Brachioradialis Reflex} A deep tendon stretch reflex elicited by briskly tapping the tendon of the brachioradialis muscle near the styloid process of the radius with a reflex hammer, producing forearm flexion and slight pronation. It is innervated by the radial nerve and tests the integrity of the C5 and C6 cervical spinal cord segments and nerve roots.

\synonyms
Supinator Reflex; Radial Reflex

\medterm{Brachium} The part of the upper limb between the shoulder and elbow. It contains the upper-arm bone, muscles that move the shoulder and elbow, and major nerves and blood vessels.

\medterm{Brachycephalic} Describes a skull that is relatively short from front to back and broad from side to side. It may be a normal family characteristic or part of a condition affecting skull growth.

\medterm{Brachycephaly} A broad, short skull shape. It may develop from pressure on an infant's head or from early joining of skull bones; assessment determines whether brain growth or treatment is affected.

\medterm{Brachydactyly} A birth condition in which one or more fingers or toes are shorter than usual because of altered development of their bones or surrounding tissues. It may occur alone or as part of an inherited syndrome.

\medterm{Brachytherapy} Internal radiation treatment in which a radioactive source is placed inside or close to a tumour. It delivers a concentrated dose to a limited area and may use a temporary applicator or a source left in place permanently.

\synonyms
Brachy Therapy
Internal Radiation Therapy

\medterm{Brachytherapy Seed} A tiny radioactive source placed in or near a tumour to deliver local radiation. Seeds may be temporary or permanent and are used most often for selected cancers, including prostate cancer.

\medterm{Bradycardia} A slower-than-usual heart rate, often below 60 beats per minute in adults. It can be normal during sleep or in trained athletes, but may result from heart-conduction disease or medicines and cause dizziness, fainting, tiredness, or breathlessness.

\synonyms
Slow Heartbeat

\medterm{Bradycardia Of Prematurity} A slow heart rate in a premature infant, often occurring with pauses in breathing and low oxygen. It commonly reflects immature breathing and heart control and is monitored and treated according to its frequency and severity.

\medterm{Bradycardia-Tachycardia Syndrome} Alternate name; see \textit{Sick Sinus Syndrome}.

\medterm{Bradykinesia} Slowness and reduced size of voluntary movements. It can make starting or repeating actions difficult and may occur with stiffness, tremor, reduced arm swing, a less expressive face, or a shuffling walk in Parkinson disease and related disorders.

\medterm{Bradykinesia and Athetosis} Hypokinetic and dyskinetic movement patterns characterized by generalized slowness of movement initiation (bradykinesia) or continuous slow, writhing, twisting involuntary movements (athetosis), frequently observed in Parkinson disease, secondary parkinsonism, or basal ganglia encephalopathies.

\synonyms
Slowness of Movement; Hypokinesia and Athetosis; Extrapyramidal Slowness
\medterm{Bradykinetic} Describes abnormally slow movement, especially the reduced speed and size of movement seen in Parkinson disease and related neurologic disorders.

\medterm{Bradykinin} A natural chemical messenger that widens some blood vessels and allows more fluid to leave them. It contributes to inflammation, swelling, pain, and rare reactions to certain blood-pressure medicines.

\medterm{Bradyphrenia} Unusually slow thinking or mental processing. It may cause delayed responses and difficulty completing tasks and can occur with neurologic disease, mental illness, medicines, depression, or other conditions.

\medterm{Bradypnea} Abnormally slow breathing, often fewer than 12 breaths per minute in an adult. It may result from sedative or opioid effects, brain disease, low body temperature, hormone disorders, or severe illness and can lead to respiratory failure.

\medterm{Bradypnoea} Abnormally slow breathing. Causes include medicines, neurologic disease, metabolic disturbance, low body temperature, or severe illness; breathing that is difficult, shallow, or associated with confusion requires urgent care.

\medterm{Braille System} A raised-dot writing and reading system used by people who are blind or have low vision. Patterns represent letters, numbers, punctuation, and symbols and can be read by touch.

\medterm{Brain} The main organ of the nervous system, protected inside the skull. It processes senses, controls movement and body functions, and supports thinking, memory, emotions, language, and consciousness.

\medterm{Brain Abscess} A pocket of pus and infection inside the brain. It can develop when infection spreads from the blood, nearby ear or sinus disease, surgery, or an injury and may cause headache, fever, seizures, or weakness.

\medterm{Brain Aneurysm} A weakened bulge in an artery supplying the brain. If it bursts, it can cause sudden severe headache, vomiting, neck stiffness, collapse, or other neurological problems.

\medterm{Brain Aneurysm Repair} Treatment to prevent a brain aneurysm from bleeding or bleeding again. A surgeon may place a clip or a specialist may seal the aneurysm from inside the artery with coils or another device.

\medterm{Brain Arteriovenous Malformation} A congenital vascular dysplasia consisting of a direct high-flow, low-resistance communication between cerebral arteries and veins through an abnormal vascular nidus without an intervening capillary bed, carrying a risk of intracranial hemorrhage, focal seizures, and neurological deficits.

\synonyms
Cerebral Arteriovenous Malformation; Brain AVM; Intracranial AVM

\medterm{Brain AVM} Alternate name; see \textit{Brain Arteriovenous Malformation}.

\medterm{Brain Cancer} Cancer affecting brain tissue or its coverings. Symptoms depend on the location and may include seizures, headache, weakness, speech or vision changes, personality change, or pressure-related vomiting; diagnosis usually requires brain imaging and sometimes tissue testing.

\synonyms
Malignant Brain Tumor; Primary Brain Neoplasm

\medterm{Brain Concussion} Alternate name; see \textit{Concussion}.

\medterm{Brain Death} Permanent loss of all brain and brainstem function, including consciousness, reflexes, and the ability to breathe independently. When formally confirmed according to law and medical standards, it means the person has died.

\medterm{Brain Disorder} Any illness or abnormality affecting the brain's structure or function, including stroke, infection, tumour, epilepsy, injury, degenerative disease, or a condition present from birth.

\medterm{Brain Edema} Swelling caused by extra fluid in brain tissue. It can increase pressure inside the skull after injury, stroke, infection, bleeding, tumour, or severe illness and may cause headache, vomiting, seizures, confusion, or reduced consciousness.

\medterm{Brain Embolism} Alternate name; see \textit{Embolic Stroke}.

\medterm{Brain Fog} An informal description of difficulty concentrating, thinking clearly, remembering, or finding words. It is a symptom rather than a diagnosis and may occur with poor sleep, stress, illness, medicines, depression, or other conditions.

\medterm{Brain Fornix} A curved bundle of nerve fibres deep in the brain that helps connect the hippocampus with other areas involved in memory. Damage may cause difficulty learning or remembering information.

\medterm{Brain Haemorrhage} Bleeding within or around the brain. It may result from high blood pressure, injury, an aneurysm or other blood-vessel problem, clotting disorders, or medicines and can cause sudden neurologic symptoms.

\medterm{Brain Health} The ability of the brain to support thinking, memory, learning, movement, emotions, communication, and consciousness. It is supported by controlling vascular risks, sleeping well, physical activity, balanced nutrition, social connection, and treatment of illness.

\medterm{Brain Hemorrhage} Bleeding within or around the brain, often caused by high blood pressure, an aneurysm, a blood-vessel malformation, injury, or blood-thinning treatment. Sudden severe headache, weakness, speech trouble, vomiting, seizure, or reduced consciousness requires emergency care.

\medterm{Brain Herniation} Dangerous displacement of brain tissue caused by severe swelling, bleeding, a tumour, or another increase in pressure inside the skull. It can compress the brainstem, impair breathing and consciousness, and requires emergency treatment.

\medterm{Brain Imaging} Tests that create pictures of brain structure or activity. CT and MRI show tissue, bleeding, swelling, tumours, and strokes; blood-vessel or nuclear scans answer more specific questions.

\synonyms
Neuroimaging; Cerebral Imaging

\medterm{Brain Injuries} Damage to the brain ranging from a concussion to severe injury, commonly caused by falls, crashes, sports impacts, or penetrating wounds. Effects depend on the area and severity and may include headache, confusion, memory loss, seizures, or unconsciousness.

\medterm{Brain Injury} Damage or impaired function of the brain caused by trauma, lack of oxygen, stroke, infection, toxins, or another disease. Effects vary with the location and severity and may affect movement, sensation, thinking, behaviour, or consciousness.

\medterm{Brain Injury} Alternate name; see \textit{Traumatic Brain Injury}.

\medterm{Brain Lesions} Areas of brain tissue that look or function abnormally because of a tumour, stroke, infection, inflammation, injury, loss of nerve insulation, or a blood-vessel disorder. Symptoms depend mainly on the lesion's location and size.

\medterm{Brain Lesions} Alternate name; see \textit{Brain Lesion}.

\medterm{Brain Malformations} Structural differences in the brain that develop before birth. They may result from genetic changes, disrupted fetal development, infection, or other prenatal influences and can affect movement, learning, seizures, or other neurologic functions.

\medterm{Brain Mapping} Methods used to identify where brain functions and structures are located. Imaging, electrical recordings, or direct stimulation may guide research or help surgeons avoid important areas during selected operations.

\medterm{Brain Metastases} Cancer deposits that have spread to the brain from another organ. They may cause headache, seizures, weakness, speech or vision changes, confusion, or no symptoms and are assessed with brain imaging.

\synonyms
Secondary Brain Cancer

\medterm{Brain Metastasis} A cancer deposit that has spread to the brain from another organ. Symptoms may include headache, seizures, weakness, speech or vision changes; treatment depends on the number, location, original cancer, and overall health.

\medterm{Brain Monitoring} Recording brain activity or pressure to assess neurologic function, seizures, sleep, or the depth of anaesthesia. The method used depends on the clinical question; a monitor result is interpreted with examination findings and other tests.

\medterm{Brain Natriuretic Peptide Test} A blood test measuring BNP or NT-proBNP, substances released when the heart is stretched. A raised result can support a diagnosis of heart failure, but it must be interpreted with symptoms, examination, age, kidney function, and other findings.

\medterm{Brain Neoplasm} An abnormal growth of cells in the brain or its coverings. It may be noncancerous or cancerous; imaging and sometimes biopsy determine its type, behaviour, extent, and treatment options.

\medterm{Brain PET Scan} A scan that uses a small amount of radioactive tracer to show chemical activity in the brain. It can help investigate selected tumours, memory disorders, seizures, or brain function when other scans do not answer the question.

\medterm{Brain Sand} Brain sand is an older name for corpora arenacea, tiny calcified deposits commonly found in the pineal gland and sometimes other brain regions.

\medterm{Brain SPECT} A nuclear-imaging scan that shows blood flow or selected brain activity using a radioactive tracer. It may help evaluate seizures, stroke-related blood-flow changes, dementia, or other specialist questions when interpreted with other findings.

\medterm{Brain Stem} The lower part of the brain connecting the cerebrum with the spinal cord and controlling essential functions such as breathing and consciousness.

\medterm{Brain Stem Glioma} A tumor arising from glial cells in the brainstem, the part of the brain controlling important functions such as breathing, swallowing, eye movement, and balance. Symptoms and treatment depend on its exact location and behavior.

\medterm{Brain Stimulation} Use of electrical, magnetic, or other energy to change brain activity. Methods include electroconvulsive therapy, transcranial magnetic stimulation, and deep-brain stimulation; their uses, benefits, and risks differ.

\medterm{Brain Surgery} An operation on the brain or nearby structures to remove or sample a tumor, relieve pressure, treat bleeding, repair an abnormal blood vessel, or control selected seizures. Risks depend on the area and procedure and can include bleeding, infection, or neurologic loss.

\medterm{Brain Tumor} An abnormal growth of cells in the brain or nearby structures. It may be noncancerous or cancerous, begin in the brain, or spread there from cancer elsewhere. It may cause seizures, headache, vomiting, weakness, changes in thinking, or no symptoms at first.

\medterm{Brain Tumor, Adult Primary} A neoplasm arising de novo within the brain parenchyma, cranial nerves, or meninges in adults, encompassing diffuse gliomas (astrocytoma, oligodendroglioma, glioblastoma), meningiomas, vestibular schwannomas, and primary central nervous system lymphomas.

\synonyms
Adult Primary Brain Tumor; Primary Intracranial Neoplasm; Adult Glioma

\medterm{Brain Tumor, Pediatric} A primary neoplasm originating within the parenchyma of the brain or surrounding meninges and cranial nerves during childhood or adolescence. Common histologic types include pilocytic astrocytoma, medulloblastoma, ependymoma, and diffuse intrinsic pontine glioma.

\synonyms
Pediatric Brain Tumor; Childhood Brain Tumor; Pediatric CNS Neoplasm

\medterm{Brain Tumour} An abnormal growth in or near the brain that may be noncancerous or cancerous. It can cause seizures, headache, vomiting, changes in thinking, weakness, or other problems by affecting nearby brain tissue.

\medterm{Brain Ventricles} Four connected spaces inside the brain that contain and circulate cerebrospinal fluid. This fluid cushions the brain and spinal cord, carries substances, and helps remove waste; blockage can cause pressure and hydrocephalus.

\medterm{Brain Wave} A pattern of electrical activity made by groups of nerve cells and recorded with an electroencephalogram. Patterns change with sleep, alertness, seizures, and some disorders of brain function.

\medterm{Brain Waves} Patterns of electrical activity produced by groups of nerve cells and recorded with an electroencephalogram. Their timing and shape can help assess sleep, seizures, brain function, and some disorders affecting consciousness.

\synonyms
Cerebral Electrical Activity; EEG Rhythms

\medterm{Brainstem} The lower part of the brain connecting the cerebrum with the spinal cord. It helps control breathing, heart rate, blood pressure, swallowing, eye movements, alertness, and pathways carrying signals between the brain and body.

\medterm{Brainstem Auditory Evoked Response} An electrodiagnostic test that records electrical potentials generated by the auditory nerve and brainstem pathways in response to acoustic click or tone stimuli. It is utilized to assess hearing acuity in newborns and uncooperative patients, and to evaluate auditory pathway lesions such as acoustic neuromas.

\synonyms
Brainstem Auditory Evoked Potential; BAEP; Brainstem Auditory Evoked Response; BAER; Auditory Brainstem Response; ABR

\medterm{Brainstem Tumor} A growth arising in the brainstem, which connects the brain with the spinal cord and controls important functions. Symptoms may include problems with eye movement, balance, swallowing, speech, strength, or breathing.

\medterm{Branchial} Relating to structures that develop from the paired tissue arches in an embryo's early head and neck. These arches contribute to parts of the face, throat, voice box, muscles, bones, blood vessels, and nerves.

\medterm{Branchial Arch} One of several paired structures in an embryo that develop into parts of the face, throat, voice box, muscles, bones, blood vessels, and nerves of the head and neck.

\medterm{Branchial Cleft} A groove between developing head-and-neck arches in an embryo. If part of it remains after birth, it may form a cyst, a narrow passage, or an opening, usually on the side of the neck.

\medterm{Branchial Cleft Anomaly} A birth-related remnant of developing head-and-neck tissue that may form a cyst, narrow passage, or skin opening. Repeated infection, swelling, or drainage from the side of the neck needs medical assessment.

\medterm{Branchial Cleft Cyst} A fluid-filled lump on the side of the neck caused by tissue left over from development before birth. It may enlarge or become painful after infection and occasionally causes drainage, swallowing difficulty, or pressure on the airway.

\medterm{Branchial Cleft Cyst} Alternate name; see \textit{Branchial Cyst}.

\medterm{Branchial Cyst} A benign, congenital lateral neck cyst that arises from the failure of obliteration of embryonic branchial cleft tissue, most commonly the second branchial cleft (accounting for over 90\% to 95\% of branchial anomalies). It typically presents in adolescents or young adults as a smooth, painless, fluctuant, non-transilluminating mass located along the anterior border of the upper third of the sternocleidomastoid muscle near the angle of the mandible, which often enlarges or becomes tender following an upper respiratory infection. Fine-needle aspiration yields characteristic turbid, straw-colored fluid containing cholesterol crystals. In adults over 40, lateral cystic neck masses require careful evaluation to rule out cystic nodal metastasis from occult oropharyngeal carcinoma. Definitive treatment is complete surgical excision of the cyst and its associated tract.

\synonyms
Branchial Cleft Cyst; Second Branchial Cleft Cyst; Lateral Cervical Cyst

\medterm{Branchio-Oto-Renal Syndrome} An inherited disorder that can cause neck cysts or openings, hearing loss, and kidney abnormalities. The combination and severity vary, and some affected people have only one or two of these features.

\medterm{Branchioma} A rare cyst or growth arising from tissue related to early head-and-neck development. It often presents as a lump in the neck and requires examination and imaging to establish its nature.

\medterm{Brass-Founder's Ague} An older name for metal-fume fever, a short-lived flu-like illness after breathing fumes from heated metals such as zinc or brass. Symptoms may include fever, chills, headache, muscle aches, cough, and nausea.

\medterm{Brassicasterol} A plant sterol found especially in rapeseed and related oils. It is mainly used to identify the composition of plant oils and has no routine role as a medical treatment.

\medterm{Brassidic Acid} A long-chain fatty acid found in some plant oils. It is mainly a chemistry term and is not an established medicine or routine clinical test.

\medterm{Braxton Hicks Contractions} Irregular tightening of the uterus during pregnancy, often becoming more noticeable later in pregnancy. They usually do not steadily strengthen or open the cervix and may ease with rest, fluids, or a change of activity.

\synonyms
False Labor; Practice Contractions

\medterm{BRCA} BRCA1 and BRCA2 are genes that help repair damaged DNA. An inherited harmful change in either gene increases the risk of breast, ovarian, fallopian-tube, primary peritoneal, prostate, and pancreatic cancers.

\medterm{BRCA1 And BRCA2} Genes that help repair breaks in DNA and limit abnormal cell growth. An inherited harmful variant raises the risk of several cancers and may affect screening, preventive surgery, and treatment choices for the person and relatives.

\medterm{BRCA1 And BRCA2 Gene Testing} A blood or saliva test looking for harmful inherited changes in BRCA1 or BRCA2. Results can guide cancer screening, preventive options, and treatment and may provide important information for biological relatives.

\medterm{BRCA1 Gene} A gene involved in repairing damaged DNA and limiting tumor formation. An inherited harmful variant increases the risk of breast, ovarian, fallopian-tube, peritoneal, prostate, and pancreatic cancers; genetic counseling helps interpret test results.

\medterm{BRCA2 Gene} A gene that helps repair serious DNA damage. An inherited harmful variant increases the risk of breast, ovarian, prostate, pancreatic, and some other cancers and may affect relatives; genetic counseling helps explain testing and options.

\medterm{Breakbone Fever} An older name for dengue fever, reflecting the severe muscle and joint pain that may occur with this mosquito-borne viral infection. Fever, headache, rash, nausea, or dangerous bleeding may also develop.

\medterm{Breakthrough Infection} An infection that occurs despite vaccination or another preventive measure intended to reduce its risk. Protection and illness severity vary with the germ, intervention, timing, and person's immune system.

\medterm{Breakthrough Pain} A short flare of pain that occurs despite otherwise controlled persistent pain. It may happen unexpectedly or after activity; treatment depends on the cause, pattern, regular pain medicines, and safety risks.

\medterm{Breast} One of the two organs on the front of the chest, made of milk-producing tissue, ducts, fat, connective tissue, blood vessels, nerves, and skin. It can be affected by benign conditions and cancer.

\medterm{Breast Abscess} A pocket of pus in breast tissue, often developing after mastitis during breastfeeding. It may cause a painful lump, redness, warmth, and fever; ultrasound can confirm it, and drainage plus antibiotics may be needed.

\medterm{Breast Anatomy} The breast contains skin, fat, milk-producing lobules, ducts, connective tissue, blood vessels, nerves, and lymph vessels that extend toward the armpit. Its appearance changes with age, hormones, pregnancy, and breastfeeding.

\medterm{Breast And Thoracic Cancer Surgery} Operations for cancers of the breast or chest, including selected cancers of the lung, pleura, food pipe, mediastinum, or chest wall. Surgery may remove a tumour and nearby lymph nodes, with further treatment based on the cancer type and stage.

\synonyms
Breast and Thoracic Oncologic Surgery

\medterm{Breast Augmentation} A procedure to increase breast size or change its shape, usually with an implant or transferred fat. Risks include infection, bleeding, scarring, altered sensation, implant problems, and the need for further surgery.

\synonyms
Breast Enlargement; Augmentation Mammoplasty

\medterm{Breast Augmentation Surgery} An operation that enlarges or reshapes the breasts using implants or transferred tissue. The surgeon discusses expected results, alternatives, risks, recovery, and possible later procedures before treatment.

\medterm{Breast Biopsy, Stereotactic} A minimally invasive percutaneous core-needle biopsy procedure utilizing digital mammographic triangulation from angled X-ray views to precisely target and sample nonpalpable microcalcifications or architectural distortions.

\synonyms
Stereotactic Breast Biopsy; Stereotactic Core Needle Biopsy; Mammographically Guided Breast Biopsy

\medterm{Breast Biopsy, Ultrasound-Guided} A percutaneous core-needle or vacuum-assisted tissue sampling procedure performed under continuous real-time sonographic guidance to biopsy solid masses, architectural abnormalities, or complex cystic breast lesions.

\synonyms
Ultrasound-Guided Breast Biopsy; Sonographically Guided Core Biopsy; Ultrasound Breast Biopsy

\medterm{Breast Bone} A common name for the breastbone, or sternum, a flat bone in the centre of the chest. It joins the ribs and helps protect the heart and lungs.

\medterm{Breast Calcifications} Small calcium deposits seen in breast tissue on a mammogram. Most are harmless, but their size, shape, and arrangement may require additional images or biopsy to exclude an early cancer.

\medterm{Breast Cancer} Cancer that begins in breast tissue. It may appear as a lump, skin or nipple change, or abnormal imaging, although early cancer may cause no symptoms; treatment depends on its type, stage, and biology.

\medterm{Breast Cancer And Lymphedema} Breast-cancer surgery or radiation can damage lymph drainage and cause lasting swelling, heaviness, tightness, or skin changes in the arm, breast, or chest. Early assessment and specialist treatment can reduce progression and complications.

\medterm{Breast Cancer Genetics} The study of inherited gene changes that raise breast-cancer risk. Important genes include BRCA1, BRCA2, PALB2, TP53, PTEN, and CHEK2; genetic counselling helps decide whether testing and extra screening are appropriate.

\medterm{Breast Cancer Prevention} Ways to lower, but not eliminate, breast-cancer risk include physical activity, limiting alcohol, avoiding tobacco, appropriate screening, and discussing inherited risk. Preventive medicine or surgery may be considered when risk is substantially high.

\medterm{Breast Cancer Recurrence} Return of breast cancer after treatment, either near the original site, in nearby lymph nodes, or in distant organs. New persistent symptoms or a lump require assessment, although many such symptoms have noncancerous causes.

\medterm{Breast Cancer Screening} Testing to find breast cancer before symptoms appear, usually with mammography. Ultrasound or MRI may be added for selected people according to age, breast density, inherited risk, and other factors.

\medterm{Breast Cancer Staging} Assessment of how far breast cancer has grown or spread, including the tumour's features, nearby lymph nodes, and distant organs. The stage helps estimate outlook and choose treatment.

\medterm{Breast Cancer, Inflammatory} Alternate name; see \textit{Inflammatory Breast Cancer}.

\medterm{Breast Cancer, Male} Alternate name; see \textit{Male Breast Cancer}.

\medterm{Breast Carcinoma} A cancer that begins in breast tissue, usually in the ducts or milk-producing lobules. Its treatment depends on the cell type, hormone and other markers, tumour size, lymph nodes, and spread.

\medterm{Breast Cyst} A benign, fluid-filled sac within breast tissue, most common in women between the ages of 35 and 50. It typically presents as a smooth, round, movable lump that can fluctuate in size and tenderness with the menstrual cycle. Ultrasound imaging distinguishes fluid-filled cysts from solid masses. Simple cysts are noncancerous and do not increase the risk of breast cancer; they generally require no intervention unless fluid is drained with a needle to relieve discomfort.

\medterm{Breast Cysts} Alternate name; see \textit{Breast Cyst}.

\medterm{Breast Cytology} Examination of individual cells from a breast lump, nipple discharge, or needle sample. It can suggest infection, a harmless change, or cancer, but a core biopsy is often needed to determine the full tissue pattern.

\medterm{Breast Density} The relative amount of milk-gland and connective tissue compared with fat on a mammogram. Dense breast tissue can hide small cancers and is linked with higher breast-cancer risk; screening plans depend on individual risk and local guidance.

\medterm{Breast Development} Growth and maturation of breast tissue, beginning mainly at puberty and changing with hormones, pregnancy, breastfeeding, ageing, and some medicines. Uneven development is common, but a new persistent change requires assessment.

\medterm{Breast Duct} A channel that carries milk from milk-producing areas toward the nipple. Blockage, inflammation, or abnormal cells in a duct can cause a lump, pain, or nipple discharge.

\medterm{Breast Engorgement} Overfilling and swelling of the breasts with milk and extra tissue fluid, usually a few days after birth. Both breasts may feel firm, warm, painful, and make the nipples harder for an infant to grasp.

\medterm{Breast Enlargement In Males} Enlargement of glandular breast tissue in a male, called gynecomastia. It may result from normal hormonal changes, medicines, obesity, liver or kidney disease, or other hormone disorders; a hard one-sided lump needs assessment.

\medterm{Breast Examination} Inspection and gentle examination of the breasts, nipples, and nearby lymph nodes to assess lumps, pain, skin or nipple changes, and other findings. It may be followed by imaging or biopsy when needed.

\medterm{Breast Feed} To feed an infant directly from the breast with human milk. Human milk provides nutrients and immune protection, but feeding may need support when the parent or infant has illness, difficulty latching, or medicine-related concerns.

\medterm{Breast Finding} Any noticeable breast change, such as a lump, pain, skin change, nipple inversion, or discharge. Its significance depends on the examination and imaging and sometimes requires a biopsy.

\medterm{Breast Health} Care and awareness aimed at recognising breast changes and detecting disease early. It includes awareness of what is normal, screening appropriate for age and risk, and prompt reporting of a new lump, discharge, or persistent skin change.

\medterm{Breast Imaging-Reporting And Data System} A standard system for describing mammogram, breast ultrasound, and breast MRI findings. Its numbered categories indicate whether results are normal, benign, probably benign, suspicious, or already known to be cancer and help guide follow-up.

\medterm{Breast Infection} Infection and inflammation of breast tissue, often occurring during breastfeeding. It can cause local pain, warmth, redness, swelling, and fever; worsening symptoms, a pus-filled lump, or feeling very unwell needs medical care because an abscess may develop.

\medterm{Breast Lift} Surgery that raises and reshapes sagging breast tissue by removing extra skin and, when needed, repositioning the nipple. It changes shape rather than necessarily increasing breast size and can leave scars or alter sensation.

\medterm{Breast Lobe} One of the larger sections of a breast's milk-producing gland. It contains smaller lobules and ducts that carry milk toward the nipple.

\medterm{Breast Lobule} A small milk-producing unit within the breast, made of tiny milk glands and ducts. Lobules enlarge and become more active during pregnancy and breastfeeding.

\medterm{Breast Lump} A swelling, mass, or thickened area in breast tissue. Most breast lumps (around 80\% to 90\%) are harmless and noncancerous, most often caused by fluid sacs (cysts), noncancerous growths (fibroadenomas), hormone changes, or past injury. Cancerous lumps are less common and usually feel hard, uneven, and fixed in place. Doctors identify the cause through physical examination, imaging such as ultrasound or mammograms, and a needle test (biopsy) if needed.

\medterm{Breast Lump Removal} An operation to remove a breast lump, usually so it can be examined or because it causes symptoms or concern. The need for surgery depends on examination, imaging, biopsy results, and the person's preferences.

\medterm{Breast Lumps} Alternate name; see \textit{Breast Lump}.

\medterm{Breast Milk} Milk made by the breasts after childbirth. It provides energy, nutrients, fluid, and immune substances that support an infant's growth and help protect against some infections.

\medterm{Breast Milk Expression and Storage} The manual or mechanical extraction of human breast milk and its subsequent storage following strict hygiene, refrigeration, and freezing protocols to preserve bioactive immunoglobulins, micronutrients, and enzymatic properties while preventing microbial contamination.

\synonyms
Pumping and Storing Breast Milk; Breast Milk Storage; Milk Expression

\medterm{Breast Milk Jaundice} Yellowing of the skin and eyes that begins after the first week in an otherwise well, feeding, and growing breastfed infant. It is usually temporary, but every newborn with jaundice needs assessment to exclude dangerous causes; feeding is generally continued unless clinicians advise otherwise.

\synonyms
Human Milk Jaundice; Late-Onset Neonatal Jaundice

\medterm{Breast MRI} An imaging test using a strong magnet and radio waves to create detailed pictures of breast tissue. Contrast is often used; MRI is added for selected screening, implant, or diagnostic questions rather than replacing mammography for everyone.

\medterm{Breast MRI Scan} A scan that uses magnetic fields and radio waves, sometimes with injected contrast, to show breast tissue in detail. It can help assess high cancer risk, an abnormality, implants, or the extent of known cancer.

\medterm{Breast Neoplasm} An abnormal growth of cells in breast tissue. It may be noncancerous or cancerous; imaging and often a biopsy determine its type, behaviour, extent, and treatment.

\medterm{Breast Pain} Pain or tenderness in one or both breasts. Common causes include hormonal changes, cysts, infection, injury, medicines, or chest-wall problems; persistent focal pain or pain with a lump needs assessment.

\synonyms
Mastalgia; Mastodynia

\medterm{Breast Rash} A change such as redness, scaling, itching, or irritation of breast skin. Causes include dermatitis, allergy, infection, or rarely inflammatory breast cancer; a persistent, one-sided, painful, or nipple-associated rash needs assessment.

\medterm{Breast Reconstruction} Breast Reconstruction is surgery to rebuild the shape and appearance of a breast, most often after mastectomy. It may use a breast implant, tissue transferred from another part of the body, or both. Reconstruction may begin during mastectomy (immediate reconstruction) or be performed later (delayed reconstruction). It may also include reconstruction of the nipple and areola and may require several procedures.

\medterm{Breast Reconstruction, Autologous} Surgical reconstruction of the breast contour utilizing the patient's own vascularized tissue flaps harvested from donor sites such as the abdomen (e.g., deep inferior epigastric perforator [DIEP] or transverse rectus abdominis myocutaneous [TRAM] flap), back (latissimus dorsi flap), or gluteal/thigh regions.

\synonyms
Autologous Breast Reconstruction; Flap Breast Reconstruction; Autogenous Tissue Reconstruction; DIEP Flap Reconstruction

\medterm{Breast Reconstruction, Implant-Based} Surgical reconstruction of the breast mound following mastectomy or trauma utilizing silicone gel-filled or sterile saline-filled prosthetic implants, often preceded by submuscular or prepectoral placement of a temporary tissue expander.

\synonyms
Prosthetic Breast Reconstruction; Implant Breast Reconstruction; Tissue Expander Breast Reconstruction

\medterm{Breast Reduction} Surgery that removes breast tissue, fat, and skin to reduce breast size and relieve symptoms such as back, neck, or shoulder pain. Possible effects include scars, altered sensation, breastfeeding difficulty, and uneven shape.

\medterm{Breast Screening} Testing people without breast symptoms to find breast cancer early, most commonly with mammography. The recommended starting age and interval depend on local guidance, personal risk, family history, and previous findings.

\medterm{Breast Self-Examination} Looking at and feeling one's own breasts to notice a new lump, skin or nipple change, or discharge. It may improve familiarity with normal changes but does not replace recommended screening or clinical assessment.

\medterm{Breast Skin And Nipple Changes} Changes such as dimpling, redness, scaling, inversion, crusting, or discharge may result from benign conditions, infection, hormonal changes, or cancer. New, persistent, one-sided, or bloody changes require assessment.

\medterm{Breast Surgeon} A surgeon who diagnoses and treats breast conditions, including benign disease, breast cancer, reconstruction, and operations that reduce cancer risk.

\medterm{Breast Surgery} An operation on the breast to remove or examine a lump, treat cancer, reduce size, reconstruct shape, or address another condition. The type of surgery depends on the diagnosis, extent of disease, and person's goals.

\medterm{Breast Ultrasound} Imaging that uses sound waves to examine breast tissue. It helps assess a lump, distinguish fluid from solid tissue, investigate a mammogram finding, guide a biopsy, and evaluate selected implants or dense breasts.

\medterm{Breast-Conserving Treatment} An oncological treatment approach for early-stage breast cancer that aims to treat the cancer effectively while preserving the natural appearance of the breast. It consists of breast-conserving surgery (a lumpectomy, wide local excision, or quadrantectomy to completely remove the tumor with clear microscopic tissue margins and assess sentinel lymph nodes) followed by adjuvant whole-breast radiation therapy to destroy any remaining microscopic cancer cells. Large randomized clinical trials demonstrate that breast-conserving therapy yields long-term survival rates equivalent to complete mastectomy. Contraindications include multicentric disease across multiple quadrants, extensive diffuse microcalcifications, prior chest wall radiation, or persistently positive surgical margins.

\synonyms
Breast-Conserving Therapy; BCT; Breast-Conserving Surgery; BCS; Lumpectomy with Radiotherapy

\medterm{Breastfeeding} Feeding an infant human milk from the breast. It provides nutrition and immune protection; many health organisations recommend exclusive breastfeeding for about six months, followed by complementary foods while breastfeeding continues as desired and feasible.

\medterm{Breastfeeding With Rheumatoid Arthritis} A person with rheumatoid arthritis may be able to breastfeed. Disease activity, fatigue, hand function, and medicine safety are reviewed because some treatments are compatible and others require avoidance or careful timing.

\medterm{Breastfeeding-Associated Nipple and Skin Changes} Dermatologic and mucosal complications of the nipple-areolar complex occurring during lactation, including painful erythema, skin fissures, milk blisters (blebs), contact dermatitis, Raynaud phenomenon of the nipple, and infectious mastitis or candidiasis.

\synonyms
Nipple Trauma in Breastfeeding; Sore Nipples; Nipple Fissures in Lactation

\medterm{Breath} Air moved into and out of the lungs. Its smell, depth, rate, pattern, and oxygen or carbon-dioxide content can provide clues about breathing, infection, metabolism, or other health problems.

\medterm{Breath Alcohol Test} A test measuring alcohol in exhaled air to estimate alcohol concentration in the blood. Results depend on timing, the device, sampling, and health factors and may be used in clinical or legal settings.

\medterm{Breath Analysis} Measurement of gases or other chemicals in exhaled air. Different tests can estimate alcohol, assess airway inflammation, detect infection, or study metabolism; results depend on the method and usually complement rather than replace standard testing.

\medterm{Breath Focus} A mindfulness or relaxation exercise that directs attention to the feeling of breathing. It may help some people reduce stress or emotional arousal but does not replace treatment for a medical or mental-health condition.

\synonyms
Breath-Focused Meditation; Breathing-Focused Mindfulness

\medterm{Breath Odor} An unpleasant smell from the mouth or exhaled air. Common causes include dental or gum disease, dry mouth, tobacco, food, medicines, and infection; persistent odor requires assessment by a dentist or clinician.

\medterm{Breath Sounds} Sounds made as air moves through the airways and lungs, heard with a stethoscope. Reduced, absent, or unusual sounds can provide clues about mucus, narrowing, fluid, collapse, or other lung problems.

\medterm{Breath Stacking} A technique in which several small breaths are taken without fully exhaling, increasing the amount of air in the lungs. It may improve cough strength in selected people with weak breathing muscles and is taught by a professional.

\medterm{Breath Test} Analysis of exhaled air to detect a substance made by the body or by microorganisms. Examples include tests for Helicobacter pylori and tests measuring hydrogen or methane after a sugar drink to assess digestion and absorption.

\medterm{Breath, Bad} Alternate name; see \textit{Bad Breath}.

\medterm{Breath-Holding Spell} A brief involuntary episode in a young child, usually after crying, pain, fear, or anger, in which breathing pauses and the child may change colour or briefly faint before recovering.

\medterm{Breath-Holding Spells} Brief involuntary episodes, usually between about six months and six years, triggered by upset, pain, or fear. A child may cry, exhale, change colour, stiffen or faint briefly, then recover; a first or atypical episode needs assessment.

\medterm{Breathalyzer} A device that estimates alcohol concentration in the blood by measuring alcohol in exhaled air. Results can be affected by timing, device accuracy, recent drinking, mouth alcohol, and sampling conditions.

\medterm{Breathing} Movement of air into and out of the lungs. It brings oxygen into the body and removes carbon dioxide through coordinated activity of the brain, nerves, muscles, and airways.

\medterm{Breathing Difficulty} Uncomfortable or difficult breathing, also called breathlessness. Causes include lung or heart disease, airway blockage, anaemia, infection, anxiety, and metabolic illness; sudden, severe, or worsening difficulty needs urgent assessment.

\medterm{Breathing Difficulty First Aid} Immediate emergency supportive care provided to an individual experiencing acute respiratory distress, including positioning the patient upright in a tripod or sitting posture, loosening restrictive clothing, assisting with prescribed bronchodilator or epinephrine autoinjector administration, and promptly contacting emergency medical services.

\synonyms
First Aid for Dyspnea; Emergency Respiratory Support; Respiratory Distress Management

\medterm{Breathing Difficulty, Recumbent} Alternate name; see \textit{Orthopnea}.

\medterm{Breathing Exercises} Deliberate changes in breathing pattern used to support relaxation, lung expansion, airway clearance, breath control, or rehabilitation. The safest technique depends on the person's condition and is taught when needed.

\medterm{Breathing, Slowed or Arrested} Pathologic depression of the respiratory rate (bradypnea) or complete cessation of ventilatory airflow (apnea), resulting from central nervous system depression, drug toxicity (such as opioid overdose), traumatic brain injury, upper airway obstruction, or terminal cardiopulmonary arrest requiring emergent airway management and resuscitation.

\synonyms
Bradypnea; Apnea; Respiratory Arrest; Hypopnea

\medterm{Breathlessness} Alternate name; see \textit{Dyspnea}.

\medterm{Breech Baby} A fetus positioned in the womb with the buttocks or feet pointing downward toward the birth canal instead of the head. Most babies turn into a head-down (vertex) position before late pregnancy, but those remaining breech near term may be delivered by cesarean section or planned vaginal breech birth, or turned via external cephalic version.
\synonyms Breech fetus; Breech presentation.

\medterm{Breech Birth} Birth in which the baby is positioned bottom or feet first rather than head first. The safest delivery plan depends on the exact position, pregnancy stage, baby and parent factors, and the team's experience.

\medterm{Breech Presentation} A fetal position in which the buttocks or feet point toward the birth canal. Types include frank, complete, and footling breech; risks include cord prolapse, difficulty delivering the head, and birth injury, so delivery needs careful planning.
\synonyms Breech position; Malpresentation.

\medterm{Bregma} The point on the skull where the main seam between the frontal and parietal bones meets the seam between the two parietal bones. In infants it lies beneath the soft spot at the front of the skull.

\medterm{Bremsstrahlung} Radiation produced when fast electrons slow down or change direction near matter. It is relevant to some imaging and radiation treatments and requires appropriate shielding and safety controls.

\medterm{Brenner Tumor} A rare ovarian growth made of cells resembling those lining the urinary tract. Most are small, one-sided, noncancerous, and found incidentally, but uncommon borderline or cancerous forms also occur.

\medterm{Brenner Tumour} A usually noncancerous ovarian growth whose cells resemble those lining the urinary tract. Rare borderline and cancerous forms occur and are distinguished by tissue examination.

\medterm{Breslau Bacillus} An older name for Salmonella Typhimurium, a bacterium that can cause foodborne gastroenteritis with diarrhoea, abdominal cramps, fever, and vomiting.

\medterm{Breslow Thickness} The depth, measured in millimetres, to which a melanoma has grown into the skin. It is an important part of staging and helps estimate the risk of spread and guide treatment.

\medterm{Bretylium Tosylate} An older medicine that can suppress certain dangerous abnormal heart rhythms, especially ventricular fibrillation. It is now rarely used because of limited availability, side effects, and newer treatments.

\medterm{Brevibacillus} A group of spore-forming bacteria found in soil, water, and food. Human infection is uncommon but can occur opportunistically, especially in people with weakened immunity or medical devices.

\medterm{Brevibacterium} A group of bacteria found on skin, in soil, and in food. They commonly contribute to body or cheese odour; rare species can cause opportunistic infection, particularly in people with serious illness or devices.

\medterm{Brevibacterium Casei} A species of Brevibacterium found in the environment and occasionally on people. It rarely causes bloodstream or device-related infection, usually in a person with major underlying illness or an indwelling device.

\medterm{Brevundimonas} A group of environmental water-associated bacteria that rarely cause human disease. Opportunistic infections, including bloodstream or wound infection, are more likely in people with weakened immunity or medical devices.

\medterm{Brevundimonas Diminuta} An environmental water-associated bacterium that rarely causes opportunistic human infection. Reported illnesses include bloodstream, wound, or device-related infection, mainly in people with serious underlying disease.

\medterm{Brevundimonas Vesicularis} An environmental bacterium that rarely causes human infection. When it is isolated from a clinical sample, symptoms, specimen quality, and other results determine whether it represents true disease or contamination.

\medterm{Bricklayer's Disease} An older occupational term for lung disease caused by breathing dust during masonry work, especially silica-containing dust. Long-term exposure can cause silicosis, which scars the lungs and impairs breathing.

\medterm{Bridge To Transplant} Temporary treatment or mechanical support that keeps a person with advanced heart failure alive and stable until a heart transplant can be performed. The device, timing, and risks depend on heart function and overall health.

\medterm{Bridge Work} A fixed dental replacement for one or more missing teeth. Artificial teeth are attached to prepared neighbouring teeth or implants, restoring some chewing function and appearance.

\medterm{Brief Psychotic Disorder} A sudden episode of delusions, hallucinations, severely disorganized speech, or markedly disorganized behavior lasting at least one day but less than one month, followed by a return toward usual functioning.

\synonyms
Brief Reactive Psychosis; Acute Transient Psychosis

\medterm{Brief Resolved Unexplained Event} A sudden, short, and resolved episode in an infant involving abnormal breathing, skin colour, muscle tone, or responsiveness, with no explanation after an initial assessment. Risk depends on the infant's age, history, examination, and whether it recurs.

\medterm{Brief Resolved Unexplained Event} A sudden, brief (under 1 minute), fully resolved episode in an infant under one year of age characterized by one or more of: cyanosis or pallor, absent, decreased, or irregular breathing, marked change in tone (hypertonia or hypotonia), and altered level of responsiveness, with no explanatory underlying cause identified following a thorough history and physical examination.

\synonyms
BRUE; Apparent Life-Threatening Event; ALTE

\medterm{Bright's Disease} An old term for several kidney diseases, especially conditions affecting the kidney's filters and causing protein in urine, swelling, or altered urine. Modern diagnosis identifies the specific kidney problem instead.

\medterm{Brill's Disease} Recurrence of epidemic typhus years after the first infection because dormant Rickettsia prowazekii becomes active again. It is usually milder than the original illness but can still be spread by body lice.

\medterm{Brill-Zinsser Disease} Recurrent epidemic typhus caused by reactivation of dormant Rickettsia prowazekii years after the original infection. The illness is usually milder but can be transmitted by body lice.

\medterm{Brimonidine Tartrate} A medicine that reduces certain nerve signals and blood-vessel activity. Eye drops lower pressure inside the eye, and skin preparations can reduce some facial redness; side effects and interactions may require medical attention.

\medterm{Brine Bath} Bathing in salty water, sometimes used for comfort or certain skin conditions. It does not treat every skin problem, and hot water or salt can worsen irritation, dry skin, wounds, or some medical conditions.

\medterm{Brittle Asthma} An uncommon term for severe asthma with sudden, unpredictable, or difficult-to-control attacks despite appropriate treatment. Specialist review is needed to confirm asthma, identify triggers or other diagnoses, and reduce risk.

\medterm{Brittle Bone Disease} An inherited disorder, usually osteogenesis imperfecta, in which bones break easily. It can also cause short stature, loose joints, blue-gray eye whites, hearing loss, dental problems, or other skeletal changes.

\medterm{Brittle Bone Disease} Alternate name; see \textit{Osteogenesis Imperfecta}.

\medterm{Brittle Diabetes} Alternate historical term; see \textit{Type 1 Diabetes}.

\medterm{Brivudine} An antiviral medicine used in some countries for shingles. It is not used with fluoropyrimidine cancer medicines or certain related treatments because the interaction can be fatal.

\medterm{Broad Ligament} A sheet of tissue extending from the sides of the uterus to the pelvic walls. It helps support the uterus, fallopian tubes, and ovaries and contains blood vessels, nerves, and connective tissue.

\medterm{Broad Nasal Bridge} An unusually wide upper part of the nose between the eyes. It may be a normal family feature or part of a condition present from birth or an inherited syndrome.

\medterm{Broad-Based Gait} An abnormal walking pattern in which the feet are placed abnormally far apart to increase the base of support and maintain balance. It is a cardinal clinical physical sign of midline cerebellar vermis lesions (cerebellar ataxia) or severe loss of lower limb proprioception (sensory ataxia).

\medterm{Broca Aphasia} A language disorder in which speech is slow, effortful, and difficult to produce, while understanding may be relatively preserved. It usually follows damage to the dominant frontal language area, often from a stroke.

\medterm{Broca's Aphasia} A nonfluent language disorder caused by damage to the dominant frontal language region. Speech is effortful and limited, while understanding may be better preserved than speaking or repeating words.

\medterm{Broca's Area} A region in the dominant frontal lobe that helps plan and produce spoken language. Damage can cause slow, effortful speech and difficulty forming words, while understanding may be relatively preserved.

\medterm{Brodie Abscess} A chronic, well-circumscribed subacute pyogenic bone abscess that represents a localized form of osteomyelitis. It characteristically develops in the metaphysis of long bones (most frequently the proximal or distal tibia) in children and young adults, presenting as a lytic bone lesion with a surrounding ring of sclerotic bone on radiography.

\medterm{Brodie's Abscess} A localized, slowly developing pocket of bone infection, often near the end of a long bone. It commonly causes persistent, sometimes night-time, bone pain with little fever and may require prolonged antibiotics or surgery.

\medterm{Brody Disease} A rare inherited muscle disorder causing painless stiffness and slow muscle relaxation after strenuous or repeated exercise. It often begins in childhood and may affect the face, eyelids, or hands; symptoms can resemble myotonia but arise from a different muscle problem.

\synonyms
Brody Myopathy; SERCA1-Related Muscle Disease

\medterm{Broken Ankle} A break in one or more bones forming the ankle joint. It usually causes pain, swelling, bruising, and difficulty bearing weight; deformity, an open wound, numbness, or poor circulation requires emergency care.

\medterm{Broken Arm} A break in one or more arm bones, usually causing pain, swelling, bruising, deformity, or loss of function. An open injury, severe deformity, numbness, or reduced circulation requires urgent care.

\medterm{Broken Blood Vessel in Eye} Alternate name; see \textit{Subconjunctival Hemorrhage}.

\medterm{Broken Bone} A partial or complete break in a bone caused by sudden injury, repeated stress, or weakened bone. Pain, swelling, bruising, deformity, or loss of function may occur and imaging usually confirms the diagnosis.

\medterm{Broken Bone} Alternate name; see \textit{Bone Fracture}.

\medterm{Broken Clavicle} Alternate name; see \textit{Broken Collarbone}.

\medterm{Broken Collarbone} A break in the clavicle, usually after a fall or direct blow. It causes pain, swelling, tenderness, and reduced shoulder movement; numbness, breathing difficulty, or a skin-threatening deformity requires urgent care.

\synonyms
Broken Clavicle

\medterm{Broken Finger} A break in a finger bone caused by a blow, crush, or twist. Pain, swelling, bruising, deformity, and reduced movement may occur; assessment checks alignment, tendon function, sensation, and circulation and guides splinting or surgery.

\synonyms
Finger Fracture

\medterm{Broken Foot} A break in one or more foot bones after a fall, blow, twist, or repeated stress. Pain, swelling, bruising, and difficulty walking are common; deformity, an open wound, numbness, or inability to walk needs urgent assessment.

\medterm{Broken Hand} A break in one or more hand bones, causing pain, swelling, bruising, deformity, or reduced movement. Assessment includes alignment, tendon function, sensation, and circulation because some breaks require reduction or surgery.

\medterm{Broken Heart Syndrome} Temporary weakening of the heart muscle, often after severe emotional or physical stress. It can resemble a heart attack with chest pain or breathlessness, but the coronary arteries usually have no blocking clot.

\medterm{Broken Hip} Alternate name; see \textit{Hip Fracture}.

\medterm{Broken Leg} A break in one or more leg bones, usually causing pain, swelling, bruising, deformity, or inability to bear weight. An open injury, numbness, or loss of circulation requires emergency care.

\synonyms
Leg Fracture

\medterm{Broken Little Finger} Alternate name; see \textit{Broken Finger}.
\medterm{Broken Nose} A break in one or more nasal bones that may cause pain, swelling, bruising, bleeding, deformity, or difficulty breathing through the nose. Persistent bleeding, clear fluid leakage, or severe deformity needs urgent assessment.

\medterm{Broken Or Dislocated Jaw} A break or displacement of the jaw that can cause pain, swelling, an altered bite, limited movement, loose teeth, or difficulty speaking, swallowing, or breathing. It needs urgent medical and dental assessment.

\medterm{Broken Or Knocked Out Tooth} A tooth damaged or completely forced out by injury. A knocked-out permanent tooth is a dental emergency; the tooth is handled by the crown, kept moist, and the root is not scrubbed.

\medterm{Broken Ribs} Breaks in one or more ribs, usually causing local chest pain that worsens with breathing, coughing, or movement. Severe breathing difficulty, increasing pain, coughing blood, or major injury requires urgent assessment.

\medterm{Broken Toe} A break in a toe bone, usually after a direct blow or crush. Pain, swelling, bruising, and tenderness are common; deformity, an open wound, numbness, or poor circulation needs prompt assessment.

\synonyms
Toe Fracture; Phalanx Fracture

\medterm{Broken Wrist} A break in one or more bones near the wrist, often causing pain, swelling, bruising, deformity, or reduced hand function. Numbness, pale or cold fingers, an open wound, or severe deformity needs urgent care.

\medterm{Bromelain} A mixture of enzymes from pineapple. It is sold as a supplement and studied for pain and inflammation, but evidence and product strength vary; it can cause allergy, stomach upset, or interactions with medicines.

\medterm{Bromhexine} A medicine that thins sticky mucus in the airways so it may be coughed up more easily. It can cause stomach upset or rash and rarely causes a serious skin reaction.

\medterm{Bromhidrosis} Unusually strong or unpleasant body odour produced when microorganisms break down sweat or skin secretions. It commonly affects the armpits or feet and may occur with excessive sweating or skin disease.

\medterm{Bromides} Salts or compounds containing bromide. Excessive or prolonged exposure can cause bromism, with symptoms affecting the nervous system, mood, digestion, or skin.

\medterm{Bromine} A reactive chemical used in industry and some flame-resistant materials. Bromine and many bromine compounds can irritate or burn tissue, while substantial exposure can cause poisoning affecting several organs.

\medterm{Bromism} Poisoning caused by prolonged or excessive exposure to bromide compounds. It may cause confusion, drowsiness, poor coordination, mood or behaviour changes, digestive symptoms, acne-like eruptions, or other neurologic problems.

\medterm{Bromobenzenes} Industrial chemicals made by attaching one or more bromine atoms to benzene. Exposure can irritate the eyes, skin, and airways or cause other toxic effects, depending on the specific compound and amount.

\medterm{Bromocriptine} A medicine that stimulates dopamine receptors and reduces prolactin release. It is used for high prolactin and selected movement or hormone disorders; nausea, dizziness, low blood pressure, hallucinations, and other effects may occur.

\medterm{Bromoderma} An inflammatory skin eruption linked to bromine or bromide exposure. It may appear as acne-like, pustular, or other irritated lesions and usually improves when the source is identified and removed.

\medterm{Bromodosis} Very unpleasant foot odour caused when sweat, moisture, and skin debris are broken down by bacteria or fungi. Washing, drying, breathable footwear, and treatment of any infection may help.

\medterm{Bromophos} An organophosphate insecticide. Exposure can overstimulate nerves and cause sweating, excess saliva, vomiting, breathing difficulty, weakness, seizures, or collapse and requires urgent poison or emergency care.

\medterm{Brompheniramine Overdose} Poisoning from taking too much brompheniramine, an antihistamine. It may cause severe sleepiness or agitation, confusion, dry mouth, rapid heartbeat, seizures, or dangerous heart rhythms; urgent help is needed.

\medterm{Bromthymol Blue} A chemical indicator that changes colour according to acidity and alkalinity. It is used in laboratories and is not a medicine or treatment.

\medterm{Bronchi} The two main airways carrying air from the windpipe into the lungs. Each divides into smaller branches; the right main bronchus is wider and more vertical, so inhaled objects are more likely to enter it.

\medterm{Bronchial} Relating to the bronchi, the large airways branching from the windpipe. Bronchial disease may cause cough, wheezing, mucus production, narrowing, or reduced airflow.

\medterm{Bronchial Artery} An artery from the body's general circulation that supplies the walls of the bronchi and supporting lung tissues. It differs from the pulmonary arteries, which carry blood to the lungs for oxygen exchange.

\medterm{Bronchial Artery Embolization} A catheter procedure that blocks bleeding arteries supplying the airways, usually to control severe or repeated coughing of blood. Imaging guides the catheter and the material used; careful technique is needed to avoid blocking arteries supplying the spinal cord.

\synonyms
BAE; Bronchial Embolization; Embolization of Bronchial Arteries

\medterm{Bronchial Asthma} A long-term condition in which sensitive, inflamed airways narrow intermittently. It can cause wheezing, cough, chest tightness, and breathlessness; symptoms often vary with triggers and improve with appropriate treatment.

\medterm{Bronchial Biopsy} Removal of a small tissue sample from an airway during bronchoscopy. Laboratory examination can investigate inflammation, infection, abnormal growths, or cancer; bleeding and breathing complications are possible.

\medterm{Bronchial Breath Sounds} Harsh, high-pitched, hollow breath sounds with a prolonged expiratory phase and an audible pause between inspiration and expiration. Normally auscultated over the trachea and main bronchi, their presence over peripheral lung fields indicates abnormal transmission through consolidated lung tissue, as in lobar pneumonia.

\medterm{Bronchial Carcinoid} A usually slow-growing neuroendocrine tumour arising in an airway. Typical carcinoids are generally less aggressive than atypical carcinoids, but either can block an airway, cause coughing or bleeding, or spread to other sites.

\medterm{Bronchial Fistula} An abnormal passage between an airway and another structure, such as the space around the lung, food pipe, or skin. It can cause persistent air leakage, coughing, infection, or breathing failure.

\medterm{Bronchial Hemorrhage} Bleeding from airway blood vessels, usually seen as coughing up blood. Heavy bleeding can block the airway or cause breathing failure and requires emergency treatment.

\medterm{Bronchial Obstruction} Partial or complete blockage of an airway by mucus, an inhaled object, a tumour, swelling, or pressure from outside. It can reduce airflow and lead to lung collapse, repeated infection, wheezing, or breathlessness.

\medterm{Bronchial Stenosis} Narrowing of an airway caused by scarring, inflammation, infection, injury, a tumour, or a previous procedure. It can cause localised wheezing, breathlessness, repeated infection, or poor mucus clearance and may require widening or surgery.

\medterm{Bronchial Stricture} Abnormal narrowing of an airway caused by scarring, inflammation, infection, a tumour, or previous instrumentation. It may produce localised wheezing, cough, repeated infection, or collapse of the lung beyond the narrowing.

\medterm{Bronchial Tree} The branching network of airways from the windpipe through the main, lobar, and smaller bronchi and bronchioles toward the air sacs. It conducts and distributes air throughout the lungs.

\medterm{Bronchial Tubes} The two main bronchi and their branches, which carry air from the windpipe into the lungs and divide into progressively smaller airways.

\synonyms
Bronchi; Bronchial Tree When Referring to the Branching Airway System

\medterm{Bronchial Vein} A vein that drains blood from the walls of the larger airways and nearby lung tissue. It empties partly into veins returning blood to the heart and partly into pulmonary veins.

\medterm{Bronchiectasis} Permanent widening and damage of the airways after repeated or severe inflammation and infection. It can cause daily cough, large amounts of mucus, breathlessness, repeated chest infections, and coughing up blood.

\medterm{Bronchiectasis, Acquired / Congenital} Alternate index label; see \textit{Bronchiectasis}.

\medterm{Bronchiole} A small airway without cartilage that branches from a bronchus and carries air toward the lung's air sacs. Its smooth muscle can narrow or widen the passage and regulate airflow.

\medterm{Bronchiolectasis} Permanent widening of the small airways, usually caused by chronic inflammation, infection, or blockage. It may occur with bronchiectasis and contribute to poor mucus clearance and repeated infection.

\medterm{Bronchioles} Small branching airways that carry air from the bronchi toward the air sacs and regulate airflow with their surrounding muscle.

\medterm{Bronchiolitis} Inflammation and blockage of the smallest airways, usually from a viral infection in infants and young children. It can cause cough, wheezing, fast breathing, chest retractions, poor feeding, and low oxygen.

\medterm{Bronchiolitis Obliterans} A rare disease in which inflammation and scarring permanently narrow or close the smallest airways. It can follow severe infection, inhaled injury, autoimmune disease, or transplantation and causes persistent cough and breathlessness.

\medterm{Bronchitis} Inflammation of the larger airways. Acute bronchitis is usually caused by a virus and causes cough, with or without mucus, while chronic bronchitis means mucus-producing cough for at least three months in two successive years.

\medterm{Bronchitis, Acute} Alternate name; see \textit{Acute Bronchitis}.

\medterm{Bronchitis, Chronic} Alternate name; see \textit{Chronic Bronchitis}.

\medterm{Bronchoalveolar Lavage} A bronchoscopy procedure in which sterile fluid is put into a small lung area and then collected. The sample can be tested for infection, inflammation, abnormal cells, or other lung disease.

\medterm{Bronchocele} Abnormal widening of a bronchus filled with mucus, usually because an airway is blocked or developed abnormally. It may cause cough, repeated infection, or no symptoms and is assessed with imaging.

\medterm{Bronchoconstriction} Narrowing of the airways because their muscle tightens, the lining swells, or mucus builds up. It increases resistance to airflow and commonly occurs in asthma, allergic reactions, or irritant exposure.

\medterm{Bronchoconstrictor Agents} Substances that narrow the airways by tightening airway muscle or causing inflammation. Examples include histamine, leukotrienes, irritant gases, and some medicines.

\medterm{Bronchodilator} A medicine that relaxes airway muscle and widens breathing tubes. It can relieve or prevent wheezing and breathlessness in asthma, chronic obstructive lung disease, and other conditions with narrowed airways.

\medterm{Bronchodilator Response} The change in airflow measured by spirometry after a bronchodilator is given. Improvement can support variable airway narrowing, but no response does not exclude asthma and a response does not prove one specific disease.

\medterm{Bronchodilator Reversibility} Improvement in measured airflow after an inhaled bronchodilator. It can support variable airway narrowing, but one result alone cannot confirm or exclude asthma or chronic obstructive pulmonary disease.

\medterm{Bronchoesophageal Fistula} An abnormal connection between an airway and the food pipe. It may be present from birth or develop from cancer, infection, injury, or treatment and can cause coughing during swallowing, repeated pneumonia, or breathing problems.

\medterm{Bronchogenic Carcinoma} A cancer that begins in the airways or nearby lung tissue. It may cause cough, coughing blood, chest pain, breathlessness, or no early symptoms; diagnosis and treatment depend on the cancer type and extent.

\medterm{Bronchogenic Cyst} A birth-related fluid-filled cyst that develops from early airway or food-pipe tissue, often near the windpipe or central chest. It may cause no symptoms or press on nearby structures and sometimes requires removal.

\medterm{Bronchomalacia} Weakness of an airway wall that allows it to collapse too much, especially during breathing out. It may cause noisy breathing, cough, breathlessness, repeated infections, or difficulty clearing mucus.

\medterm{Bronchophony} An examination finding in which spoken sounds are heard unusually clearly through a stethoscope over part of the lung. It can occur when lung tissue is solidified by infection or another process but is not diagnostic by itself.

\medterm{Bronchopleural Fistula} An abnormal connection between an airway and the space around the lung, often after lung surgery or severe infection. It can cause a persistent air leak, breathlessness, infection, or air under the skin.

\medterm{Bronchopneumonia} Infection causing scattered areas of inflammation and consolidation around the small airways in one or more lung lobes. It may cause fever, cough, mucus, chest pain, fast breathing, and low oxygen.

\medterm{Bronchopulmonary Dysplasia} A chronic lung disease in some premature infants whose developing lungs have been affected by immaturity, inflammation, oxygen, or breathing support. It may cause ongoing oxygen need, fast breathing, or wheezing.

\medterm{Bronchopulmonary Segments} Natural sections of a lung lobe, each supplied by its own airway branch and pulmonary artery branch. Their separation helps doctors describe disease and allows some operations to remove only an affected segment.

\medterm{Bronchoscope} A thin tube with a light and camera used to examine the windpipe and airways. Instruments passed through it can remove an object or collect tissue, fluid, or mucus for testing.

\medterm{Bronchoscopic Culture} Laboratory testing of fluid or tissue collected during bronchoscopy to look for bacteria, fungi, or mycobacteria when sputum or other ordinary samples are not sufficient.

\medterm{Bronchoscopy} Examination of the windpipe and airways with a flexible or rigid camera tube. It can investigate persistent cough, infection, bleeding, or a lung mass and can remove foreign material or collect samples; sedation or anaesthesia is used.

\medterm{Bronchospasm} Sudden tightening of airway muscles that narrows the breathing passages. It can cause wheezing, cough, chest tightness, and difficulty breathing and may be triggered by asthma, allergy, infection, irritants, or airway procedures.

\medterm{Bronchostenosis} Narrowing of a bronchus that restricts airflow. Causes include scarring, inflammation, infection, a tumour, an inhaled object, pressure from outside, or a previous procedure and may cause wheezing or repeated infection.

\medterm{Bronchus} One of the two main branches of the windpipe that carries air into a lung. Each divides into smaller branches; the right main bronchus is wider and more vertical, making it a more common site for inhaled objects to lodge.

\medterm{Bronchus Disease} Any condition affecting one or more bronchi, the air passages leading into the lungs. Infection, asthma, chronic inflammation, blockage, scarring, or a tumour may reduce airflow and cause cough, wheezing, or repeated infection.

\medterm{Bronchus Neoplasm} An abnormal growth arising in a bronchus. It may be noncancerous or cancerous; imaging, bronchoscopy, and tissue testing help determine its type, extent, and treatment.

\medterm{Bronze Diabetes} Historical name for hereditary hemochromatosis; see \textit{Hemochromatosis}.

\medterm{Brooke Disease} A rare inherited skin disorder involving abnormal thickening of the outer skin or hair follicles. Because the eponym has been used for more than one condition, the original source and clinical features are checked.

\medterm{Brooke-Spiegler Syndrome} An inherited disorder causing multiple usually noncancerous skin tumours, often on the head and neck. It is commonly linked to changes in the CYLD gene and may include cylindromas, spiradenomas, or trichoepitheliomas.

\medterm{Broselow Tape} A colour-coded tape that estimates a child's weight from body length during an emergency. It helps select medicine doses, fluid amounts, equipment sizes, and defibrillator settings, but clinical assessment and an actual weight are used when available.

\synonyms
Broselow Pediatric Emergency Tape; Pediatric Length-Based Resuscitation Tape

\medterm{Brow Presentation} A fetal position in which the forehead leads because the head is partly extended. The head presents a wide diameter, so persistent brow presentation at term may prevent vaginal birth and requires specialist delivery planning.

\medterm{Brown Fat} A type of body fat containing many energy-producing structures that generate heat. It is especially important for keeping newborns warm and can also be activated by cold in adults.

\medterm{Brown Lung Disease} An occupational breathing disorder caused by repeated exposure to cotton, flax, or hemp dust, traditionally called byssinosis. It can cause chest tightness, cough, wheezing, and work-related narrowing of the airways.

\medterm{Brown Recluse Spider} A venomous spider whose bite may initially be mild but can later cause painful skin damage, blistering, or tissue death. Rarely, its venom causes fever, red-cell destruction, kidney injury, or other systemic illness.

\medterm{Brown Recluse Spider Bite} Alternate name; see \textit{Spider Bites}.

\medterm{Brown Sequard Syndrome} A spinal-cord injury affecting mainly one side. It typically causes weakness and loss of vibration or position sense on the injured side and loss of pain and temperature sensation on the opposite side below the injury.

\medterm{Brown Stain} A brown deposit on teeth caused by pigments from food, drinks, tobacco, or bacteria. Professional cleaning may remove surface stains, but decay or enamel disease needs separate treatment.

\medterm{Brown Syndrome} Limited ability to lift the eye when it is turned inward, because the superior oblique tendon cannot glide freely through its pulley, the trochlea. The affected eye appears lower than the other in that gaze direction and double vision may occur. It can be present at birth or develop from inflammation, injury, sinus disease or surgery, or rheumatic disease such as rheumatoid arthritis. Some cases resolve on their own; persistent double vision is treated by an eye specialist.

\medterm{Brown Teeth} Brown discolouration of teeth caused by surface stains, decay, enamel defects, excess fluoride, medicines, injury, or other conditions. A dental examination identifies the cause and appropriate treatment.

\medterm{Brown-SéQuard Syndrome} A spinal-cord injury affecting mainly one side. It causes weakness and loss of position or vibration sense on the injured side and loss of pain and temperature sensation on the opposite side below the injury.

\medterm{Bruce Protocol} A graded treadmill test in which speed and slope increase at set stages. It assesses exercise capacity, symptoms, heart-rate and blood-pressure responses, and signs of reduced blood flow to the heart.

\medterm{Brucella} A group of bacteria that usually infect animals but can infect people through unpasteurised dairy products, inhalation, or contact with infected animals or tissues. Human infection causes brucellosis.

\medterm{Brucella Abortus} A Brucella species mainly associated with cattle. It can infect people exposed to infected animals, tissues, or unpasteurised dairy products and cause prolonged fever and other brucellosis symptoms.

\medterm{Brucella Canis} A Brucella species associated with dogs. It rarely infects people after contact with infected dogs or reproductive material and may cause prolonged fever or other brucellosis symptoms.

\medterm{Brucella Melitensis} A Brucella species mainly associated with goats and sheep and a frequent cause of severe human brucellosis, usually acquired through unpasteurised dairy products or infected animals.

\medterm{Brucella Suis} A Brucella species associated with pigs and some wildlife. It can infect people through animal contact or contaminated food and may cause prolonged illness or infection in bones and organs.

\medterm{Brucellaceae} A family of bacteria that includes Brucella and related animal-associated species. Some members cause brucellosis, an infection that can pass from animals to people.

\medterm{Brucellosis} An infection caused by Brucella bacteria, usually acquired from infected animals, animal tissues, or unpasteurised dairy products. It may cause prolonged fever, sweats, tiredness, headache, and muscle or joint pain.

\medterm{Bruch's Membrane} A thin supporting layer between the retina and the blood-rich tissue behind it. It helps nourish and maintain the retina and can become damaged in ageing, diabetes, and some retinal diseases.

\medterm{Bruck's Syndrome} A rare inherited disorder combining fragile bones with joint stiffness or fixed contractures present from birth. It results from changes affecting connective tissue and may cause repeated fractures and limited movement.

\medterm{Brudzinski Sign} Involuntary bending of the hips and knees when the neck is gently bent forward. It can occur with irritation of the tissues around the brain and spinal cord, including meningitis, but does not confirm the diagnosis.

\medterm{Brugada Phenocopy} A temporary ECG pattern resembling Brugada syndrome but caused by another problem, such as fever, abnormal blood salts, reduced heart blood flow, medicines, or poisoning. The pattern resolves when the cause is corrected and is distinguished from inherited disease.

\synonyms
Acquired Brugada Pattern; BrP

\medterm{Brugada Syndrome} An inherited heart-rhythm disorder that can cause dangerous ventricular arrhythmias and sudden cardiac death, often during rest, sleep, or fever. Diagnosis relies on a characteristic ECG pattern and clinical assessment.

\medterm{Brugia} A group of parasitic worms spread by mosquitoes that can cause lymphatic filariasis. The infection may damage lymph vessels and lead to recurrent fever, swelling, or long-term enlargement of limbs or genital tissues.

\medterm{Bruise} Bleeding under unbroken skin after small blood vessels are damaged, usually by a blow, pressure, or injury. It commonly changes colour as it heals; unexplained or frequent bruising needs medical assessment.

\medterm{Bruises} Areas of bleeding beneath the skin after small blood vessels are damaged. Frequent, large, spontaneous, painful, or unusually persistent bruises may result from medicines, platelet or clotting problems, liver disease, nutritional deficiency, or injury.

\medterm{Bruising} Purplish, blue, green, or yellow skin discolouration caused by blood leaking into tissue, usually after injury. Colour changes as the blood is broken down; widespread or unexplained bruising needs assessment.

\medterm{Bruit} An abnormal whooshing sound heard over a blood vessel, usually with a stethoscope, caused by turbulent blood flow. It may suggest narrowing but does not by itself establish the cause or severity.

\medterm{Brunner Gland} A gland in the first part of the small intestine that releases mucus and alkaline fluid. These secretions protect the intestinal lining and help neutralise acid arriving from the stomach.

\medterm{Brunner's Gland} A mucus- and bicarbonate-producing gland in the duodenum, the first part of the small intestine. Its alkaline secretion protects the lining from stomach acid and supports digestion.

\medterm{Brunner's Glands} Glands in the duodenum that release mucus and bicarbonate-rich fluid to protect the intestinal lining from stomach acid and help create a suitable environment for digestive enzymes.

\medterm{Brush Burns} Superficial friction injuries caused by scraping skin against a rough surface. They cause pain, redness, and grazes; cleaning and protecting the area usually allow healing, but deep, extensive, or infected injuries need medical care.

\medterm{Brushfield's Spots} Small gray or white spots arranged around the coloured part of the eye. They are common in people with Down syndrome but can also occur in healthy people and usually do not affect vision.

\medterm{Bruxism} Repeated clenching or grinding of the teeth, often during sleep. It can wear or fracture teeth and cause jaw, face, ear, or headache pain; dental review can assess damage and possible sleep, stress, or medicine-related causes.

\synonyms
Teeth Grinding; Tooth Clenching

\medterm{Bruxism, Sleep / Awake} Alternate index label; see \textit{Bruxism}.

\medterm{Bryant's Traction} An older treatment in which an infant's or young child's legs are gently elevated with traction to treat selected hip or thigh conditions. It is rarely used and requires careful specialist monitoring.

\medterm{Bubble Bath Soap Poisoning} Illness after swallowing or inhaling bubble-bath product. A small amount usually causes mouth or stomach irritation, but repeated vomiting, coughing, breathing difficulty, or drowsiness requires emergency assessment.

\medterm{Bubo} A tender, enlarged lymph node, usually in the groin, armpit, or neck, caused by infection or inflammation. Buboes are classically linked with plague but also occur in sexually transmitted and other infections.

\medterm{Bubonic Plague} A serious infection caused by Yersinia pestis, usually spread by infected flea bites. It causes fever and painful swollen lymph nodes and can spread to the blood or lungs without prompt antibiotic treatment.

\medterm{Buccal} Relating to the cheek or its inner lining. A buccal medicine is placed between the cheek and gum so it dissolves and is absorbed through the mouth.

\medterm{Buccal Gland} A small salivary gland in the cheek that releases mucus into the mouth. The secretion helps keep the cheek lining moist and supports chewing and swallowing.

\medterm{Buccal Mucosa} The non-keratinized stratified squamous mucous membrane lining the inside of the cheeks, bounded by the labial commissures anteriorly, the retromolar trigone posteriorly, and the alveolar sulci superiorly and inferiorly. It contains the opening of the parotid duct (Stensen duct) opposite the maxillary second molar and serves as a site for clinical examination of exanthems (such as Koplik spots in measles), oral lichen planus (Wickham striae), aphthous stomatitis, leukoplakia, and drug absorption.

\medterm{Buccal Smear} A sample of cells gently collected from the inside of the cheek for laboratory testing. It may be used for genetic analysis or, less commonly, to investigate infection or abnormal cells.

\medterm{Buccalmucosa} A spelling variant of buccal mucosa, the moist lining on the inside of the cheeks.

\medterm{Buccinator} A muscle in the cheek that presses the cheek against the teeth. It helps keep food between the teeth and assists sucking, blowing, chewing, and speaking.

\medterm{Buckle Fracture} Alternate name; see \textit{Torus Fracture}.

\medterm{Budd-Chiari Syndrome} Blockage of veins draining blood from the liver, usually from a blood clot or a condition that makes clotting more likely. It can cause abdominal pain, enlarged liver, fluid in the abdomen, jaundice, and liver failure.

\medterm{Budesonide} A corticosteroid that reduces inflammation. It is used in selected inhaled, nasal, digestive, or rectal treatments for conditions such as asthma, chronic obstructive lung disease, allergic symptoms, and inflammatory bowel disease; side effects depend on the route and dose.

\medterm{Budvicia} A group of water-associated bacteria related to Yersinia. They rarely cause human disease, but some isolates have been associated with urinary or other opportunistic infections.

\medterm{Budvicia Aquatica} A water-associated bacterium that has rarely been isolated from human urinary infections, including infections in people with urinary catheters. Its significance depends on symptoms and specimen quality.

\medterm{Buerger Disease} An inflammatory disease that blocks small and medium blood vessels, especially in the hands and feet. It is strongly linked to tobacco use and can cause pain, ulcers, poor circulation, and loss of tissue.

\medterm{Buerger Test} A bedside physical examination assessment used to evaluate the severity of peripheral arterial occlusive disease in the lower extremities. The patient's leg is elevated to 45 or 90 degrees until maximum pallor develops (Buerger's angle of circulatory sufficiency), and is then placed in a dependent sitting position; a prolonged delay in the return of pink skin color followed by a deep red or purple reactive erythema (rubor of dependency) indicates advanced arterial insufficiency.

\medterm{Buerger's Disease} A tobacco-associated inflammatory disease of small and medium arteries and veins, also called thromboangiitis obliterans. Reduced blood flow can cause pain, colour changes, ulcers, gangrene, and loss of fingers or toes.

\medterm{Buffalo Hump} A pad of fat at the back of the neck and upper back. It may occur with obesity, long-term corticosteroid excess, some HIV treatments, or Cushing syndrome and is assessed together with other findings.

\medterm{Buffer} A chemical solution that resists a change in acidity when small amounts of acid or base are added. Buffers help maintain stable conditions in blood, cells, medicines, and laboratory tests.

\medterm{Buffy Coat} The thin pale layer of white blood cells and platelets that forms between plasma and red blood cells when a blood sample is spun. It can be collected for laboratory testing.

\medterm{Bufonidae} The family of amphibians commonly called true toads. Some species produce bufotoxins in their skin glands; swallowing or absorbing these toxins can affect the heart and nervous system.

\medterm{Buformin} An older biguanide medicine for diabetes that has been withdrawn or restricted in many countries because it can cause dangerous lactic-acid buildup, especially when kidney function is poor.

\medterm{Bug Spray Poisoning} Illness after exposure to an insecticide. Effects depend on the chemical and may include skin or eye irritation, vomiting, breathing difficulty, unusual sweating, confusion, seizures, or other nervous-system problems; severe poisoning requires urgent medical treatment.

\medterm{Bulbar} Relating to the lower brainstem or to the nerves controlling swallowing, speaking, chewing, coughing, and movement of the tongue and soft palate. Bulbar weakness can allow food or saliva to enter the airway.

\medterm{Bulbar Conjunctiva} The part of the eye's clear moist covering that lies over the white of the eye. It joins the inner surface of the eyelids around the front of the eye and can become inflamed or injured.

\medterm{Bulbar Palsy} Weakness or paralysis of muscles controlled by lower cranial nerves. It can cause slurred speech, difficulty swallowing, weak coughing, poor handling of saliva, and sometimes breathing problems.

\medterm{Bulbar Paralysis} Weakness or paralysis of muscles used for speaking, swallowing, chewing, coughing, and moving the tongue or palate. It may result from brainstem, motor-neuron, or neuromuscular disease and can threaten breathing.

\medterm{Bulbocavernosus} A paired muscle in the pelvic floor that helps empty the urethra and contributes to erection and ejaculation. It surrounds structures at the base of the penis or the vaginal opening.

\medterm{Bulbospongiosus} A paired pelvic-floor muscle that helps empty the urethra and assists erection and the release of genital secretions during sexual activity.

\medterm{Bulbourethral Gland} One of two small glands below the prostate that release a slippery, alkaline fluid into the urethra during sexual arousal. The fluid helps lubricate the passage and reduce acidity.

\medterm{Bulbourethral Glands} Two small male reproductive glands, also called Cowper glands, that release mucus-like fluid into the urethra during sexual arousal. The fluid lubricates the urethra and may reduce acidity.

\medterm{Bulging} Projecting or swelling outward beyond the usual shape. A bulge may involve a disc, membrane, organ, skin area, or body wall, and its significance depends on the structure and cause.

\medterm{Bulging Eyes} Alternate name; see \textit{Exophthalmos}.
\medterm{Bulging Fontanelle} An outward protrusion or tense swelling of an infant's soft spot (fontanelle) when the baby is calm and upright, indicating elevated intracranial pressure from conditions such as meningitis, encephalitis, hydrocephalus, or intracranial hemorrhage.

\synonyms
Tense Fontanelle; Fontanelle Protrusion; Raised Intracranial Pressure in Infants
\medterm{Bulimia} Alternate name; see \textit{Bulimia Nervosa}.

\medterm{Bulimia Nervosa} An eating disorder involving recurrent binge-eating episodes, a sense of loss of control, and recurrent inappropriate compensatory behaviors such as vomiting, laxative misuse, fasting, or excessive exercise. Self-evaluation is unduly influenced by body shape or weight, and the pattern causes clinically significant distress or impairment.

\synonyms
Bulimia; Bulimic Disorder

\medterm{Bulla} A large fluid-filled blister on the skin or mucous membrane, or a large air-filled space in an organ such as the lung. The cause and treatment depend on its location.

\medterm{Bullae} Large fluid-filled blisters on or beneath the skin or a mucous membrane. They can result from friction, burns, infection, allergy, medicines, or autoimmune disease and may need care if widespread or infected.

\medterm{Bullectomy} An operation to remove one or more large air-filled spaces in damaged lung tissue. It may be considered when a giant bulla compresses healthier lung, causes severe breathlessness, or contributes to repeated pneumothorax.

\synonyms
Excision of Pulmonary Bullae; Resection of Giant Bulla

\medterm{Bullet Track} The path a firearm projectile takes through the body, from the entry wound through injured tissues to an exit wound or a retained bullet. The path may damage skin, organs, blood vessels, nerves, or bone.

\medterm{Bullet Wounds} Penetrating or perforating injuries caused by firearm projectiles. They may damage skin, organs, blood vessels, nerves, and bones and require urgent assessment even when the external wound looks small.

\medterm{Bullous} Describes a condition involving large, raised, fluid-filled blisters. Bullous disease can result from infection, burns, medicines, allergy, or an autoimmune reaction.

\medterm{Bullous Impetigo} A contagious superficial bacterial skin infection, usually caused by Staphylococcus, that produces fragile fluid-filled blisters and yellow or varnish-like crusts. It is especially common in infants and young children.

\medterm{Bullous Pemphigoid} An autoimmune blistering disease, usually in older adults, causing intensely itchy skin and firm blisters on normal or reddened skin. It results from antibodies attacking proteins that hold the skin layers together.

\synonyms
BP; Paraneoplastic Pemphigoid

\medterm{Bully} A person who repeatedly harms, threatens, or intimidates someone with less power. Bullying can cause injury, anxiety, depression, sleep problems, school avoidance, and other health effects.

\medterm{Bullying} Repeated physical, verbal, social, or online aggression involving a power imbalance. It can harm mental and physical health; safety planning, support, and intervention by responsible adults are important.

\medterm{Bumetanide} A diuretic that makes the kidneys remove extra salt and water in urine. It is used for swelling related to heart, liver, or kidney disease; excessive fluid loss can cause dehydration and abnormal blood salts.

\medterm{Bumps} Raised or palpable areas on or under the skin. Causes range from harmless cysts, blocked follicles, and insect bites to injury, allergy, infection, or tumours; a growing, painful, hard, or persistent bump requires assessment.

\medterm{BUN} Alternate name; see \textit{Blood Urea Nitrogen Test}.

\medterm{Bundle Branch Block} Delay or interruption of electrical signals through one of the heart's lower-chamber conduction branches. It may cause no symptoms or reflect heart disease and is diagnosed on an ECG.

\medterm{Bundle Of His} A short pathway of specialised heart tissue that carries electrical signals from the atria to the ventricles. It divides into right and left branches so the lower chambers contract in a coordinated way.

\medterm{Bundle-Branch Block} Delayed or interrupted electrical signalling through one of the heart's lower-chamber conduction branches. It can change the ECG and heartbeat pattern and may be harmless or related to underlying heart disease.

\medterm{Bungarotoxins} Poisons in the venom of krait snakes that disrupt nerve signals to muscles. A bite can cause drooping eyelids, weakness, progressive paralysis, and breathing failure and requires emergency treatment.

\medterm{Bunion} A change in the joint at the base of the big toe that pushes the toe toward the others and creates a bump on the inner foot. It may cause pain, redness, swelling, stiffness, and difficulty wearing shoes.

\medterm{Bunion Removal} An operation to correct a painful bunion by realigning bone and, when needed, balancing nearby soft tissues. It is considered when symptoms persist despite suitable shoes and other nonsurgical care.

\medterm{Bunionectomy} Surgery to correct a painful bunion by cutting or shifting bone and adjusting surrounding tissues. The operation is chosen according to the deformity, symptoms, activity, other health conditions, and response to nonsurgical care.

\medterm{Bunionette} A bony prominence at the outer side of the joint below the little toe. It may become painful, red, or swollen from pressure, tight shoes, or altered foot mechanics.

\synonyms
Tailor's Bunion; Bunion of the Fifth Toe

\medterm{Bunions} Deformities at the base of the big toe that push the toe toward the others and create a bump on the inner foot. They may cause pain, swelling, stiffness, and difficulty finding comfortable shoes.

\medterm{Bunioplasty} Surgery to correct a bunion by realigning the toe joint and, when needed, balancing nearby soft tissues. The operation is selected according to the deformity and symptoms.

\medterm{Bunyaviridae} An older virus-family grouping that included several viruses spread by insects, ticks, or rodents. Modern classification has changed; individual viruses differ in the illnesses they cause.

\medterm{Bunyaviruses} Alternate name; see \textit{Viral Hemorrhagic Fevers}.

\medterm{Buphthalmos} Enlargement of an infant's or young child's eye caused by high pressure inside the eye while its outer coat is still flexible. It is usually associated with congenital or infantile glaucoma and needs urgent eye care.

\medterm{Bupivacaine} A long-acting local anaesthetic that temporarily blocks nerve signals to prevent pain during local, nerve-block, or epidural procedures. Too much in the bloodstream can cause seizures, dangerous heart effects, or cardiac arrest.

\medterm{Bupivacaine Hydrochloride} The hydrochloride form of bupivacaine, a long-acting local anaesthetic used to block pain signals. Accidental injection into a blood vessel can cause severe neurologic or heart toxicity.

\medterm{Buprenorphine} A medicine that partly activates opioid receptors and is used to treat opioid-use disorder and some pain. It can cause dependence and overdose, especially with other sedatives, although respiratory suppression is limited compared with full opioid agonists.

\medterm{Bupropion} A medicine used for depression and, in some settings, to help stop smoking. It changes norepinephrine and dopamine signalling and may cause insomnia, dry mouth, or anxiety; it increases seizure risk and is unsuitable for some seizure or eating disorders.

\medterm{Burkholderia} A group of environmental bacteria found in soil and water. Some species cause melioidosis or glanders, while others cause difficult-to-treat lung, bloodstream, or wound infections, especially in people with chronic lung disease or weakened immunity.

\medterm{Burkholderia Cenocepacia} A bacterium in the Burkholderia cepacia complex that can cause severe, persistent lung infection, especially in people with cystic fibrosis or other chronic lung disease.

\medterm{Burkholderia Cepacia} An environmental bacterium in the Burkholderia cepacia complex. It can cause opportunistic lung, bloodstream, wound, or device-related infection, particularly in people with cystic fibrosis or weakened immunity.

\medterm{Burkholderia Gladioli} An environmental bacterium that can rarely cause lung, bloodstream, or wound infection, especially in people with cystic fibrosis or weakened immunity.

\medterm{Burkholderia Mallei} The bacterium that causes glanders, an infection mainly affecting horses and related animals. People can become infected through contact with infected animals or tissues.

\medterm{Burkholderia Pseudomallei} The bacterium that causes melioidosis. Infection follows contact with contaminated soil or water and may cause pneumonia, abscesses, bloodstream infection, or long-term illness.

\medterm{Burkholderiaceae} A family of mainly environmental bacteria that includes Burkholderia. Some members cause serious infections in people or animals, while others are harmless or used in environmental research.

\medterm{Burkitt Lymphoma} A very fast-growing cancer of B lymphocytes, often involving the abdomen, jaw, bone marrow, or nervous system. It requires prompt intensive treatment; the pattern differs between endemic, sporadic, and immunodeficiency-associated forms.

\synonyms
Burkitt's Tumor

\medterm{Burkitt's Lymphoma} A very fast-growing cancer of B lymphocytes. It may affect the abdomen, jaw, bone marrow, or brain and spinal cord and needs prompt intensive treatment with careful supportive care.

\medterm{Burkitt's Tumor} Alternate name; see \textit{Burkitt Lymphoma}.

\medterm{Burn} An injury to the skin or deeper tissue caused by heat, chemicals, electricity, radiation, friction, or extreme cold. Severity depends on depth, size, location, and associated injury; extensive or deep burns need urgent medical care.

\medterm{Burn Injury} Alternate name; see \textit{Burns}.

\medterm{Burn Surgery} Surgical care for serious burns, including assessing depth and area, removing dead tissue, covering wounds, performing skin grafts, and treating scars or movement-limiting contractures.

\medterm{Burn, Oral} Thermal, chemical, or electrical injury to the oral mucosal lining, palate, tongue, or lips, resulting in acute pain, erythema, epithelial desquamation, ulceration, and potential upper airway compromise.

\synonyms
Oral Burn; Oral Cavity Burn; Mouth Burn; Thermal Palatal Burn

\medterm{Burning Feet} Burning, tingling, or painful sensations in the feet, often caused by nerve damage from diabetes, nutritional deficiency, alcohol, medicines, or nerve compression. New, progressive, or weakness-associated symptoms need assessment.

\medterm{Burning Mouth Syndrome} A chronic oral pain disorder causing daily burning, scalding, or tingling discomfort in the mouth, often involving the tongue, lips, or palate, usually with a normal-appearing mucosa. Altered taste or dry mouth may occur. It may be primary, with no identifiable cause, or secondary to conditions or exposures such as nutritional deficiency, oral infection, medication effects, allergy, diabetes, or hormonal change; evaluation excludes treatable causes.

\synonyms
Scalded Mouth Syndrome
Stomatodynia

\medterm{Burning Thigh Pain} Alternate name; see \textit{Meralgia Paresthetica}.

\medterm{Burnout} A work-related state of emotional exhaustion, mental distance or cynicism, and reduced effectiveness caused by chronic workplace stress. It is not the same as depression, although both may occur together.

\medterm{Burns} Tissue injuries caused by heat, chemicals, electricity, radiation, friction, or extreme cold. Seriousness depends on depth, area, and location; burns affecting the face, airway, hands, genitals, major joints, or large areas need prompt care.

\synonyms
Thermal Injury; Burn Trauma

\medterm{Burns, First Aid} Alternate name; see \textit{Burns}.

\medterm{Burnt-Out Joint} An older description of a joint in which active inflammation has subsided but permanent damage, deformity, pain, or restricted movement remains.

\medterm{Burping} Release of swallowed air or gas from the stomach through the mouth. Occasional burping is normal; frequent, painful, or newly worsening burping may result from swallowed air, reflux, food intolerance, or another digestive disorder.

\medterm{Burr Cells} Erythrocytes that exhibit numerous (usually 10 to 30) short, blunt, evenly spaced, and symmetrical spicules across their entire surface membrane. They are typically observed in end-stage renal disease and uremia, severe hypomagnesemia, and pyruvate kinase deficiency.

\synonyms
Echinocytes; Crenated Red Cells

\medterm{Burr Hole Evacuation} An emergency or urgent neurosurgical procedure in which one or two small holes (approximately 10 to 14 mm in diameter) are drilled through the skull bone to drain an abnormal fluid or blood collection compressing the brain. Most frequently performed to treat chronic subdural hematomas (which commonly occur in older adults after minor, often forgotten head injuries), a neurosurgeon incises the scalp and dura mater, washes out the liquefied old blood and clots with warm sterile saline, and places a temporary flexible drain to prevent fluid reaccumulation. By rapidly decompressing the brain, burr hole drainage reverses life-threatening brain herniation, confusion, focal weakness, and lethargy with significantly less surgical trauma, blood loss, and recovery time than a formal craniotomy.

\synonyms
Burr Hole Drainage; Trephination; Chronic Subdural Hematoma Drainage; Burr Hole Craniostomy

\medterm{Burroughs} An inconsistent surname or eponym in medical sources rather than a precise diagnosis; the associated condition or person must be identified from context.

\medterm{Bursa} A small fluid-filled sac that reduces friction between moving tissues, such as a tendon and bone or skin and bone, especially near a joint. Inflammation of a bursa is called bursitis.

\synonyms
Bursae
Synovial Bursae

\medterm{Bursa Inflammation} Alternate name; see \textit{Bursitis}.

\medterm{Bursectomy} Surgical removal of a bursa, usually considered for persistent painful bursitis, recurrent infection, or a structural problem that has not improved with appropriate nonsurgical treatment.

\medterm{Bursitis} Inflammation of a bursa, the small fluid-filled sac that cushions movement between tendons, muscles, bones, and skin. It causes localized pain, tenderness, swelling, or reduced movement and may result from pressure, injury, infection, or arthritis.

\medterm{Bursitis Of The Heel} Inflammation of a bursa near the heel, causing localized tenderness, swelling, and pain when walking or when pressure is applied. It may follow repeated friction, injury, inflammatory arthritis, or infection.

\medterm{Bursitis of the Knee} Alternate name; see \textit{Knee Bursitis}.

\medterm{Bursitis, Calcific} Alternate name; see \textit{Calcific Bursitis}.

\medterm{Bursitis, Shoulder} Alternate name; see \textit{Shoulder Bursitis}.

\medterm{Bursopathy} Any disorder of a bursa, including inflammation, infection, bleeding, thickening, or degeneration. Affected bursae may cause localized pain, swelling, tenderness, and restricted movement around a joint.

\medterm{Buruli Ulcer} A skin infection caused by Mycobacterium ulcerans, usually affecting the arms or legs. It often begins as a painless lump or swelling and progresses to tissue destruction and a large ulcer; untreated disease can damage deeper tissue and cause permanent disability.

\medterm{Burundanga} A colloquial term for scopolamine (hyoscine hydrobromide), a tropane alkaloid and muscarinic receptor antagonist that induces central anticholinergic effects, anterograde amnesia, sedation, and severe toxicity.

\synonyms
Scopolamine; Hyoscine; Hyoscine Hydrobromide

\medterm{Buserelin} A synthetic gonadotropin-releasing hormone analogue. Continuous treatment first briefly increases and then suppresses pituitary reproductive hormones, lowering ovarian or testicular hormone production in selected fertility and hormone-sensitive conditions.

\medterm{Buspirone} A nonbenzodiazepine medicine used mainly for generalized anxiety disorder. It acts on serotonin signaling and usually takes days to weeks to help; it does not provide immediate relief from an acute panic attack.

\medterm{Busulfan} An alkylating chemotherapy medicine used mainly to prepare the body for blood-forming stem-cell transplantation and sometimes to treat chronic myeloid leukemia. It can suppress bone marrow and may cause seizures, infertility, lung injury, or liver damage.

\medterm{Butanol} A flammable industrial alcohol used as a solvent. Swallowing it or substantial exposure to its vapor can irritate the eyes, skin, and airways and may cause dizziness, vomiting, drowsiness, or other nervous-system effects.

\medterm{Butazolidin Overdose} Toxic effects from excessive phenylbutazone, an older nonsteroidal anti-inflammatory drug. Severe exposure may cause gastrointestinal bleeding, kidney injury, bone-marrow suppression, metabolic disturbances, or seizures.

\medterm{Butocarboxim} A carbamate insecticide that inhibits acetylcholinesterase. Significant exposure can cause cholinergic poisoning, with sweating, salivation, vomiting, diarrhea, wheezing, muscle weakness, seizures, or respiratory failure.

\medterm{Butter, Margarine, And Cooking Oils} Butter, margarine, and cooking oils are dietary fat sources. Their nutritional effects depend on portion size, fatty-acid composition, and the overall eating pattern.

\medterm{Buttiauxella} A genus of environmental gram-negative bacteria in the Enterobacterales; human infection is uncommon and usually opportunistic.

\medterm{Buttiauxella Agrestis} A species of environmental Gram-negative bacteria found in water, soil, and other natural settings. It is rarely reported from human clinical samples; when detected, the specimen source and the patient’s symptoms determine whether it represents infection or harmless contamination.

\medterm{Buttiauxella Ferragutiae} A species of environmental Gram-negative bacteria. It is uncommon in human disease, and a laboratory finding requires interpretation with the specimen source, repeat cultures, symptoms, and other evidence before deciding whether treatment is needed.

\medterm{Buttiauxella Gaviniae} A species of Buttiauxella bacteria found mainly in environmental samples. Human infection is uncommon; identification from a clinical specimen should be reviewed with the microbiology laboratory and interpreted according to the patient’s condition.

\medterm{Buttiauxella Noackiae} A species of environmental Gram-negative bacteria that is rarely associated with human infection. Its significance depends on where it was found, whether it was recovered repeatedly, and whether the patient has compatible symptoms or signs of infection.

\medterm{Buttock} The buttock is the rounded posterior region overlying the gluteal muscles, used for movement, sitting, examination, and selected injections.

\medterm{Buttocks} The paired fleshy regions overlying the gluteal muscles at the back of the pelvis, formed mainly by muscle and subcutaneous fat.

\medterm{Button Batteries} Button batteries are small batteries that can cause severe tissue injury if swallowed or inserted into the nose or ear. Suspected ingestion requires emergency evaluation.

\medterm{Button Magnets In The Nasal Cavity} Small magnets placed in or inhaled into the nose can attract across the nasal partition and cut off its blood supply. They may cause tissue death or a hole and require urgent medical removal.

\medterm{Buttonhole Deformity} Buttonhole deformity is a finger deformity with flexion of the proximal interphalangeal joint and hyperextension of the distal interphalangeal joint, often due to central slip injury or rheumatoid disease.

\medterm{Buttressing} Buttressing is reinforcement or thickening that supports a weakened structure; in bone it may describe cortical stress adaptation or surgical structural support.

\medterm{Butylamine} Butylamine is an irritating industrial amine whose vapor or liquid can injure the eyes, skin, and respiratory tract after exposure.

\medterm{Butyric Acid} Butyric acid is a short-chain fatty acid produced by gut bacteria during fiber fermentation and used by colon cells as an energy source; concentrated exposure irritates tissue.

\medterm{Bx} A common medical abbreviation for biopsy, the removal of cells or tissue for microscopic or laboratory examination.

\medterm{By-Pass} A route made around a blocked or narrowed passage. In medicine, a bypass may redirect blood, food, bile, urine, or another fluid through a natural or surgically created pathway.

\medterm{Bypass} Creation of an alternative route around a blocked or narrowed passage. A bypass can restore blood flow or redirect body contents, using a natural pathway, tube, vein, or artificial graft depending on the organ involved.

\medterm{Bypass Graft} A vein or artificial tube surgically connected to route blood around a blocked or weakened artery. Grafts are used in coronary, leg, and aortic operations; blood flow, infection, clotting, and long-term narrowing affect their durability.

\medterm{Bypass Surgery} An operation that creates a new route around a blocked or diseased passage. Examples include coronary and leg-artery bypasses and gastric bypass; possible complications include bleeding, infection, clotting, and narrowing of the connection.

\medterm{Byssinosis} An occupational lung disease caused by repeated inhalation of dust from cotton, flax, hemp, or other textile fibers. It can cause chest tightness, cough, and breathing difficulty that worsen during or after work exposure.
